% ============================================================
% RAPPORT DE PROJET DE FIN D'ÉTUDES
% Plateforme de Gouvernance et Monitoring des Modèles de Langage (LLM)
% Institut Supérieur de Gestion de Tunis — 2024/2025
% ============================================================

\documentclass[12pt, a4paper, oneside]{report}

% ── Encodage et langue ──────────────────────────────────────
\usepackage[utf8]{inputenc}
\usepackage[T1]{fontenc}
\usepackage[french]{babel}
\usepackage{lmodern}

% ── Géométrie de page ───────────────────────────────────────
\usepackage[
  top=2.5cm,
  bottom=2.5cm,
  left=3cm,
  right=2.5cm,
  headheight=16pt
]{geometry}

% ── En-têtes et pieds de page ───────────────────────────────
\usepackage{fancyhdr}
\pagestyle{fancy}
\fancyhf{}
\fancyhead[L]{\leftmark}
\fancyhead[R]{\thepage}
\renewcommand{\headrulewidth}{0.4pt}
\fancypagestyle{plain}{
  \fancyhf{}
  \fancyfoot[C]{\thepage}
  \renewcommand{\headrulewidth}{0pt}
}

% ── Packages graphiques ─────────────────────────────────────
\usepackage{graphicx}
\usepackage{float}
\usepackage{subcaption}
\graphicspath{{figures/}}

% ── Tableaux ────────────────────────────────────────────────
\usepackage{longtable}
\usepackage{booktabs}
\usepackage{array}
\usepackage{multirow}
\usepackage{tabularx}
\usepackage{colortbl}

% ── Couleurs ────────────────────────────────────────────────
\usepackage[dvipsnames,table]{xcolor}

% ── Code source ─────────────────────────────────────────────
\usepackage{listings}
\lstset{
  basicstyle=\ttfamily\small,
  keywordstyle=\color{MidnightBlue}\bfseries,
  stringstyle=\color{ForestGreen},
  commentstyle=\color{gray}\itshape,
  numberstyle=\tiny\color{gray},
  numbers=left,
  numbersep=8pt,
  frame=single,
  framerule=0.5pt,
  rulecolor=\color{gray!40},
  backgroundcolor=\color{gray!5},
  breaklines=true,
  breakatwhitespace=false,
  showstringspaces=false,
  tabsize=4,
  captionpos=b,
  xleftmargin=15pt,
  framexleftmargin=15pt,
  aboveskip=10pt,
  belowskip=10pt,
  literate={é}{{\'e}}1 {è}{{\`e}}1 {ê}{{\^e}}1 {ë}{{\"e}}1
           {à}{{\`a}}1 {â}{{\^a}}1 {ù}{{\`u}}1 {û}{{\^u}}1
           {ô}{{\^o}}1 {ç}{{\c{c}}}1 {î}{{\^i}}1 {ï}{{\"i}}1
}

% Définition des langages
\lstdefinelanguage{Python3}{
  language=Python,
  morekeywords={async, await, f, True, False, None, self, cls, Mapped, mapped_column},
}

\lstdefinelanguage{TypeScript}{
  language=Java,
  morekeywords={interface, type, const, let, async, await, export, import, from, default, string, number, boolean, void, null, undefined, any},
}

% ── Mathématiques ───────────────────────────────────────────
\usepackage{amsmath,amssymb,amsfonts}

% ── Divers ──────────────────────────────────────────────────
\usepackage{enumitem}
\usepackage{pifont}
\usepackage{caption}
\usepackage{setspace}
\usepackage{parskip}
\usepackage{titlesec}
\usepackage{appendix}
\usepackage{url}

% ── Liens hypertexte (DOIT être chargé AVANT tocloft) ──────
\usepackage[
  colorlinks=true,
  linkcolor=MidnightBlue,
  citecolor=ForestGreen,
  urlcolor=RoyalBlue,
  bookmarks=true,
  bookmarksnumbered=true,
  pdfstartview=FitH
]{hyperref}

% ── Table des matières personnalisée (APRÈS hyperref) ───────
\usepackage{tocloft}

% ── Nomenclature / Acronymes ───────────────────────────────
\usepackage[intoc, french]{nomencl}
\makenomenclature

% ── Personnalisation des titres de chapitres ────────────────
\titleformat{\chapter}[display]
  {\normalfont\huge\bfseries\color{MidnightBlue}}
  {\chaptertitlename\ \thechapter}{20pt}{\Huge}
\titlespacing*{\chapter}{0pt}{-20pt}{30pt}

% ── Interligne 1.5 ─────────────────────────────────────────
\onehalfspacing

% ── Profondeur de la table des matières ────────────────────
\setcounter{tocdepth}{3}
\setcounter{secnumdepth}{3}


% ============================================================
% DÉBUT DU DOCUMENT
% ============================================================
\begin{document}

% ── Pages préliminaires (numérotation romaine) ─────────────
\pagenumbering{roman}

% Page de garde
% ============================================================
% PAGE DE GARDE
% ============================================================

\begin{titlepage}
\begin{center}

% ── Logo et en-tête institutionnel ──────────────────────────
\vspace*{-1cm}

% Espace réservé pour le logo de l'Université de Tunis / ISGT
% \includegraphics[width=4cm]{logo_universite_tunis.png}
% \hfill
% \includegraphics[width=4cm]{logo_isgt.png}

\textsc{\Large Université de Tunis}\\[0.3cm]
\textsc{\large Institut Supérieur de Gestion de Tunis}\\[0.2cm]
\textsc{Département des Technologies de l'Information et de la Communication}\\[1.5cm]

% ── Ligne de séparation ────────────────────────────────────
\rule{\textwidth}{1.5pt}\\[0.4cm]

% ── Type de document ───────────────────────────────────────
{\Large \textbf{Rapport de Projet de Fin d'Études}}\\[0.3cm]
{\large Présenté en vue de l'obtention du}\\[0.2cm]
{\large \textbf{Diplôme National de Licence en Informatique de Gestion}}\\[0.5cm]

\rule{\textwidth}{1.5pt}\\[1.2cm]

% ── Titre du projet ────────────────────────────────────────
{\LARGE \textbf{Plateforme de Gouvernance et Monitoring}}\\[0.3cm]
{\LARGE \textbf{des Modèles de Langage (LLM)}}\\[0.4cm]
{\large \textit{LLM Dashboard — AI Governance Platform}}\\[1.5cm]

% ── Informations ───────────────────────────────────────────
\begin{minipage}[t]{0.45\textwidth}
\begin{flushleft}
\textit{Élaboré par :}\\[0.3cm]
\textbf{Houcine LAGHMANI}\\[1cm]

\textit{Entreprise d'accueil :}\\[0.3cm]
\textbf{TIM Tunisie}\\
\url{https://tim-tunisie.com/}
\end{flushleft}
\end{minipage}
\hfill
\begin{minipage}[t]{0.45\textwidth}
\begin{flushright}
\textit{Encadrant académique :}\\[0.3cm]
\textbf{Mme Salsabil HADJI}\\[1cm]

\textit{Encadrant professionnel :}\\[0.3cm]
\textbf{TIM Tunisie}\\
\end{flushright}
\end{minipage}

\vfill

% ── Espace réservé pour les logos ──────────────────────────
% \includegraphics[width=3cm]{logo_tim_tunisie.png}

% ── Année académique ───────────────────────────────────────
\rule{\textwidth}{0.5pt}\\[0.3cm]
{\Large \textbf{Année Universitaire : 2025 -- 2026}}

\end{center}
\end{titlepage}

% Page blanche après la page de garde
\newpage\thispagestyle{empty}\mbox{}\newpage


% Dédicaces
% ============================================================
% DÉDICACES
% ============================================================

\newpage
\thispagestyle{empty}

\vspace*{3cm}

\begin{flushright}
\Large\textit{Dédicaces}
\end{flushright}

\vspace{2cm}

\begin{quote}
\large
\textit{Je dédie ce travail à\ldots}

\vspace{1cm}

\textit{À mes chers parents,}\\
\textit{pour leur soutien inconditionnel, leurs sacrifices}\\
\textit{et leur amour tout au long de mon parcours académique.}

\vspace{0.8cm}

\textit{À ma famille,}\\
\textit{pour leur encouragement permanent et leur confiance en moi.}

\vspace{0.8cm}

\textit{À mes amis et collègues,}\\
\textit{pour les moments partagés et l'entraide précieuse.}

\vspace{0.8cm}

\textit{À tous ceux qui ont contribué, de près ou de loin,}\\
\textit{à la réalisation de ce projet.}

\vspace{1.5cm}

\begin{flushright}
\textbf{Houcine LAGHMANI}
\end{flushright}
\end{quote}

\newpage


% Remerciements
% ============================================================
% REMERCIEMENTS
% ============================================================

\newpage
\thispagestyle{empty}

\begin{center}
{\Large\textbf{Remerciements}}
\end{center}

\vspace{1.5cm}

Au terme de ce projet de fin d'études, je tiens à exprimer ma profonde gratitude envers toutes les personnes qui ont contribué à la réalisation de ce travail.

\vspace{0.5cm}

Je remercie tout d'abord \textbf{\underline{\hspace{4cm}}}, mon encadrant académique à l'Institut Supérieur de Gestion de Tunis, pour ses conseils avisés, son suivi rigoureux et sa disponibilité tout au long de ce projet.

\vspace{0.5cm}

Je tiens également à remercier \textbf{\underline{\hspace{4cm}}}, mon encadrant professionnel au sein de \textbf{TIM Tunisie}, pour la confiance qu'il m'a accordée, son accompagnement technique et ses orientations précieuses dans la compréhension des besoins métier liés aux systèmes \textbf{QALITAS} et \textbf{GMAO PRO}.

\vspace{0.5cm}

Mes remerciements s'adressent aussi à l'ensemble de l'équipe de \textbf{TIM Tunisie} pour l'accueil chaleureux, l'environnement de travail stimulant et les échanges enrichissants qui ont facilité la réalisation de cette plateforme.

\vspace{0.5cm}

Je remercie sincèrement les membres du jury pour l'honneur qu'ils me font en acceptant d'évaluer ce travail. Leurs remarques et suggestions seront d'un apport considérable.

\vspace{0.5cm}

Enfin, je ne saurais oublier mes parents, ma famille et mes amis pour leur soutien moral indéfectible et leurs encouragements permanents.

\vspace{1cm}

\begin{flushright}
\textit{Houcine LAGHMANI}\\
\textit{Tunis, Juin 2026}
\end{flushright}

\newpage


% Résumé / Abstract
% ============================================================
% RÉSUMÉ
% ============================================================

\chapter*{Résumé}
\addcontentsline{toc}{chapter}{Résumé}

\vspace{0.5cm}

Le présent projet de fin d'études, réalisé au sein de l'entreprise \textbf{TIM Tunisie}, porte sur la conception et le développement d'une \textbf{plateforme SaaS de gouvernance et de monitoring des modèles de langage de grande taille (LLM)}. Face à l'intégration croissante de l'intelligence artificielle dans les produits industriels de l'entreprise --- notamment les systèmes QALITAS et GMAO PRO --- le besoin d'un outil centralisé de supervision, de maîtrise des coûts et de traçabilité s'est imposé comme une nécessité stratégique.

La solution développée, baptisée \textbf{LLM Dashboard}, se positionne comme un proxy intelligent entre les applications clientes et les fournisseurs de LLM (OpenAI, Anthropic, Google, Groq, Cohere). Elle intègre \textbf{onze modules fonctionnels} couvrant l'ensemble du cycle de vie des interactions LLM~: authentification et contrôle d'accès (RBAC), Gateway proxy avec gouvernance temps réel, détection d'anomalies par apprentissage automatique (STL, Isolation Forest, CUSUM), explication automatique des anomalies par IA, prévision budgétaire à trois scénarios, optimisation des coûts, analytics et agrégation, facturation contractuelle, traçage des sessions, gestion des prompts et tableau de bord exécutif interactif.

L'architecture technique repose sur une stack moderne~: FastAPI (Python~3.13), Next.js~16 (React~19), PostgreSQL avec TimescaleDB et pgvector, Redis~7, Celery~5 et Docker. Le développement a été conduit selon la méthodologie Scrum, structuré en trois sprints, et validé par 180~tests unitaires et d'intégration.

\vspace{0.5cm}

\textbf{Mots clés~:} LLM, Gouvernance, Monitoring, Multi-tenant, SaaS, Détection d'anomalies, Apprentissage automatique, FastAPI, TimescaleDB, Scrum.

\newpage

% ============================================================
% ABSTRACT
% ============================================================

\chapter*{Abstract}
\addcontentsline{toc}{chapter}{Abstract}

\vspace{0.5cm}

This graduation project, carried out within \textbf{TIM Tunisie}, focuses on the design and development of a \textbf{multi-tenant SaaS platform for Large Language Model (LLM) governance and monitoring}. As the company progressively integrates AI capabilities into its industrial products --- QALITAS (quality management) and GMAO PRO (maintenance management) --- the need for a centralized tool to supervise, control costs, and ensure traceability of LLM usage has become a strategic imperative.

The developed solution, named \textbf{LLM Dashboard}, acts as an intelligent proxy positioned between client applications and LLM providers (OpenAI, Anthropic, Google, Groq, Cohere). It integrates \textbf{eleven functional modules} covering the complete lifecycle of LLM interactions: JWT authentication with 4-level RBAC, intelligent gateway proxy with real-time governance, anomaly detection using three complementary machine learning algorithms (STL decomposition, Isolation Forest, CUSUM), AI-powered anomaly explanation, three-scenario budget forecasting, cost optimization recommendations, daily and monthly analytics aggregation, contractual billing with automatic invoice generation, session tracing with waterfall visualization, prompt template management, and an interactive executive dashboard with real-time WebSocket updates.

The technical architecture is built on a modern stack: FastAPI (Python~3.13), Next.js~16 (React~19), PostgreSQL with TimescaleDB and pgvector extensions, Redis~7, Celery~5, and Docker containerization. Development followed the Scrum methodology across three sprints and was validated by 180~unit and integration tests.

\vspace{0.5cm}

\textbf{Keywords:} LLM, Governance, Monitoring, Multi-tenant, SaaS, Anomaly Detection, Machine Learning, FastAPI, TimescaleDB, Scrum.

\newpage


% Tables des matières, figures, tableaux
\tableofcontents
\listoffigures
\listoftables

% Glossaire
% ============================================================
% GLOSSAIRE ET ACRONYMES
% ============================================================

\chapter*{Glossaire}
\addcontentsline{toc}{chapter}{Glossaire}
\markboth{Glossaire}{Glossaire}

\vspace{0.3cm}

\begin{longtable}{p{3.5cm} p{11cm}}

\textbf{API} & \textit{Application Programming Interface} --- Interface de programmation applicative permettant la communication entre systèmes logiciels. \\[6pt]

\textbf{ASGI} & \textit{Asynchronous Server Gateway Interface} --- Standard d'interface pour les serveurs web Python asynchrones. \\[6pt]

\textbf{BPE} & \textit{Byte Pair Encoding} --- Algorithme de tokenisation sous-mot utilisé par les LLM pour décomposer le texte en unités élémentaires. \\[6pt]

\textbf{CI/CD} & \textit{Continuous Integration / Continuous Deployment} --- Pratique d'intégration et de déploiement continus automatisés. \\[6pt]

\textbf{CORS} & \textit{Cross-Origin Resource Sharing} --- Mécanisme HTTP permettant à un serveur d'autoriser les requêtes provenant d'origines différentes. \\[6pt]

\textbf{CRUD} & \textit{Create, Read, Update, Delete} --- Les quatre opérations fondamentales de gestion des données persistantes. \\[6pt]

\textbf{CUSUM} & \textit{Cumulative Sum Control Chart} --- Algorithme de contrôle statistique de processus pour la détection de changements de moyenne. \\[6pt]

\textbf{DTO} & \textit{Data Transfer Object} --- Objet de transfert de données utilisé pour structurer les échanges entre les couches applicatives. \\[6pt]

\textbf{GHCR} & \textit{GitHub Container Registry} --- Registre de conteneurs Docker intégré à GitHub pour le stockage et la distribution d'images. \\[6pt]

\textbf{GMAO} & \textit{Gestion de la Maintenance Assistée par Ordinateur} --- Système logiciel de planification et de suivi de la maintenance industrielle. \\[6pt]

\textbf{HSTS} & \textit{HTTP Strict Transport Security} --- En-tête de sécurité HTTP imposant l'utilisation exclusive de HTTPS. \\[6pt]

\textbf{JWT} & \textit{JSON Web Token} --- Standard ouvert (RFC~7519) pour la transmission sécurisée d'informations entre parties sous forme de jeton signé. \\[6pt]

\textbf{KPI} & \textit{Key Performance Indicator} --- Indicateur clé de performance permettant de mesurer l'atteinte des objectifs. \\[6pt]

\textbf{LLM} & \textit{Large Language Model} --- Modèle de langage de grande taille, réseau de neurones profond pré-entraîné sur des corpus textuels massifs. \\[6pt]

\textbf{LLMOps} & \textit{Large Language Model Operations} --- Ensemble de pratiques pour l'exploitation des LLM en production (monitoring, gouvernance, sécurité). \\[6pt]

\textbf{MoSCoW} & Méthode de priorisation des besoins~: \textbf{M}ust have, \textbf{S}hould have, \textbf{C}ould have, \textbf{W}on't have. \\[6pt]

\textbf{ORM} & \textit{Object-Relational Mapping} --- Technique de correspondance entre les objets d'un langage de programmation et les tables d'une base de données relationnelle. \\[6pt]

\textbf{PBKDF2} & \textit{Password-Based Key Derivation Function 2} --- Fonction de dérivation de clé utilisée pour le hachage sécurisé des mots de passe. \\[6pt]

\textbf{RBAC} & \textit{Role-Based Access Control} --- Modèle de contrôle d'accès basé sur les rôles attribués aux utilisateurs. \\[6pt]

\textbf{REST} & \textit{Representational State Transfer} --- Style d'architecture logicielle pour la conception d'interfaces de programmation web. \\[6pt]

\textbf{RLHF} & \textit{Reinforcement Learning from Human Feedback} --- Technique d'entraînement des LLM par renforcement à partir de retours humains. \\[6pt]

\textbf{ROI} & \textit{Return on Investment} --- Retour sur investissement, indicateur de rentabilité d'une action ou d'un projet. \\[6pt]

\textbf{RPM} & \textit{Requests Per Minute} --- Nombre de requêtes par minute, métrique de fréquence d'appels API. \\[6pt]

\textbf{SaaS} & \textit{Software as a Service} --- Modèle de distribution logicielle où l'application est hébergée par le fournisseur et accessible via Internet. \\[6pt]

\textbf{SSR / CSR} & \textit{Server-Side Rendering / Client-Side Rendering} --- Stratégies de rendu des pages web, respectivement côté serveur et côté client. \\[6pt]

\textbf{STL} & \textit{Seasonal and Trend decomposition using Loess} --- Technique de décomposition des séries temporelles en composantes de tendance, de saisonnalité et de résidu. \\[6pt]

\textbf{TLS} & \textit{Transport Layer Security} --- Protocole de sécurité assurant le chiffrement des communications réseau. \\[6pt]

\textbf{TTL} & \textit{Time To Live} --- Durée de validité d'un élément dans un cache ou d'un enregistrement DNS. \\[6pt]

\textbf{UML} & \textit{Unified Modeling Language} --- Langage de modélisation graphique utilisé en ingénierie logicielle pour la conception de systèmes. \\[6pt]

\textbf{XSS} & \textit{Cross-Site Scripting} --- Type de vulnérabilité de sécurité web permettant l'injection de scripts malveillants. \\[6pt]

\end{longtable}


% ── Corps du rapport (numérotation arabe) ──────────────────
\newpage
\pagenumbering{arabic}

% Introduction générale
% ============================================================
% INTRODUCTION GÉNÉRALE
% ============================================================

\chapter*{Introduction Générale}
\addcontentsline{toc}{chapter}{Introduction Générale}
\markboth{Introduction Générale}{Introduction Générale}

\vspace{0.5cm}

L'essor fulgurant de l'intelligence artificielle générative constitue sans doute la mutation technologique la plus significative de cette décennie. Au c\oe ur de cette révolution se trouvent les modèles de langage de grande taille --- communément désignés par l'acronyme LLM (\textit{Large Language Models}) --- dont les capacités de compréhension et de génération de texte ont atteint un niveau de sophistication sans précédent. Les travaux fondateurs de Vaswani \textit{et al.}~\cite{vaswani2017} sur l'architecture Transformer, puis les recherches de Brown \textit{et al.}~\cite{brown2020} sur les propriétés émergentes des modèles à grande échelle, ont ouvert la voie à une nouvelle génération d'outils cognitifs : GPT-4 d'OpenAI~\cite{openai2023}, Claude d'Anthropic~\cite{anthropic2024}, Gemini de Google ou encore LLaMA de Meta~\cite{touvron2023}.

L'adoption de ces technologies par le monde industriel est désormais massive. Selon le rapport annuel de McKinsey \& Company publié en mars 2024, plus de 72\,\% des organisations à l'échelle mondiale ont intégré au moins une forme d'intelligence artificielle dans leurs processus opérationnels~\cite{mckinsey2024}. Bommasani \textit{et al.}~\cite{bommasani2022} qualifient les LLM de \og modèles fondamentaux \fg{} (\textit{foundation models}), soulignant leur caractère transversal et leur capacité à transformer simultanément de multiples secteurs d'activité.

\vspace{0.3cm}

Toutefois, cette adoption croissante engendre des problématiques majeures en termes de \textbf{gouvernance}, de \textbf{maîtrise des coûts} et de \textbf{traçabilité}. La consommation de tokens --- unité élémentaire de facturation des services LLM --- représente désormais un poste budgétaire significatif pour les entreprises, souvent opaque et difficilement prévisible. L'absence d'outils de supervision centralisés expose les organisations à des dépassements budgétaires non détectés, à des utilisations non conformes et à une incapacité d'identifier les anomalies de consommation en temps réel. Chen \textit{et al.}~\cite{chen2024llmops} identifient ces défis comme les principaux freins à l'industrialisation des solutions fondées sur les LLM.

\vspace{0.3cm}

C'est dans ce contexte que s'inscrit notre projet de fin d'études, réalisé au sein de \textbf{TIM Tunisie}, entreprise spécialisée dans l'édition de solutions logicielles pour le secteur industriel. TIM Tunisie développe et commercialise deux systèmes d'information de référence~:

\begin{itemize}[noitemsep]
    \item \textbf{QALITAS}~: plateforme intégrée de gestion de la qualité industrielle~;
    \item \textbf{GMAO PRO}~: système de gestion de la maintenance assistée par ordinateur.
\end{itemize}

Ces deux produits intègrent progressivement des fonctionnalités alimentées par les LLM --- analyse automatique de documents qualité, assistance conversationnelle pour les auditeurs, génération de rapports de conformité --- ce qui génère un besoin croissant de gouvernance et de monitoring de ces ressources d'intelligence artificielle.

\vspace{0.3cm}

Notre mission consiste à concevoir et à développer \textbf{LLM Dashboard}, une plateforme SaaS multi-tenant de gouvernance et de monitoring des modèles de langage. Positionnée comme un \textbf{proxy intelligent} entre les applications clientes et les fournisseurs de LLM, cette plateforme assure~:

\begin{itemize}[noitemsep]
    \item un \textbf{suivi en temps réel} de la consommation de tokens et des coûts associés~;
    \item un \textbf{moteur de gouvernance} fondé sur des règles configurables de quotas, de budgets et de contrôle d'accès~;
    \item un système de \textbf{détection d'anomalies} reposant sur trois algorithmes d'apprentissage automatique complémentaires~;
    \item un module d'\textbf{explication automatique} des anomalies par intelligence artificielle~;
    \item un moteur de \textbf{prévision budgétaire} intégrant une analyse de risques à trois scénarios~;
    \item un système de \textbf{facturation contractuelle} avec génération automatisée de factures~;
    \item un \textbf{tableau de bord interactif} enrichi de visualisations en temps réel.
\end{itemize}

\vspace{0.3cm}

Le développement de cette plateforme a été conduit selon la méthodologie agile \textbf{Scrum}~\cite{schwaber2020}, structuré en trois sprints correspondant aux trois releases fonctionnelles du projet. Cette approche itérative a permis de livrer des incréments fonctionnels réguliers et de recueillir un retour continu de la part de l'entreprise d'accueil.

\vspace{0.3cm}

Le présent rapport est organisé comme suit~:

\begin{itemize}[noitemsep]
    \item Le \textbf{premier chapitre} présente l'étude du projet~: l'organisme d'accueil, la problématique identifiée, la solution proposée et la méthodologie de développement adoptée.
    \item Le \textbf{deuxième chapitre} dresse un état de l'art des technologies mobilisées~: les modèles de langage, la gouvernance des LLM, les architectures multi-tenant et la détection d'anomalies par apprentissage automatique.
    \item Le \textbf{troisième chapitre} détaille le Sprint~0, consacré à l'identification des besoins fonctionnels et non fonctionnels, à la modélisation du cas d'utilisation global, à la définition du product backlog et à l'environnement de travail.
    \item Le \textbf{quatrième chapitre} couvre le Sprint~1, dédié à la mise en place de l'infrastructure technique et des modules fondamentaux~: authentification, gateway LLM, gouvernance et tableau de bord.
    \item Le \textbf{cinquième chapitre} présente le Sprint~2, consacré aux composants d'intelligence artificielle~: détection d'anomalies, explication IA, prévision budgétaire et optimisation des coûts.
    \item Le \textbf{sixième chapitre} couvre le Sprint~3, portant sur les fonctionnalités entreprise --- facturation, traçage des sessions, gestion des prompts --- ainsi que sur le déploiement, les tests et la validation.
\end{itemize}

Le rapport se conclut par un bilan global, une discussion des limites identifiées et une présentation des perspectives d'évolution de la plateforme.


% Chapitre 1 — Étude du Projet
% ============================================================
% CHAPITRE 1 — ÉTUDE DU PROJET
% ============================================================

\chapter{Étude du Projet}

\section*{Plan}

\begin{tabular}{r l r}
\textbf{1} & \textbf{Organisme d'accueil} & \textbf{\pageref{sec:organisme}} \\
\textbf{2} & \textbf{Présentation du projet} & \textbf{\pageref{sec:projet_overview}} \\
\textbf{3} & \textbf{Méthodologie de développement} & \textbf{\pageref{sec:methodologie}} \\
\textbf{4} & \textbf{Méthodologies d'ingénierie complémentaires} & \textbf{\pageref{sec:methodologies_complementaires}} \\
\textbf{5} & \textbf{Langage de modélisation~: UML} & \textbf{\pageref{sec:uml}} \\
\end{tabular}

\section*{Introduction}

Ce chapitre introductif pose les fondations de notre projet de fin d'études. Nous débutons par une présentation détaillée de l'organisme d'accueil, \textbf{TIM Tunisie}, en décrivant sa structure organisationnelle et ses produits phares. Par la suite, nous exposons la problématique qui motive ce travail, la solution proposée pour y répondre ainsi que les objectifs visés. Enfin, nous décrivons la méthodologie de développement retenue pour conduire ce projet dans un cadre structuré et itératif.

% ────────────────────────────────────────────────────────────
\section{Organisme d'accueil}
\label{sec:organisme}
% ────────────────────────────────────────────────────────────

\subsection{TIM Tunisie --- Présentation générale}

\textbf{TIM Tunisie} (\textit{Technologies, Informatique et Management}) est une entreprise tunisienne spécialisée dans la conception et le développement de solutions logicielles à destination du secteur industriel. Fondée avec pour vocation d'accompagner la transformation numérique des entreprises industrielles en Tunisie et dans la région MENA, TIM Tunisie a acquis une expertise reconnue dans deux domaines stratégiques~:

\begin{itemize}[noitemsep]
    \item la \textbf{gestion de la qualité} industrielle conformément aux normes internationales~;
    \item la \textbf{gestion de la maintenance} assistée par ordinateur (GMAO).
\end{itemize}

L'entreprise se positionne aujourd'hui comme un acteur clé de l'innovation technologique industrielle, en intégrant progressivement des composantes d'intelligence artificielle --- et notamment des modèles de langage de grande taille --- au sein de ses produits existants.

\begin{figure}[H]
    \centering
    \fbox{\parbox{6cm}{\centering\vspace{1cm}\textit{Logo de TIM Tunisie}\vspace{1cm}}}
    \caption{Logo de TIM Tunisie}
    \label{fig:logo_tim}
\end{figure}

\begin{table}[H]
\centering
\caption{Fiche signalétique de TIM Tunisie}
\label{tab:fiche_tim}
\begin{tabularx}{\textwidth}{|l|X|}
\hline
\textbf{Raison sociale} & TIM Tunisie \\
\hline
\textbf{Secteur d'activité} & Édition de logiciels industriels \\
\hline
\textbf{Produits principaux} & QALITAS, GMAO PRO \\
\hline
\textbf{Domaines} & Qualité industrielle, Maintenance, Intelligence Artificielle \\
\hline
\textbf{Site web} & \url{https://tim-tunisie.com/} \\
\hline
\textbf{Localisation} & Tunis, Tunisie \\
\hline
\end{tabularx}
\end{table}


\subsection{Les produits de TIM Tunisie}

L'activité de TIM Tunisie s'articule autour de deux produits logiciels complémentaires, chacun répondant à des besoins industriels spécifiques et intégrant des fonctionnalités d'intelligence artificielle.

\subsubsection{QALITAS --- Plateforme de gestion de la qualité}

\textbf{QALITAS} est un système d'information intégré dédié à la gestion de la qualité industrielle. Cette plateforme permet aux entreprises de structurer et d'automatiser l'ensemble de leurs processus qualité~:

\begin{itemize}[noitemsep]
    \item gestion des processus de contrôle qualité et d'audit interne~;
    \item suivi des non-conformités et des actions correctives et préventives~;
    \item production de rapports de conformité aux normes internationales (ISO~9001, ISO~14001, ISO~45001)~;
    \item automatisation des workflows de validation et d'approbation documentaire.
\end{itemize}

Avec l'intégration progressive de l'intelligence artificielle, QALITAS exploite désormais des \textbf{modèles de langage (LLM)} pour l'analyse automatique de documents qualité, la génération de rapports de conformité et l'assistance conversationnelle pour les auditeurs internes.

\subsubsection{GMAO PRO --- Gestion de la maintenance assistée par ordinateur}

\textbf{GMAO PRO} constitue la solution de TIM Tunisie pour la gestion de la maintenance industrielle. Ses fonctionnalités principales couvrent~:

\begin{itemize}[noitemsep]
    \item la planification et le suivi des opérations de maintenance préventive et corrective~;
    \item la gestion du parc d'équipements et des pièces de rechange~;
    \item le suivi des indicateurs clés de performance de maintenance (MTBF, MTTR, taux de disponibilité)~;
    \item la génération automatisée de bons de travail et de rapports d'intervention.
\end{itemize}

GMAO PRO intègre également des fonctionnalités d'IA pour la \textbf{maintenance prédictive}, s'appuyant sur les LLM pour analyser les historiques de pannes et formuler des recommandations d'intervention anticipée.


% ────────────────────────────────────────────────────────────
\section{Présentation du projet}
\label{sec:projet_overview}
% ────────────────────────────────────────────────────────────

Ce projet s'inscrit dans le cadre d'un stage de fin d'études en licence en Informatique de Gestion à l'Institut Supérieur de Gestion de Tunis. Le stage a été effectué au sein de TIM Tunisie sous la supervision académique de \textbf{Mme Salsabil Hadji}.

\subsection{Étude de l'existant}

Avant le développement de notre plateforme, l'utilisation des LLM au sein des produits de TIM Tunisie se caractérisait par plusieurs lacunes structurelles~:

\begin{itemize}[noitemsep]
    \item \textbf{Appels directs}~: chaque module applicatif (QALITAS, GMAO PRO) effectuait des appels directs aux API des fournisseurs LLM (OpenAI, Anthropic, Google) sans intermédiaire centralisé~;
    \item \textbf{Absence de proxy}~: aucun point de passage unique ne permettait d'intercepter, de monitorer ou de gouverner les appels~;
    \item \textbf{Clés API partagées}~: les clés d'accès étaient souvent partagées entre plusieurs services, sans isolation par client~;
    \item \textbf{Pas de comptage de tokens}~: la consommation réelle n'était pas tracée de manière systématique~;
    \item \textbf{Facturation manuelle}~: les coûts LLM étaient estimés manuellement à partir des factures des fournisseurs.
\end{itemize}

\subsection{Critique de l'existant}

L'analyse approfondie de la situation existante a révélé des \textbf{lacunes critiques} compromettant la viabilité économique et opérationnelle de l'utilisation des LLM~:

\begin{table}[H]
\centering
\caption{Critique de l'existant --- Problèmes identifiés}
\label{tab:critique_existant}
\begin{tabularx}{\textwidth}{|c|p{4.5cm}|X|}
\hline
\textbf{\#} & \textbf{Problème} & \textbf{Impact} \\
\hline
1 & Pas de visibilité sur la consommation & Impossible de savoir quel agent, modèle ou client consomme le plus \\
\hline
2 & Pas de contrôle des coûts & Risque de dépassements budgétaires non détectés \\
\hline
3 & Pas d'isolation multi-tenant & Un client peut monopoliser les ressources sans limitation \\
\hline
4 & Pas de gouvernance & Aucune règle de quota, de budget ou de restriction de modèle \\
\hline
5 & Pas de détection d'anomalies & Les pics de consommation anormaux passent inaperçus \\
\hline
6 & Pas de prévision budgétaire & Impossible d'anticiper les coûts des mois suivants \\
\hline
7 & Pas de facturation automatisée & Les factures produites manuellement sont sources d'erreurs \\
\hline
8 & Pas d'observabilité & Aucune métrique, aucun tableau de bord pour les décideurs \\
\hline
\end{tabularx}
\end{table}


\subsection{Problématique}

La problématique centrale de ce projet peut être formulée ainsi~:

\begin{quote}
\textit{\textbf{Comment concevoir et implémenter une plateforme centralisée capable de gouverner, monitorer, analyser et optimiser l'utilisation des modèles de langage (LLM) dans un contexte multi-tenant industriel, tout en assurant la traçabilité complète, la maîtrise des coûts et la détection proactive des anomalies de consommation~?}}
\end{quote}

Cette problématique se décline en sous-questions de recherche~:

\begin{enumerate}[noitemsep]
    \item Comment intercepter et journaliser tous les appels LLM de manière transparente pour les applications clientes~?
    \item Comment appliquer des règles de gouvernance (quotas, budgets, restrictions de modèles) en temps réel~?
    \item Comment détecter automatiquement les anomalies de consommation à l'aide d'algorithmes d'apprentissage automatique~?
    \item Comment fournir des explications intelligibles des anomalies détectées via l'IA~?
    \item Comment prédire les coûts futurs et alerter sur les risques de dépassement budgétaire~?
    \item Comment automatiser la facturation par client avec des contrats flexibles~?
    \item Comment offrir un tableau de bord exécutif interactif et en temps réel~?
\end{enumerate}


\subsection{Solution proposée}

Pour répondre à la problématique identifiée, nous proposons le développement de \textbf{LLM Dashboard}, une plateforme SaaS de gouvernance et de monitoring des modèles de langage. Cette solution s'articule autour de \textbf{onze modules fonctionnels} interconnectés, formant un écosystème cohérent pour la gestion complète du cycle de vie des appels LLM~:

\begin{table}[H]
\centering
\caption{Modules fonctionnels de la plateforme LLM Dashboard}
\label{tab:modules}
\begin{tabularx}{\textwidth}{|c|l|X|}
\hline
\textbf{N°} & \textbf{Module} & \textbf{Description} \\
\hline
M1 & Gateway & Proxy LLM intelligent avec routage, gouvernance temps réel et journalisation complète \\
\hline
M2 & Analytics & Agrégation des coûts quotidiens et mensuels, pricing versionné par modèle \\
\hline
M3 & Detection & Détection d'anomalies par ML (STL, Isolation Forest, CUSUM) \\
\hline
M4 & Explainer & Explication automatique des anomalies détectées via LLM multi-fournisseur \\
\hline
M5 & Optimizer & Recommandations d'optimisation des prompts et des modèles avec simulation ROI \\
\hline
M6 & Forecasting & Prévisions budgétaires à trois scénarios et analyse de risques \\
\hline
M7 & Auth & Authentification JWT, RBAC à 4 rôles, gestion des clés API \\
\hline
M8 & Dashboard & Tableau de bord exécutif avec WebSocket temps réel et assistant IA \\
\hline
M9 & Tracing & Suivi des sessions conversationnelles, A/B testing, alertes multi-canal \\
\hline
M10 & Billing & Facturation contractuelle, génération PDF, webhooks de notification \\
\hline
M11 & Prompts & CMS de templates de prompts versionné avec variables dynamiques \\
\hline
\end{tabularx}
\end{table}


\subsection{Objectifs du projet}

Les objectifs de ce projet sont définis selon une hiérarchie de priorités~:

\begin{enumerate}[noitemsep]
    \item \textbf{Centraliser} tous les appels LLM à travers un proxy intelligent unique~;
    \item \textbf{Journaliser} chaque interaction avec les modèles (tokens, coûts, latence, statut) dans une base de données optimisée pour les séries temporelles~;
    \item \textbf{Gouverner} l'utilisation des LLM via un moteur de règles configurables~;
    \item \textbf{Détecter} automatiquement les anomalies de consommation grâce à trois algorithmes ML complémentaires~;
    \item \textbf{Expliquer} les anomalies détectées de manière intelligible via un module d'IA dédié~;
    \item \textbf{Optimiser} les coûts en recommandant des modèles alternatifs et des restructurations de prompts~;
    \item \textbf{Prévoir} les budgets futurs avec trois scénarios (optimiste, moyen, pessimiste)~;
    \item \textbf{Facturer} automatiquement chaque tenant selon son type de contrat~;
    \item \textbf{Visualiser} toutes les métriques dans un tableau de bord interactif~;
    \item \textbf{Sécuriser} la plateforme avec un système d'authentification robuste et une isolation stricte par tenant.
\end{enumerate}


% ────────────────────────────────────────────────────────────
\section{Méthodologie de développement}
\label{sec:methodologie}
% ────────────────────────────────────────────────────────────

Le choix de la méthodologie de développement constitue une décision structurante pour la réussite d'un projet logiciel. Cette section présente la méthodologie \textbf{Scrum} que nous avons adoptée pour conduire ce projet, en justifiant ce choix par rapport aux alternatives existantes.

\subsection{Méthodologies agiles --- Contexte}

Les méthodologies agiles sont nées en réponse aux limites des approches traditionnelles de développement logiciel, dites \og en cascade \fg{} (\textit{waterfall}). Le \textit{Manifesto for Agile Software Development}, publié en 2001, a posé les fondements d'une nouvelle philosophie de développement centrée sur quatre valeurs fondamentales~: les individus et les interactions, les logiciels fonctionnels, la collaboration avec le client et la réponse au changement.

Parmi les frameworks agiles les plus répandus, on distingue Scrum, Kanban, Extreme Programming (XP) et SAFe. Pour notre projet, nous avons retenu \textbf{Scrum} en raison de sa structure claire, de son orientation vers les livrables incrémentaux et de sa capacité à encadrer un projet de taille moyenne avec une équipe restreinte.

\subsection{Le framework Scrum}

Scrum est un cadre de travail agile conçu pour le développement itératif et incrémental de produits complexes. Défini formellement par Schwaber et Sutherland~\cite{schwaber2020} dans le \textit{Scrum Guide}, ce framework repose sur trois piliers fondamentaux~: la \textbf{transparence}, l'\textbf{inspection} et l'\textbf{adaptation}.

\begin{figure}[H]
    \centering
    \fbox{\parbox{12cm}{\centering\vspace{2cm}\textit{Diagramme du processus Scrum}\vspace{2cm}}}
    \caption{Le processus Scrum~\cite{schwaber2020}}
    \label{fig:scrum_process}
\end{figure}

Le fonctionnement de Scrum s'organise autour de cycles itératifs appelés \textbf{sprints}, dont la durée est typiquement comprise entre une et quatre semaines~\cite{rubin2012}. Chaque sprint constitue une unité de travail autonome qui produit un \textbf{incrément fonctionnel} potentiellement livrable. Cette approche permet de recueillir un retour régulier des parties prenantes et d'ajuster les priorités en fonction de l'évolution des besoins.

\subsubsection{Les rôles Scrum}

Le framework Scrum définit trois rôles distincts dont les responsabilités sont clairement délimitées~\cite{schwaber2020}~:

\begin{itemize}[noitemsep]
    \item \textbf{Product Owner}~: responsable de la maximisation de la valeur du produit. Il gère et priorise le \textit{Product Backlog}, un artefact qui recense l'ensemble des fonctionnalités souhaitées~;
    \item \textbf{Scrum Master}~: garant du respect du cadre Scrum, il facilite les cérémonies et lève les obstacles rencontrés par l'équipe~;
    \item \textbf{Équipe de développement}~: un groupe auto-organisé et pluridisciplinaire chargé de la conception, du développement et des tests des incréments.
\end{itemize}

\subsubsection{Les événements Scrum}

Chaque sprint est ponctué par quatre événements formels~\cite{schwaber2020}~:

\begin{enumerate}[noitemsep]
    \item \textbf{Sprint Planning}~: réunion de planification où l'équipe sélectionne les éléments du backlog à réaliser durant le sprint~;
    \item \textbf{Daily Scrum}~: mêlée quotidienne de 15 minutes pour la synchronisation de l'équipe~;
    \item \textbf{Sprint Review}~: présentation de l'incrément fonctionnel aux parties prenantes à la fin du sprint~;
    \item \textbf{Sprint Retrospective}~: session de rétrospection pour identifier les améliorations du processus.
\end{enumerate}

\subsubsection{Les artefacts Scrum}

Le framework repose sur trois artefacts essentiels~\cite{rubin2012}~:

\begin{itemize}[noitemsep]
    \item \textbf{Product Backlog}~: liste ordonnée de l'ensemble des fonctionnalités, améliorations et corrections identifiées pour le produit~;
    \item \textbf{Sprint Backlog}~: sous-ensemble du Product Backlog sélectionné pour le sprint en cours, accompagné d'un plan de réalisation~;
    \item \textbf{Incrément}~: le résultat concret du sprint, constitué de l'ensemble des éléments du Product Backlog achevés, ajoutés aux incréments des sprints précédents.
\end{itemize}


\subsection{Justification du choix de Scrum}

Le choix de Scrum comme méthodologie de développement se justifie par plusieurs facteurs inhérents à la nature de notre projet~:

\begin{enumerate}[noitemsep]
    \item \textbf{Complexité du système}~: la plateforme LLM Dashboard comporte onze modules interconnectés. Scrum permet de décomposer cette complexité en sprints autonomes, chacun produisant un sous-ensemble fonctionnel cohérent~;
    \item \textbf{Besoins évolutifs}~: les exigences en matière de gouvernance des LLM évoluent rapidement, notamment avec l'apparition continue de nouveaux modèles et fournisseurs. L'approche itérative de Scrum permet d'intégrer ces évolutions au fil des sprints~;
    \item \textbf{Retour continu}~: les sprint reviews permettent de présenter régulièrement les fonctionnalités développées à l'entreprise d'accueil et d'ajuster les priorités~;
    \item \textbf{Taille de l'équipe}~: Scrum est particulièrement adapté aux petites équipes, ce qui correspond à la configuration de notre projet de fin d'études.
\end{enumerate}

\subsection{Application de Scrum au projet}

L'application de Scrum à notre projet s'est traduite par la définition de \textbf{trois sprints} couvrant l'intégralité du développement~:

\begin{table}[H]
\centering
\caption{Découpage du projet en sprints}
\label{tab:sprints}
\begin{tabularx}{\textwidth}{|c|l|X|}
\hline
\textbf{Sprint} & \textbf{Intitulé} & \textbf{Périmètre} \\
\hline
Sprint~1 & Infrastructure et Gateway LLM & Authentification, Gateway proxy, Gouvernance, Tableau de bord, base de données \\
\hline
Sprint~2 & Intelligence Artificielle & Détection d'anomalies, Explication IA, Prévision budgétaire, Analytics, Optimisation \\
\hline
Sprint~3 & Entreprise et Déploiement & Facturation, Traçage, Gestion des prompts, Tenants, Déploiement Docker, Tests \\
\hline
\end{tabularx}
\end{table}


% ────────────────────────────────────────────────────────────
\section{Méthodologies architecturales et d'ingénierie complémentaires}
\label{sec:methodologies_complementaires}
% ────────────────────────────────────────────────────────────

Si Scrum fournit le cadre organisationnel du développement itératif, la complexité technique de la plateforme LLM Dashboard nécessite l'adoption de méthodologies d'ingénierie complémentaires pour guider les choix architecturaux, opérationnels et liés à l'intelligence artificielle. Cette section présente les cinq approches retenues et leur rôle spécifique dans notre projet.

\subsection{Domain-Driven Design (DDD)}

Le \textbf{Domain-Driven Design}, introduit par Eric Evans~\cite{evans2003}, est une méthodologie architecturale centrée sur la modélisation fidèle du domaine métier dans le code source. Le principe fondamental consiste à aligner l'implémentation logicielle sur les concepts et la terminologie propres au domaine d'affaires, afin de produire un système maintenable et évolutif.

\textbf{Application au projet.} Notre plateforme couvre plusieurs domaines métier distincts~: l'\textit{économie des tokens} (pricing, agrégation, facturation), la \textit{gouvernance} (règles, quotas, restrictions), la \textit{détection d'anomalies} (baselines, scores ML) et l'\textit{identité multi-tenant} (utilisateurs, rôles, clés API). Le DDD nous a conduit à structurer le backend en \textbf{contextes délimités} (\textit{bounded contexts}), chacun encapsulé dans un module Python autonome avec ses propres modèles, services et repositories~:

\begin{itemize}[noitemsep]
    \item \texttt{modules/gateway/} --- Contexte \og Routage et proxy LLM \fg{}~;
    \item \texttt{modules/detection/} --- Contexte \og Analyse statistique et ML \fg{}~;
    \item \texttt{modules/billing/} --- Contexte \og Facturation contractuelle \fg{}.
\end{itemize}

Cette séparation empêche la complexité d'un domaine (par exemple la logique de calcul de factures) de polluer un autre domaine (la détection d'anomalies), et facilite l'ajout futur de nouveaux contextes à mesure que la réglementation de l'IA évolue.

\subsection{Platform Engineering}

Le \textbf{Platform Engineering} constitue l'évolution moderne du DevOps, centrée sur la création d'une \og plateforme interne de développement \fg{} (\textit{Internal Developer Platform --- IDP}) offrant un \og chemin doré \fg{} (\textit{golden path}) aux développeurs consommateurs~\cite{ford2017}.

\textbf{Application au projet.} Le LLM Dashboard agit précisément comme une \textbf{plateforme de services} pour les développeurs de QALITAS et GMAO PRO. Plutôt que de gérer eux-mêmes leurs clés API, leur monitoring et leur gouvernance, les équipes de développement de TIM Tunisie consomment un service centralisé et auto-administrable~:

\begin{itemize}[noitemsep]
    \item \textbf{Self-service}~: les administrateurs de tenants créent et gèrent leurs clés API, leurs règles de gouvernance et leurs contrats de facturation de manière autonome via le tableau de bord~;
    \item \textbf{Abstraction de la complexité}~: les agents applicatifs n'ont qu'un seul endpoint à consommer (\texttt{/v1/gateway/chat/completions}), indépendamment du fournisseur LLM sous-jacent~;
    \item \textbf{Observabilité intégrée}~: les métriques, les alertes et les rapports sont produits automatiquement, sans intervention des équipes consommatrices.
\end{itemize}


\subsection{Site Reliability Engineering (SRE)}

Le \textbf{Site Reliability Engineering}, formalisé par Google~\cite{beyer2016}, est une approche d'ingénierie logicielle appliquée aux opérations. Elle repose sur la définition de \textbf{Service Level Objectives} (SLO) et la gestion de budgets d'erreur, en opposition à l'objectif irréaliste de 100\,\% de disponibilité.

\textbf{Application au projet.} La plateforme LLM Dashboard joue un rôle critique dans la chaîne applicative de TIM Tunisie~: si le proxy Gateway est indisponible, \textbf{aucune application cliente ne peut accéder aux LLM}. Les principes SRE ont guidé plusieurs décisions architecturales~:

\begin{itemize}[noitemsep]
    \item \textbf{Métriques de latence}~: le Gateway mesure et journalise la latence de bout en bout (\textit{Time to First Token --- TTFT}) pour chaque appel~;
    \item \textbf{Health checks automatisés}~: chaque service Docker expose un endpoint \texttt{/health} interrogé par le scheduler~;
    \item \textbf{Circuit breaker}~: le mécanisme de fallback multi-fournisseur (Groq $\rightarrow$ OpenAI $\rightarrow$ Anthropic) garantit la résilience face aux pannes d'un fournisseur~;
    \item \textbf{Monitoring temps réel}~: les WebSockets du tableau de bord fournissent une visibilité instantanée sur les quatre signaux dorés (\textit{golden signals})~: latence, trafic, erreurs et saturation.
\end{itemize}


\subsection{MLOps}

Le \textbf{MLOps} (\textit{Machine Learning Operations}) intègre les pratiques DevOps à l'ingénierie du Machine Learning pour traiter les modèles comme des artefacts vivants nécessitant un monitoring et une mise à jour continus~\cite{sculley2015}.

\textbf{Application au projet.} La plateforme embarque trois algorithmes d'apprentissage automatique (STL, Isolation Forest, CUSUM) exécutés automatiquement par le pipeline Celery Beat. L'approche MLOps se manifeste dans~:

\begin{itemize}[noitemsep]
    \item \textbf{Monitoring continu des modèles}~: les baselines statistiques (\texttt{llm\_token\_baseline}) sont recalculées quotidiennement pour s'adapter à l'évolution naturelle des patterns de consommation, prévenant ainsi la dérive des données (\textit{data drift})~;
    \item \textbf{Pipeline automatisé}~: l'agrégation quotidienne, la détection d'anomalies et la génération de prévisions sont orchestrées par Celery Beat sans intervention humaine~;
    \item \textbf{Traçabilité}~: chaque anomalie détectée est persistée avec l'algorithme utilisé, le z-score calculé et les valeurs attendue/réelle, assurant la reproductibilité des résultats~;
    \item \textbf{Boucle de rétroaction}~: les explications IA générées par le module Explainer enrichissent la compréhension des anomalies et orientent l'ajustement des seuils de détection.
\end{itemize}


\subsection{GitOps}

Le \textbf{GitOps} est une méthodologie opérationnelle moderne dans laquelle le dépôt Git constitue la \og source de vérité unique \fg{} pour l'ensemble de l'infrastructure et de la configuration applicative~\cite{limoncelli2014}. Tout changement d'infrastructure ou de configuration est versionné, auditable et déployé de manière automatisée.

\textbf{Application au projet.} La gouvernance des LLM implique des configurations sensibles (règles de quota, limites budgétaires, restrictions de modèles) dont la traçabilité est essentielle. Notre approche GitOps se concrétise par~:

\begin{itemize}[noitemsep]
    \item \textbf{Infrastructure as Code}~: la configuration Docker Compose (6 services), les fichiers d'environnement et les scripts de migration sont versionnés dans le dépôt Git~;
    \item \textbf{CI/CD automatisé}~: le pipeline GitHub Actions exécute les tests, construit les images Docker et les déploie automatiquement~;
    \item \textbf{Piste d'audit}~: chaque modification de configuration est associée à un commit Git horodaté, ce qui est essentiel pour la conformité réglementaire du module de gouvernance.
\end{itemize}


\subsection{Synthèse comparative des méthodologies}

Le tableau~\ref{tab:methodologies_comparatif} résume les méthodologies adoptées et leur champ d'application dans notre projet.

\begin{table}[H]
\centering
\caption{Synthèse comparative des méthodologies adoptées}
\label{tab:methodologies_comparatif}
\begin{tabularx}{\textwidth}{|l|l|X|}
\hline
\textbf{Méthodologie} & \textbf{Focus principal} & \textbf{Application dans le projet} \\
\hline
Scrum & Vélocité et itération & Cadre organisationnel, livraison incrémentale UI/Backend \\
\hline
DDD & Complexité logique & Structure modulaire (Gateway, Billing, Detection) \\
\hline
SRE & Fiabilité et performance & Proxy Gateway critique, fallback multi-fournisseur \\
\hline
MLOps & Cycle de vie des modèles ML & Pipeline de détection d'anomalies, recalcul des baselines \\
\hline
Platform Eng. & Expérience développeur & Dashboard self-service pour les tenants \\
\hline
GitOps & Audit de configuration & Versionnement Docker, CI/CD, traçabilité gouvernance \\
\hline
\end{tabularx}
\end{table}


% ────────────────────────────────────────────────────────────
\section{Langage de modélisation~: UML}
\label{sec:uml}
% ────────────────────────────────────────────────────────────

Le \textbf{Unified Modeling Language} (UML) est le standard de modélisation graphique adopté pour la conception de notre système. Défini par l'Object Management Group (OMG), UML offre un ensemble de diagrammes complémentaires permettant de représenter la structure statique et le comportement dynamique d'un système logiciel~\cite{gamma1994}.

Dans le cadre de ce projet, nous utilisons les types de diagrammes suivants~:

\begin{itemize}[noitemsep]
    \item \textbf{Diagrammes de cas d'utilisation}~: pour identifier les acteurs et leurs interactions avec le système (global et par sprint)~;
    \item \textbf{Diagrammes de classes}~: pour modéliser la structure des entités métier et leurs relations~;
    \item \textbf{Diagrammes de séquence}~: pour détailler les flux d'interaction entre les composants du système~;
    \item \textbf{Diagrammes de composants}~: pour représenter l'architecture modulaire du backend~;
    \item \textbf{Diagrammes d'activité}~: pour illustrer les processus asynchrones (pipeline Celery)~;
    \item \textbf{Diagrammes d'états}~: pour modéliser les cycles de vie des entités (contrats, factures)~;
    \item \textbf{Diagrammes de déploiement}~: pour documenter l'architecture d'infrastructure Docker.
\end{itemize}

Les diagrammes ont été produits à l'aide de l'outil \textbf{PlantUML}, qui permet de générer des représentations UML à partir de descriptions textuelles versionnées dans le dépôt Git du projet.


\section*{Conclusion}

Ce premier chapitre a posé le cadre général de notre projet de fin d'études. La présentation de TIM Tunisie a mis en lumière le contexte industriel dans lequel s'inscrit notre travail, tandis que l'analyse de l'existant a révélé des lacunes critiques dans la gestion des LLM. La problématique formulée a guidé la définition d'une solution ambitieuse --- la plateforme LLM Dashboard --- dont les onze modules fonctionnels répondent de manière exhaustive aux besoins identifiés. La méthodologie Scrum, combinée aux approches architecturales complémentaires (DDD, SRE, MLOps, Platform Engineering, GitOps), garantit un développement à la fois agile, fiable et maintenable. Le chapitre suivant dresse un état de l'art des technologies fondamentales qui sous-tendent cette plateforme.


% Chapitre 2 — État de l'Art
% ============================================================
% CHAPITRE 2 — ÉTAT DE L'ART
% ============================================================

\chapter{État de l'Art}

\section*{Plan}

\begin{tabular}{r l r}
\textbf{1} & \textbf{Les modèles de langage de grande taille (LLM)} & \textbf{\pageref{sec:llm}} \\
\textbf{2} & \textbf{Gouvernance et observabilité des LLM} & \textbf{\pageref{sec:gouvernance}} \\
\textbf{3} & \textbf{Architecture API Gateway et pattern Reverse Proxy} & \textbf{\pageref{sec:apigateway}} \\
\textbf{4} & \textbf{Séries temporelles et bases de données temporelles} & \textbf{\pageref{sec:timeseries}} \\
\textbf{5} & \textbf{Cache sémantique et embeddings vectoriels} & \textbf{\pageref{sec:embeddings}} \\
\textbf{6} & \textbf{Architectures multi-tenant SaaS} & \textbf{\pageref{sec:multitenant}} \\
\textbf{7} & \textbf{Sécurité des API~: JWT, hachage et bonnes pratiques} & \textbf{\pageref{sec:securite}} \\
\textbf{8} & \textbf{Détection d'anomalies par apprentissage automatique} & \textbf{\pageref{sec:anomaly}} \\
\textbf{9} & \textbf{Prévision de séries temporelles} & \textbf{\pageref{sec:forecasting}} \\
\textbf{10} & \textbf{Architecture événementielle et tâches asynchrones} & \textbf{\pageref{sec:eventdriven}} \\
\textbf{11} & \textbf{Conteneurisation et intégration continue} & \textbf{\pageref{sec:conteneurisation}} \\
\textbf{12} & \textbf{Solutions existantes et analyse comparative} & \textbf{\pageref{sec:comparatif}} \\
\end{tabular}

\section*{Introduction}

Ce chapitre dresse un panorama exhaustif des fondements théoriques et technologiques qui sous-tendent notre plateforme de gouvernance des LLM. Nous explorons successivement les modèles de langage de grande taille et l'architecture Transformer qui les rend possibles, les enjeux de gouvernance et d'observabilité en production, l'architecture des API Gateways et des reverse proxies, les systèmes de gestion de séries temporelles, les embeddings vectoriels et le cache sémantique, les architectures multi-tenant SaaS, la sécurité des API par JWT et hachage cryptographique, les algorithmes de détection d'anomalies par apprentissage automatique, les méthodes de prévision de séries temporelles, les paradigmes événementiels et les tâches asynchrones, la conteneurisation et l'intégration continue, et enfin les solutions existantes sur le marché. Cette revue de la littérature, s'appuyant sur des sources scientifiques et des standards industriels, permet de situer notre contribution par rapport à l'état actuel des connaissances et de justifier rigoureusement chacun des choix techniques opérés dans les sprints de développement.

% ────────────────────────────────────────────────────────────
\section{Les modèles de langage de grande taille (LLM)}
\label{sec:llm}
% ────────────────────────────────────────────────────────────

\subsection{Définition et principes fondamentaux}

Un modèle de langage de grande taille (\textit{Large Language Model}, LLM) est un réseau de neurones profond pré-entraîné sur des corpus textuels massifs, capable de comprendre et de générer du langage naturel avec une fluidité remarquable. Ces modèles reposent sur l'apprentissage de distributions de probabilité sur des séquences de tokens --- unités élémentaires de texte correspondant à des mots, des sous-mots ou des caractères selon le schéma de tokenisation employé~\cite{brown2020}.

La puissance des LLM réside dans leur capacité à capturer des relations sémantiques complexes et à effectuer du \textbf{raisonnement en contexte} (\textit{in-context learning}), c'est-à-dire à s'adapter à de nouvelles tâches sans réentraînement, uniquement à partir d'exemples fournis dans le prompt~\cite{brown2020}. Cette propriété émergente, observée à partir d'une certaine échelle de paramètres, a transformé le paradigme de l'intelligence artificielle appliquée.

\subsection{L'architecture Transformer}

L'architecture \textbf{Transformer}, introduite par Vaswani \textit{et al.} en 2017~\cite{vaswani2017}, constitue le socle technique de tous les LLM modernes. Contrairement aux architectures récurrentes (RNN, LSTM) qui traitent les séquences de manière séquentielle, le Transformer s'appuie sur un mécanisme d'\textbf{attention multi-tête} (\textit{multi-head self-attention}) qui permet de modéliser les dépendances à longue distance de manière parallélisée.

Le mécanisme d'attention peut être formalisé comme suit~\cite{vaswani2017}~:

\begin{equation}
\text{Attention}(Q, K, V) = \text{softmax}\left(\frac{QK^T}{\sqrt{d_k}}\right)V
\end{equation}

\noindent où $Q$ (Query), $K$ (Key) et $V$ (Value) sont des projections linéaires de l'entrée, et $d_k$ est la dimension des clés, servant de facteur de normalisation. Cette formulation permet à chaque position de la séquence d'accéder directement à toutes les autres positions, éliminant ainsi le goulot d'étranglement séquentiel inhérent aux architectures récurrentes.

L'architecture Transformer se compose de deux blocs principaux~:

\begin{itemize}[noitemsep]
    \item l'\textbf{encodeur}~: transforme la séquence d'entrée en une représentation contextuelle dense~;
    \item le \textbf{décodeur}~: génère la séquence de sortie token par token, en s'appuyant sur la représentation produite par l'encodeur et sur les tokens déjà générés.
\end{itemize}

Les LLM modernes adoptent généralement une architecture de type \textbf{décodeur seul} (\textit{decoder-only}), où le modèle est entraîné de manière auto-régressive à prédire le token suivant dans une séquence~\cite{brown2020}.

\begin{figure}[H]
    \centering
    \fbox{\parbox{10cm}{\centering\vspace{3cm}\textit{Schéma de l'architecture Transformer}\vspace{3cm}}}
    \caption{Architecture Transformer originale~\cite{vaswani2017}}
    \label{fig:transformer}
\end{figure}

\subsection{Évolution historique des LLM}

L'évolution des modèles de langage peut être retracée à travers plusieurs jalons fondamentaux~:

\begin{itemize}[noitemsep]
    \item \textbf{GPT-1} (2018, OpenAI)~: premier modèle démontrant l'efficacité du pré-entraînement non supervisé suivi d'un fine-tuning supervisé, avec 117 millions de paramètres~;
    \item \textbf{BERT} (2018, Google)~: modèle bidirectionnel exploitant l'encodeur du Transformer, révolutionnant la compréhension du langage naturel~;
    \item \textbf{GPT-3} (2020, OpenAI)~: avec 175 milliards de paramètres, ce modèle a démontré des capacités émergentes de \textit{few-shot learning}~\cite{brown2020}~;
    \item \textbf{ChatGPT} (2022, OpenAI)~: fine-tuné par apprentissage par renforcement à partir de retours humains (RLHF), il a démocratisé l'usage des LLM~;
    \item \textbf{GPT-4} (2023, OpenAI)~: modèle multimodal aux performances proches de l'expertise humaine sur de nombreux benchmarks~\cite{openai2023}~;
    \item \textbf{LLaMA} (2023, Meta)~: famille de modèles open-source performants, contribuant à la démocratisation de la recherche~\cite{touvron2023}~;
    \item \textbf{Claude} (2023--2024, Anthropic)~: modèle axé sur la sécurité et l'alignement, avec des fenêtres de contexte étendues~\cite{anthropic2024}.
\end{itemize}

\subsection{Tokenisation et facturation}

La \textbf{tokenisation} est le processus de décomposition du texte en unités élémentaires appelées \textbf{tokens}. La plupart des LLM modernes utilisent des algorithmes de tokenisation sous-mot tels que \textbf{BPE} (\textit{Byte Pair Encoding}) ou \textbf{SentencePiece}. La facturation des services LLM s'effectue au token, avec des tarifs différenciés entre les tokens d'entrée (\textit{input tokens}) et les tokens de sortie (\textit{output tokens}). Cette structure tarifaire rend la maîtrise de la consommation de tokens essentielle pour les entreprises, justifiant l'existence d'une plateforme de monitoring dédiée.


% ────────────────────────────────────────────────────────────
\section{Gouvernance et observabilité des LLM}
\label{sec:gouvernance}
% ────────────────────────────────────────────────────────────

\subsection{Les défis de la production LLM}

Le déploiement de modèles de langage en production soulève des défis spécifiques que Bommasani \textit{et al.}~\cite{bommasani2022} et Chen \textit{et al.}~\cite{chen2024llmops} regroupent sous le terme de \textbf{LLMOps} (\textit{Large Language Model Operations}). Ces défis couvrent~:

\begin{itemize}[noitemsep]
    \item la \textbf{maîtrise des coûts}~: la facturation au token peut engendrer des coûts exponentiels si la consommation n'est pas surveillée et plafonnée~;
    \item la \textbf{traçabilité}~: chaque appel LLM doit être journalisé avec ses métadonnées (modèle utilisé, tokens consommés, latence, statut) pour l'audit et l'optimisation~;
    \item la \textbf{gouvernance}~: des règles de contrôle d'accès, de quotas et de budgets doivent être appliquées en temps réel~;
    \item la \textbf{sécurité}~: la protection des clés API, la prévention des injections de prompts et l'isolation des données entre tenants sont des impératifs~;
    \item l'\textbf{observabilité}~: la mise à disposition de métriques, de tableaux de bord et d'alertes est nécessaire pour une exploitation sereine.
\end{itemize}

\subsection{Le concept de proxy LLM}

Un \textbf{proxy LLM} est un composant architectural positionné entre les applications clientes et les fournisseurs de modèles. Il agit comme un intermédiaire transparent qui intercepte les requêtes, applique des politiques de gouvernance, journalise les interactions et route les appels vers le fournisseur approprié. Ce pattern architectural, inspiré des API Gateways traditionnelles~\cite{fielding2000}, constitue le c\oe ur de notre plateforme.

\begin{figure}[H]
    \centering
    \fbox{\parbox{12cm}{\centering\vspace{3cm}\textit{Architecture d'un proxy LLM avec gouvernance}\vspace{3cm}}}
    \caption{Principe du proxy LLM intelligent}
    \label{fig:llm_proxy}
\end{figure}


% ────────────────────────────────────────────────────────────
\section{Architecture API Gateway et pattern Reverse Proxy}
\label{sec:apigateway}
% ────────────────────────────────────────────────────────────

\subsection{Définition et rôle architectural}

Dans le contexte des architectures distribuées modernes, l'\textbf{API Gateway} désigne un composant qui sert de point d'entrée unique pour l'ensemble des requêtes adressées à un système composé de multiples services~\cite{richardson2018}. Richardson définit ce pattern comme un intermédiaire qui découple les clients de la topologie interne du système, en assurant des fonctions transversales telles que le routage, la composition de réponses, l'authentification et la limitation de débit.

Ce concept s'inscrit dans la lignée du pattern \textbf{Reverse Proxy}, formalisé dès les premières spécifications HTTP~\cite{fielding2000}. Contrairement à un proxy classique qui agit au nom du client, un reverse proxy se positionne côté serveur et reçoit les requêtes entrantes pour les redistribuer vers les services appropriés, tout en masquant l'infrastructure sous-jacente. Nginx, utilisé dans notre plateforme, est l'implémentation de référence de ce pattern, offrant des fonctionnalités natives de terminaison TLS, de compression et de \textit{rate limiting}~\cite{nginx2024}.

\subsection{Patterns de résilience}

Les systèmes distribués sont intrinsèquement exposés aux défaillances partielles. Pour y répondre, la littérature en ingénierie logicielle a formalisé plusieurs patterns de résilience que Richardson~\cite{richardson2018} et Nygard décrivent dans le contexte des microservices~:

\begin{itemize}[noitemsep]
    \item le \textbf{Circuit Breaker}~: inspiré de l'électrotechnique, ce pattern surveille le taux d'erreur d'un service aval et <<~ouvre le circuit~>> temporairement lorsqu'un seuil est dépassé, évitant ainsi la propagation en cascade des défaillances~;
    \item le \textbf{Rate Limiting}~: mécanisme de limitation du nombre de requêtes par unité de temps, typiquement implémenté par un compteur à fenêtre glissante dans un stockage en mémoire tel que Redis~;
    \item le \textbf{Retry with Backoff}~: stratégie de réessai avec délai exponentiel croissant, permettant de gérer les erreurs transitoires sans surcharger le service défaillant~;
    \item le \textbf{Bulkhead}~: isolation des ressources par domaine fonctionnel, empêchant qu'un composant défaillant ne monopolise l'ensemble des ressources du système.
\end{itemize}

\subsection{Application à notre plateforme}

Notre plateforme implémente le pattern API Gateway à deux niveaux. Le premier niveau est assuré par \textbf{Nginx}, qui agit comme reverse proxy en entrée de l'infrastructure, gérant la terminaison TLS, le \textit{rate limiting} global et le routage vers les services internes. Le second niveau est le \textbf{Gateway LLM}, un composant applicatif développé en FastAPI, qui intercepte chaque appel destiné aux fournisseurs LLM pour y appliquer la gouvernance, la journalisation et le cache sémantique avant de router la requête vers le provider sélectionné.


% ────────────────────────────────────────────────────────────
\section{Séries temporelles et bases de données temporelles}
\label{sec:timeseries}
% ────────────────────────────────────────────────────────────

\subsection{Définition et propriétés}

Une \textbf{série temporelle} est une séquence ordonnée de mesures indexées par le temps, formellement représentée par $\{(t_i, v_i)\}_{i=1}^{n}$ où $t_i$ désigne l'instant d'observation et $v_i$ la valeur mesurée~\cite{jensen2017}. Les séries temporelles possèdent des propriétés caractéristiques qui les distinguent des données relationnelles classiques~:

\begin{itemize}[noitemsep]
    \item l'\textbf{append-only}~: les nouvelles données sont exclusivement insérées en fin de série, les mises à jour étant exceptionnelles~;
    \item la \textbf{cardinalité temporelle}~: le volume croît linéairement avec le temps, pouvant atteindre plusieurs milliards d'enregistrements~;
    \item la \textbf{dépendance temporelle}~: les observations consécutives sont généralement corrélées, propriété exploitée par les algorithmes de détection d'anomalies~;
    \item la \textbf{granularité variable}~: les données sont souvent agrégées à différentes résolutions (minute, heure, jour, mois) pour des besoins analytiques distincts.
\end{itemize}

\subsection{Systèmes de gestion de séries temporelles}

Jensen, Pedersen et Thomsen~\cite{jensen2017} proposent une taxonomie des systèmes de gestion de séries temporelles (\textit{Time Series Management Systems}, TSMS) en distinguant trois grandes familles~:

\begin{table}[H]
\centering
\caption{Comparaison des approches de stockage de séries temporelles}
\label{tab:tsms_comparaison}
\begin{tabularx}{\textwidth}{|l|X|X|l|}
\hline
\textbf{Approche} & \textbf{Principe} & \textbf{Exemple} & \textbf{Cas d'usage} \\
\hline
SGBDR étendu & Extension d'un SGBDR classique avec des types et opérateurs temporels & TimescaleDB & Analytics, facturation \\
\hline
Base native & Base conçue spécifiquement pour les séries temporelles & InfluxDB & Monitoring infra \\
\hline
Stockage clé-valeur & Paires (timestamp, valeur) dans un stockage distribué & Prometheus & Métriques système \\
\hline
\end{tabularx}
\end{table}

\subsection{TimescaleDB et le concept de hypertable}

\textbf{TimescaleDB} est une extension open-source de PostgreSQL qui transforme une table relationnelle classique en une \textbf{hypertable}~\cite{timescaledb2024}. Une hypertable est une abstraction qui partitionne automatiquement les données en \textit{chunks} selon la dimension temporelle, offrant des performances d'insertion et de requête adaptées aux séries temporelles tout en préservant la compatibilité SQL complète de PostgreSQL.

Les avantages principaux de cette approche sont~:

\begin{itemize}[noitemsep]
    \item la \textbf{compression automatique}~: les chunks anciens sont compressés avec un ratio pouvant atteindre 95\,\%, réduisant considérablement l'empreinte disque~;
    \item les \textbf{agrégations continues} (\textit{continuous aggregates})~: des vues matérialisées qui se mettent à jour automatiquement à chaque nouvel enregistrement~;
    \item le \textbf{partitionnement transparent}~: les requêtes SQL standard bénéficient du partitionnement sans modification du code applicatif~;
    \item la \textbf{rétention de données}~: des politiques de suppression automatique des données anciennes, essentielles pour la conformité réglementaire.
\end{itemize}

Dans notre plateforme, la table \texttt{llm\_token\_log} --- qui enregistre chaque appel LLM avec ses métadonnées --- est configurée en hypertable partitionnée par la colonne \texttt{timestamp}, permettant de gérer efficacement un flux continu d'événements à haute fréquence tout en maintenant des performances de requête compatibles avec les tableaux de bord temps réel.


% ────────────────────────────────────────────────────────────
\section{Cache sémantique et embeddings vectoriels}
\label{sec:embeddings}
% ────────────────────────────────────────────────────────────

\subsection{Représentations vectorielles distribuées}

Les \textbf{embeddings vectoriels} constituent une avancée fondamentale en traitement du langage naturel. Mikolov \textit{et al.}~\cite{mikolov2013} ont démontré qu'il est possible de représenter des mots, des phrases ou des documents sous forme de vecteurs denses dans un espace de dimension réduite ($\mathbb{R}^d$, typiquement $d \in [128, 1536]$), de sorte que la proximité géométrique entre vecteurs reflète la similarité sémantique entre les entités représentées.

La \textbf{similarité cosinus} est la mesure privilégiée pour comparer deux embeddings $\mathbf{a}$ et $\mathbf{b}$~:

\begin{equation}
\text{sim}(\mathbf{a}, \mathbf{b}) = \frac{\mathbf{a} \cdot \mathbf{b}}{\|\mathbf{a}\| \cdot \|\mathbf{b}\|} = \frac{\sum_{i=1}^{d} a_i b_i}{\sqrt{\sum_{i=1}^{d} a_i^2} \cdot \sqrt{\sum_{i=1}^{d} b_i^2}}
\end{equation}

\noindent Une valeur de 1 indique une identité sémantique parfaite, tandis qu'une valeur de 0 signale l'absence de relation sémantique. Les modèles de type \textbf{Sentence-BERT} ou \textbf{text-embedding} des fournisseurs LLM étendent ce principe à des séquences entières de texte.

\subsection{Recherche par plus proches voisins}

La recherche de similarité dans des espaces vectoriels de grande dimension pose un défi algorithmique majeur. Johnson, Douze et Jégou~\cite{johnson2019} ont proposé des approches de recherche approximative de plus proches voisins (\textit{Approximate Nearest Neighbor}, ANN) accélérées par GPU, permettant de parcourir des milliards de vecteurs en temps quasi-constant. Les structures d'index les plus courantes sont~:

\begin{itemize}[noitemsep]
    \item \textbf{IVF} (\textit{Inverted File Index})~: partitionnement de l'espace en cellules de Voronoï, permettant de limiter la recherche aux cellules voisines du vecteur requête~;
    \item \textbf{HNSW} (\textit{Hierarchical Navigable Small World})~: graphe de navigation hiérarchique offrant un excellent compromis précision/vitesse pour les requêtes en ligne~;
    \item \textbf{Product Quantization}~: compression des vecteurs en sous-quantificateurs, réduisant l'empreinte mémoire tout en préservant la précision de recherche.
\end{itemize}

\subsection{Bases de données vectorielles}

L'essor des LLM a engendré une nouvelle catégorie de systèmes de stockage~: les \textbf{bases de données vectorielles}. Notre plateforme utilise \textbf{pgvector}, une extension PostgreSQL qui ajoute un type natif \texttt{VECTOR(n)} et des opérateurs de recherche par similarité directement dans le SGBDR existant. Cette approche offre l'avantage de ne pas introduire un composant d'infrastructure supplémentaire, les embeddings cohabitant avec les données relationnelles dans une même base.

\subsection{Le cache sémantique}

Le \textbf{cache sémantique} représente une évolution du cache exact traditionnel. Là où un cache clé-valeur classique ne peut retourner un résultat que pour une requête strictement identique à une requête déjà mise en cache, un cache sémantique exploite les embeddings vectoriels pour identifier les requêtes \textit{sémantiquement proches} d'une requête déjà traitée. Si la similarité cosinus entre le vecteur de la requête entrante et celui d'une entrée en cache dépasse un seuil prédéfini (typiquement $\theta \geq 0{,}95$), la réponse mise en cache est retournée directement, économisant à la fois la latence et le coût de l'appel LLM. Dans notre plateforme, ce mécanisme est intégré au Gateway et permet de réduire significativement les coûts pour les prompts récurrents au sein d'un même tenant.


% ────────────────────────────────────────────────────────────
\section{Architectures multi-tenant SaaS}
\label{sec:multitenant}
% ────────────────────────────────────────────────────────────

\subsection{Définition du multi-tenant}

Une architecture \textbf{multi-tenant} (\textit{multi-locataire}) est un modèle de conception logicielle dans lequel une seule instance d'une application dessert simultanément plusieurs organisations clientes, appelées \textbf{tenants}, tout en garantissant l'isolation stricte de leurs données respectives~\cite{bezemer2010}. Ce modèle est au c\oe ur des solutions SaaS (\textit{Software as a Service}) modernes.

Bezemer et Zaidman~\cite{bezemer2010} identifient trois niveaux d'isolation dans une architecture multi-tenant~:

\begin{enumerate}[noitemsep]
    \item \textbf{Base de données partagée, schéma partagé}~: tous les tenants partagent les mêmes tables, la séparation étant assurée par une colonne \texttt{tenant\_id}~;
    \item \textbf{Base de données partagée, schéma séparé}~: chaque tenant dispose de son propre schéma au sein de la même instance de base de données~;
    \item \textbf{Base de données séparée}~: chaque tenant dispose de sa propre instance de base de données.
\end{enumerate}

Dans notre plateforme, nous avons adopté le premier niveau --- base de données partagée avec colonne \texttt{tenant\_id} --- pour sa simplicité de gestion et sa scalabilité, tout en implémentant une isolation stricte à trois niveaux (API, service, base de données).

\subsection{Contrôle d'accès basé sur les rôles (RBAC)}

Le modèle RBAC (\textit{Role-Based Access Control}) est un mécanisme de contrôle d'accès qui attribue des permissions non pas directement aux utilisateurs, mais à des rôles que les utilisateurs endossent. Conformément aux recommandations de l'OWASP~\cite{owasp2021}, notre plateforme implémente un RBAC à quatre niveaux hiérarchiques~: Super Administrateur, Administrateur Tenant, Observateur et Agent API.


% ────────────────────────────────────────────────────────────
\section{Sécurité des API~: JWT, hachage et bonnes pratiques}
\label{sec:securite}
% ────────────────────────────────────────────────────────────

La sécurisation des interfaces programmatiques constitue un enjeu critique pour toute plateforme SaaS exposée sur Internet. Cette section présente les fondements théoriques des mécanismes d'authentification et de protection déployés dans notre système.

\subsection{JSON Web Tokens (JWT)}

Le standard \textbf{JWT} (\textit{JSON Web Token}), formalisé dans la RFC~7519~\cite{rfc7519}, définit un format compact et auto-contenu pour transmettre des assertions de sécurité entre deux parties. Un JWT se compose de trois segments encodés en Base64URL et séparés par des points~:

\begin{equation}
\texttt{JWT} = \texttt{Base64(Header)}.\texttt{Base64(Payload)}.\texttt{Base64(Signature)}
\end{equation}

\begin{itemize}[noitemsep]
    \item le \textbf{Header} spécifie l'algorithme de signature (HS256, RS256) et le type de token~;
    \item le \textbf{Payload} contient les \textit{claims} --- des assertions telles que l'identité de l'utilisateur (\texttt{sub}), la date d'expiration (\texttt{exp}), le rôle (\texttt{role}) et l'identifiant du tenant (\texttt{tenant\_id})~;
    \item la \textbf{Signature} garantit l'intégrité du token par un calcul HMAC ou RSA appliqué aux deux premiers segments.
\end{itemize}

L'architecture \textbf{access/refresh token} est un pattern complémentaire recommandé par le cadre OAuth~2.0~\cite{oauth2012}. L'access token, de courte durée de vie (typiquement 15 à 30 minutes), autorise l'accès aux ressources protégées, tandis que le refresh token, de plus longue durée, permet d'obtenir un nouvel access token sans réauthentification. Cette séparation limite la fenêtre d'exposition en cas de compromission d'un access token.

\subsection{Hachage sécurisé des mots de passe}

Le stockage des mots de passe en clair constitue une vulnérabilité critique classée dans le Top~10 OWASP~\cite{owasp2021}. Les algorithmes de \textbf{hachage adaptatif} répondent à cette menace en produisant une empreinte irréversible dont le coût de calcul est volontairement élevé, rendant les attaques par force brute économiquement prohibitives.

L'algorithme \textbf{PBKDF2} (\textit{Password-Based Key Derivation Function 2}), spécifié dans la RFC~2898~\cite{rfc2898}, applique une fonction pseudo-aléatoire (typiquement HMAC-SHA256) de manière itérative~:

\begin{equation}
\texttt{DK} = \text{PBKDF2}(\texttt{PRF}, \texttt{Password}, \texttt{Salt}, c, \texttt{dkLen})
\end{equation}

\noindent où \texttt{PRF} est la fonction pseudo-aléatoire, \texttt{Salt} un sel aléatoire unique par utilisateur, $c$ le nombre d'itérations (recommandation OWASP~: $c \geq 600\,000$ pour SHA-256) et \texttt{dkLen} la longueur de la clé dérivée. Le sel empêche les attaques par tables arc-en-ciel (\textit{rainbow tables}), et le nombre élevé d'itérations impose un coût prohibitif à l'attaquant.

\subsection{Blacklisting de tokens et sécurité de session}

La révocation de tokens JWT pose un défi inhérent à leur nature stateless~: un token émis reste valide jusqu'à son expiration, même si l'utilisateur se déconnecte. La solution adoptée dans notre plateforme consiste à maintenir une \textbf{liste noire} (\textit{blacklist}) des tokens révoqués dans Redis, avec une clé TTL (\textit{Time To Live}) correspondant à la durée de vie résiduelle du token. Cette approche préserve la nature stateless de l'authentification tout en offrant un mécanisme de révocation immédiat, conformément aux recommandations OWASP pour les architectures API~\cite{owasp2021}.


% ────────────────────────────────────────────────────────────
\section{Détection d'anomalies par apprentissage automatique}
\label{sec:anomaly}
% ────────────────────────────────────────────────────────────

La détection d'anomalies constitue un domaine fondamental de l'apprentissage automatique. Chandola \textit{et al.}~\cite{chandola2009} définissent une anomalie comme une observation qui s'écarte significativement du comportement attendu dans un jeu de données. Dans le contexte de notre plateforme, les anomalies correspondent à des pics ou des creux inhabituels de consommation de tokens.

Notre système implémente \textbf{trois algorithmes complémentaires} pour maximiser la couverture de détection~:

\subsection{Décomposition STL}

La décomposition STL (\textit{Seasonal and Trend decomposition using Loess}), proposée par Cleveland \textit{et al.}~\cite{cleveland1990}, est une technique de décomposition des séries temporelles en trois composantes~:

\begin{equation}
Y_t = T_t + S_t + R_t
\end{equation}

\noindent où $Y_t$ est la valeur observée au temps $t$, $T_t$ la composante de tendance, $S_t$ la composante saisonnière et $R_t$ le résidu. Une anomalie est détectée lorsque le résidu $R_t$ dépasse un seuil défini en termes de z-scores~:

\begin{equation}
z_t = \frac{R_t - \mu_R}{\sigma_R}
\end{equation}

\noindent Un z-score $|z_t| > 3$ indique une anomalie statistiquement significative.

\subsection{Isolation Forest}

L'\textbf{Isolation Forest}, proposé par Liu \textit{et al.}~\cite{liu2008}, est un algorithme de détection d'anomalies non supervisé particulièrement efficace pour les données à haute dimensionnalité. Son principe fondamental repose sur l'observation que les anomalies, étant rares et différentes, sont plus facilement \textbf{isolables} que les observations normales.

L'algorithme construit un ensemble d'\textbf{arbres d'isolation} (\textit{isolation trees}) en sélectionnant aléatoirement un attribut et une valeur de coupure. La profondeur moyenne à laquelle une observation est isolée dans l'ensemble des arbres constitue son \textbf{score d'anomalie}~: les observations isolées rapidement (à faible profondeur) sont considérées comme anomales~\cite{liu2008}.

Le score d'anomalie est normalisé selon la formule~:

\begin{equation}
s(x, n) = 2^{-\frac{E[h(x)]}{c(n)}}
\end{equation}

\noindent où $E[h(x)]$ est la profondeur moyenne d'isolation de l'observation $x$, et $c(n)$ est la profondeur moyenne dans un arbre binaire de recherche de $n$ éléments, servant de facteur de normalisation.

\begin{figure}[H]
    \centering
    \fbox{\parbox{10cm}{\centering\vspace{3cm}\textit{Principe de l'Isolation Forest : les anomalies sont isolées en moins de coupures}\vspace{3cm}}}
    \caption{Principe de l'Isolation Forest~\cite{liu2008}}
    \label{fig:isolation_forest}
\end{figure}

\subsection{CUSUM (Cumulative Sum Control Chart)}

L'algorithme \textbf{CUSUM} (\textit{Cumulative Sum}), proposé par Page~\cite{page1954}, est une technique de contrôle statistique de processus conçue pour détecter des changements de moyenne dans une série temporelle. Il calcule une somme cumulative des écarts par rapport à une valeur de référence~:

\begin{equation}
S_t = \max(0, S_{t-1} + (x_t - \mu_0 - k))
\end{equation}

\noindent où $x_t$ est la valeur observée, $\mu_0$ la moyenne de référence et $k$ un paramètre de sensibilité. Une alarme est déclenchée lorsque $S_t$ dépasse un seuil prédéfini $h$.

L'intérêt du CUSUM dans notre contexte réside dans sa sensibilité aux \textbf{dérives progressives} de consommation, complémentaire de l'Isolation Forest qui excelle dans la détection de points isolés.

\subsection{Paradigme d'apprentissage et pipeline de données}

Une question fondamentale en apprentissage automatique concerne la nature de la supervision~: les algorithmes de notre plateforme relèvent-ils de l'apprentissage supervisé (nécessitant des données étiquetées <<~anomalie~>> / <<~normal~>>) ou de l'apprentissage non supervisé~? Dans le contexte de la gouvernance LLM, l'obtention de telles étiquettes est impraticable~: elle exigerait qu'un expert humain qualifie rétrospectivement chaque journée de consommation comme normale ou anormale pour chaque tenant. C'est pourquoi notre système repose intégralement sur un \textbf{apprentissage non supervisé} et \textbf{statistique}~\cite{chandola2009}.

\subsubsection{Phase 1~: Apprentissage hors ligne --- le Baseline}

Le Baseline constitue le <<~profil de référence~>> d'un tenant, calculé périodiquement (chaque nuit) à partir de l'historique de consommation. Le processus d'apprentissage s'effectue en quatre étapes~:

\begin{enumerate}[noitemsep]
    \item \textbf{Extraction des données.} Les agrégats quotidiens sont récupérés depuis la table \texttt{llm\_cost\_daily} sur une fenêtre glissante de 30 jours (minimum~: 14 jours). Chaque jour produit un \textbf{vecteur de caractéristiques à 6 dimensions}~:
    \begin{equation}
    \mathbf{x}_t = \begin{bmatrix} \text{total\_tokens} \\ \text{call\_count} \\ \text{avg\_tokens\_per\_call} \\ \text{cost\_usd} \\ \text{error\_rate} \\ \text{avg\_latency\_ms} \end{bmatrix}
    \end{equation}
    Ce vecteur multi-dimensionnel est essentiel, car il permet de distinguer un pic de tokens causé par davantage d'appels d'un pic causé par des prompts individuellement plus longs --- deux situations aux implications opérationnelles radicalement différentes.

    \item \textbf{Calcul des statistiques descriptives.} La moyenne $\mu$, l'écart-type $\sigma$, le minimum et le maximum des tokens quotidiens sont calculés sur la fenêtre. Ces statistiques définissent quantitativement ce qui constitue un <<~comportement normal~>> pour ce tenant spécifique.

    \item \textbf{Apprentissage structurel implicite.} L'Isolation Forest est entraîné (\texttt{fit}) sur la matrice $\mathbf{X} \in \mathbb{R}^{30 \times 6}$ des vecteurs historiques. L'algorithme construit 100 arbres d'isolation qui <<~apprennent~>> la structure de distribution des données sans jamais recevoir d'étiquettes~: il identifie les régions denses (normalité) et les régions creuses (anomalies potentielles).

    \item \textbf{Calibration adaptative du seuil.} Plutôt qu'un seuil fixe, le système calibre un seuil propre à chaque tenant en calculant le 5\textsuperscript{e} percentile des scores de l'Isolation Forest sur les données d'entraînement elles-mêmes. Ainsi, un tenant à consommation très variable aura un seuil plus permissif qu'un tenant à consommation stable.
\end{enumerate}

Le Baseline résultant est persisté dans la table \texttt{llm\_token\_baseline} et constitue la <<~mémoire apprise~>> du système.

\subsubsection{Phase 2~: Inférence en ligne --- la Détection}

À chaque nouvel appel LLM, le système compare la consommation courante au Baseline appris~:

\begin{enumerate}[noitemsep]
    \item \textbf{STL}~: décompose les 30 derniers jours + la valeur du jour en tendance + saisonnalité + résidu. Il <<~sait~>> qu'un lundi a un profil différent d'un samedi. Seul le résidu est évalué par un z-score~; un $|z| \geq 2{,}5$ déclenche un vote d'anomalie.
    \item \textbf{Isolation Forest}~: le vecteur $\mathbf{x}_{\text{today}}$ est scoré par les arbres déjà entraînés. Un score inférieur au seuil calibré signifie que le point est <<~facile à isoler~>>, donc anormal.
    \item \textbf{CUSUM}~: met à jour deux accumulateurs persistants $C^+$ et $C^-$ avec la valeur normalisée du jour. L'état des accumulateurs est sauvegardé dans le Baseline après chaque évaluation, ce qui confère à CUSUM une \textbf{mémoire inter-jours} que les deux autres algorithmes ne possèdent pas.
\end{enumerate}

\subsubsection{Vote par ensemble majoritaire}

Les trois détecteurs émettent chacun un \textbf{vote binaire} (anomalie / normal). La décision finale suit une règle de vote majoritaire~: au moins 2 votes sur 3 sont nécessaires pour confirmer une anomalie. La sévérité est ensuite graduée selon le nombre de votes et la magnitude des scores individuels~:

\begin{table}[H]
\centering
\caption{Règle de décision de l'ensemble de détection}
\label{tab:ensemble_voting}
\begin{tabularx}{\textwidth}{|c|c|X|}
\hline
\textbf{Votes concordants} & \textbf{Sévérité} & \textbf{Interprétation} \\
\hline
0--1 / 3 & Normal & Signal faible ou bruit statistique \\
\hline
2 / 3 & Warning ou High & Anomalie probable, deux algorithmes indépendants concordent \\
\hline
3 / 3 & High ou Critical & Anomalie confirmée avec très haute confiance \\
\hline
\end{tabularx}
\end{table}

Cette architecture en deux phases --- apprentissage hors ligne du Baseline, puis inférence en ligne avec vote majoritaire --- permet au système de s'auto-adapter à chaque client sans intervention humaine et sans données étiquetées, le modèle ré-apprenant continuellement à mesure que de nouvelles données sont intégrées dans la fenêtre glissante.

\subsection{Complémentarité des trois algorithmes}

La combinaison de ces trois algorithmes permet de couvrir un spectre large de types d'anomalies~:

\begin{table}[H]
\centering
\caption{Complémentarité des algorithmes de détection d'anomalies}
\label{tab:algo_complement}
\begin{tabularx}{\textwidth}{|l|X|X|}
\hline
\textbf{Algorithme} & \textbf{Forces} & \textbf{Type d'anomalies détectées} \\
\hline
STL & Capture la saisonnalité et les tendances & Déviations par rapport aux patterns temporels \\
\hline
Isolation Forest & Non supervisé, robuste aux données à haute dimension & Points aberrants isolés, clusters anormaux \\
\hline
CUSUM & Sensible aux changements graduels de moyenne & Dérives progressives, changements de régime \\
\hline
\end{tabularx}
\end{table}


% ────────────────────────────────────────────────────────────
\section{Prévision de séries temporelles}
\label{sec:forecasting}
% ────────────────────────────────────────────────────────────

Si la détection d'anomalies vise à identifier les comportements passés aberrants, la \textbf{prévision} (\textit{forecasting}) cherche à projeter les tendances futures à partir des données historiques. Dans le contexte d'une plateforme de gouvernance LLM, la prévision budgétaire permet aux décideurs d'anticiper les dépassements de coûts et de planifier les enveloppes budgétaires allouées à chaque tenant.

\subsection{Méthodes statistiques de prévision}

Makridakis, Spiliotis et Assimakopoulos~\cite{makridakis2018} ont conduit une étude comparative exhaustive des méthodes de prévision, concluant que les approches statistiques classiques surpassent souvent les modèles d'apprentissage profond pour les séries temporelles univariées de taille modérée --- précisément le profil de données rencontré dans notre contexte de consommation quotidienne de tokens.

Les méthodes les plus établies sont~:

\begin{itemize}[noitemsep]
    \item la \textbf{régression linéaire}~: modèle de base qui estime une tendance linéaire $\hat{y}_t = \alpha + \beta t$, où $\beta$ représente le taux de croissance quotidien. Sa simplicité en fait un outil de première approximation pour les projections à court terme~;
    \item les \textbf{modèles ARIMA} (\textit{AutoRegressive Integrated Moving Average})~: modèles paramétriques qui capturent à la fois les corrélations auto-régressives et les moyennes mobiles des séries, après stationnarisation par différenciation~;
    \item le \textbf{lissage exponentiel} (Holt-Winters)~: famille de modèles qui décomposent la série en niveau, tendance et saisonnalité, chaque composante étant mise à jour de manière récursive avec un poids exponentiel décroissant sur les observations passées.
\end{itemize}

\subsection{Prévision à l'échelle et analyse de scénarios}

Taylor et Letham~\cite{taylor2018} ont proposé \textbf{Prophet}, un modèle de prévision conçu pour les séries temporelles métier présentant des tendances non linéaires et des saisonnalités multiples. Le modèle décompose la série en trois composantes additives~:

\begin{equation}
y(t) = g(t) + s(t) + h(t) + \varepsilon_t
\end{equation}

\noindent où $g(t)$ est la fonction de tendance (linéaire ou logistique), $s(t)$ la composante saisonnière modélisée par une série de Fourier, $h(t)$ l'effet des jours exceptionnels, et $\varepsilon_t$ le terme d'erreur. L'intérêt de cette approche réside dans sa capacité à produire des intervalles de confiance, naturellement exploitables pour construire des scénarios multiples.

Notre plateforme adopte une approche d'\textbf{analyse de scénarios} à trois niveaux~: un scénario \textit{optimiste} (croissance freinée), un scénario \textit{moyen} (tendance linéaire prolongée) et un scénario \textit{pessimiste} (croissance accélérée). Chaque scénario est croisé avec les plafonds contractuels du tenant pour évaluer le risque de dépassement budgétaire et déclencher des alertes préventives.


% ────────────────────────────────────────────────────────────
\section{Architecture événementielle et tâches asynchrones}
\label{sec:eventdriven}
% ────────────────────────────────────────────────────────────

\subsection{Le paradigme publish/subscribe}

L'architecture \textbf{événementielle} (\textit{event-driven architecture}) est un paradigme de conception dans lequel les composants d'un système communiquent par l'émission et la consommation d'événements, plutôt que par des appels synchrones directs. Hohpe et Woolf~\cite{hohpe2003} formalisent ce paradigme à travers le pattern \textbf{Publish/Subscribe}, où un producteur publie un message sur un canal (ou \textit{topic}) et un ou plusieurs consommateurs s'y abonnent de manière découplée.

Kleppmann~\cite{kleppmann2017} souligne que ce découplage offre trois avantages fondamentaux dans les systèmes de traitement de données~:

\begin{itemize}[noitemsep]
    \item la \textbf{résilience}~: la défaillance d'un consommateur n'affecte pas le producteur ni les autres consommateurs~;
    \item la \textbf{scalabilité}~: de nouveaux consommateurs peuvent s'abonner à un flux existant sans modifier le producteur~;
    \item l'\textbf{extensibilité}~: l'ajout d'une nouvelle fonctionnalité (par exemple, la détection d'anomalies) se réduit à la souscription à un événement existant (\texttt{llm.call.completed}).
\end{itemize}

\subsection{Tâches asynchrones et orchestration}

Les opérations dont la latence est incompatible avec une réponse HTTP synchrone --- telles que l'agrégation de coûts sur des millions d'enregistrements ou l'appel à un LLM pour expliquer une anomalie --- doivent être déléguées à des \textbf{workers asynchrones}. Le pattern \textbf{Task Queue}~\cite{hohpe2003} repose sur trois composants~:

\begin{enumerate}[noitemsep]
    \item un \textbf{producteur} qui enfile une tâche dans une file de messages~;
    \item un \textbf{broker} (Redis ou RabbitMQ) qui stocke la tâche en mémoire~;
    \item un \textbf{worker} qui consomme la tâche, l'exécute et persiste le résultat.
\end{enumerate}

\textbf{Celery} est le framework de référence en Python pour l'orchestration de tâches asynchrones distribuées~\cite{celery2024}. Son module \textbf{Celery Beat} ajoute un planificateur (\textit{scheduler}) qui déclenche automatiquement des tâches périodiques selon une expression crontab. Dans notre plateforme, Celery Beat orchestre trois tâches critiques~: l'agrégation quotidienne des coûts (chaque nuit), le rollup mensuel (chaque premier du mois) et le scan d'optimisation hebdomadaire.

\subsection{Pipelines de données}

Les pipelines de données de notre plateforme suivent le paradigme \textbf{ELT} (\textit{Extract, Load, Transform})~\cite{kleppmann2017}, où les données brutes sont d'abord chargées dans la base de données (table \texttt{llm\_token\_log}), puis transformées par des tâches Celery en agrégats exploitables (\texttt{llm\_cost\_daily}, \texttt{llm\_cost\_monthly}). Cette approche est préférable au paradigme ETL traditionnel dans un contexte cloud, car elle exploite la puissance de calcul du SGBDR pour les transformations et évite les erreurs de conversion en amont.


% ────────────────────────────────────────────────────────────
\section{Conteneurisation et intégration continue}
\label{sec:conteneurisation}
% ────────────────────────────────────────────────────────────

\subsection{Principes de la conteneurisation}

La \textbf{conteneurisation} est une technique de virtualisation au niveau du système d'exploitation qui encapsule une application et ses dépendances dans une unité d'exécution isolée appelée \textbf{conteneur}~\cite{merkel2014}. Contrairement à la virtualisation traditionnelle par hyperviseur, les conteneurs partagent le noyau de l'hôte et ne transportent que les bibliothèques strictement nécessaires, offrant un démarrage quasi instantané et une empreinte mémoire minimale.

Merkel~\cite{merkel2014} identifie trois propriétés fondamentales qui justifient l'adoption massive de Docker dans l'industrie~:

\begin{itemize}[noitemsep]
    \item la \textbf{reproductibilité}~: un conteneur construit à partir d'un \texttt{Dockerfile} produit un environnement d'exécution identique, indépendamment de la machine hôte~;
    \item l'\textbf{isolation}~: chaque conteneur dispose de son propre espace de noms (\textit{namespace}) pour les processus, le réseau et le système de fichiers~;
    \item la \textbf{portabilité}~: une image Docker peut être transférée et exécutée sur tout système supportant le runtime Docker, du poste de développement au serveur de production.
\end{itemize}

\subsection{Orchestration multi-conteneurs}

Les applications modernes se composent de multiples services interdépendants. \textbf{Docker Compose} est un outil d'orchestration qui permet de décrire, dans un fichier déclaratif YAML, l'ensemble des services, réseaux et volumes d'une application multi-conteneurs~\cite{merkel2014}. Pour les architectures de plus grande envergure nécessitant de l'auto-scaling et de la haute disponibilité, \textbf{Kubernetes} propose une orchestration de niveau production avec des mécanismes de \textit{self-healing}, de \textit{rolling updates} et de \textit{service discovery}~\cite{burns2018}.

Notre plateforme utilise Docker Compose pour orchestrer six services conteneurisés --- PostgreSQL (avec TimescaleDB et pgvector), Redis, l'API FastAPI, le worker Celery, le scheduler Celery Beat et Nginx --- au sein d'un réseau interne isolé avec des volumes persistants pour les données.

\subsection{Intégration continue et déploiement continu (CI/CD)}

L'\textbf{intégration continue} (\textit{Continuous Integration}, CI) est une pratique d'ingénierie logicielle qui consiste à fusionner fréquemment le code dans un tronc commun, chaque fusion déclenchant automatiquement une suite de vérifications~: analyse syntaxique (\textit{linting}), exécution des tests unitaires et d'intégration, et construction de l'image Docker. Le \textbf{déploiement continu} (\textit{Continuous Deployment}, CD) prolonge cette chaîne en automatisant la mise en production des artefacts validés~\cite{ford2017}.

Notre pipeline CI/CD, implémenté via \textbf{GitHub Actions}, automatise quatre étapes à chaque push~: la vérification syntaxique de tous les fichiers Python, l'exécution des 180 tests avec des services TimescaleDB et Redis éphémères, la construction et le push de l'image Docker multi-stage vers GitHub Container Registry (GHCR), et la notification de déploiement. Cette automatisation s'inscrit dans les principes \textbf{GitOps} décrits dans le chapitre~1, où l'état souhaité de l'infrastructure est entièrement versionné dans le dépôt Git.


% ────────────────────────────────────────────────────────────
\section{Solutions existantes et analyse comparative}
\label{sec:comparatif}
% ────────────────────────────────────────────────────────────

Plusieurs solutions de monitoring et de gouvernance des LLM existent sur le marché. Nous présentons ici une analyse comparative des principaux acteurs~:

\begin{table}[H]
\centering
\caption{Analyse comparative des solutions existantes}
\label{tab:comparatif}
\begin{tabularx}{\textwidth}{|p{2.2cm}|X|X|X|}
\hline
\textbf{Solution} & \textbf{Points forts} & \textbf{Limites} & \textbf{Différenciation LLM Dashboard} \\
\hline
LangSmith (LangChain) & Traçage détaillé des chaînes LLM, évaluation & Pas de gouvernance, pas de facturation, couplé à LangChain & Gouvernance temps réel + facturation intégrée + indépendant du framework \\
\hline
Helicone & Proxy léger, analytics de base & Pas de détection d'anomalies ML, pas de multi-tenant complet & 3 algorithmes ML + isolation multi-tenant + RBAC 4 niveaux \\
\hline
PortKey AI & Gateway multi-fournisseur, cache sémantique & Pas de facturation contractuelle, pas de prévision budgétaire & Contrats \textit{pay-as-you-go}/forfait/hybride + prévision à 3 scénarios \\
\hline
LiteLLM & Proxy open-source, 100+ modèles & Pas d'UI, pas d'analytics, pas de détection d'anomalies & Dashboard exécutif complet + WebSocket temps réel + assistant IA \\
\hline
\end{tabularx}
\end{table}

Notre plateforme se distingue par son approche \textbf{intégrée de bout en bout}, combinant dans un seul système le proxying, la gouvernance, la détection d'anomalies par ML, l'explication IA, la prévision budgétaire, la facturation contractuelle et un tableau de bord exécutif interactif --- une combinaison qu'aucune solution existante ne propose de manière unifiée.


\section*{Conclusion}

Ce chapitre a établi un socle théorique complet pour les choix architecturaux et algorithmiques de notre plateforme. L'étude de l'architecture Transformer et des mécanismes de tokenisation a fondé la compréhension du modèle économique des LLM et la nécessité d'un monitoring granulaire. L'analyse des enjeux de gouvernance et du pattern API Gateway a justifié le positionnement d'un proxy intelligent entre les applications clientes et les fournisseurs de modèles. La présentation des séries temporelles, de TimescaleDB et des hypertables a posé les bases de notre système d'agrégation de coûts, tandis que l'étude des embeddings vectoriels et du cache sémantique a introduit les mécanismes d'optimisation exploités par le module Gateway et le module Explainer. Les sections consacrées à la sécurité API (JWT, PBKDF2), à l'architecture multi-tenant et au RBAC ont défini le cadre de protection des données et d'isolation des tenants. Les algorithmes de détection d'anomalies (STL, Isolation Forest, CUSUM) et les méthodes de prévision budgétaire fournissent les fondements mathématiques des modules d'intelligence artificielle. Enfin, l'étude de l'architecture événementielle, des tâches asynchrones Celery, de la conteneurisation Docker et de l'intégration continue CI/CD a complété le panorama des paradigmes d'ingénierie logicielle mis en \oe uvre. L'analyse comparative des solutions existantes a confirmé l'originalité de notre approche intégrée de bout en bout. Le chapitre suivant détaille le Sprint~0, consacré à l'identification des besoins et à la définition de l'environnement de travail.


% Chapitre 3 — Sprint 0 : Identification des besoins
% ============================================================
% CHAPITRE 3 — SPRINT 0 : IDENTIFICATION DES BESOINS
% ============================================================

\chapter{Sprint 0 --- Identification des Besoins}

\section*{Plan}

\begin{tabular}{r l r}
\textbf{1} & \textbf{Identification des rôles} & \textbf{\pageref{sec:roles}} \\
\textbf{2} & \textbf{Identification des besoins} & \textbf{\pageref{sec:besoins}} \\
\textbf{3} & \textbf{Product Backlog} & \textbf{\pageref{sec:backlog}} \\
\textbf{4} & \textbf{Environnement de travail} & \textbf{\pageref{sec:environnement}} \\
\end{tabular}

\section*{Introduction}

Ce chapitre constitue la phase préparatoire du projet, correspondant au Sprint~0 dans le cadre méthodologique Scrum. Nous y définissons les rôles au sein de l'équipe Scrum, identifions de manière exhaustive les besoins fonctionnels et non fonctionnels du système, présentons le diagramme de cas d'utilisation global, élaborons le Product Backlog, et décrivons l'environnement de travail retenu --- tant sur le plan logiciel que matériel.

% ────────────────────────────────────────────────────────────
\section{Identification des rôles}
\label{sec:roles}
% ────────────────────────────────────────────────────────────

L'équipe Scrum constituée pour ce projet est composée des membres suivants~:

\begin{figure}[H]
    \centering
    \fbox{\parbox{10cm}{\centering\vspace{2cm}\textit{Diagramme de l'équipe Scrum}\vspace{2cm}}}
    \caption{Composition de l'équipe Scrum}
    \label{fig:scrum_team}
\end{figure}

\begin{table}[H]
\centering
\caption{Rôles de l'équipe Scrum}
\label{tab:scrum_roles}
\begin{tabularx}{\textwidth}{|l|l|X|}
\hline
\textbf{Rôle Scrum} & \textbf{Membre} & \textbf{Responsabilités} \\
\hline
Product Owner & TIM Tunisie & Définition des priorités, validation des incréments \\
\hline
Scrum Master & Mme Salsabil Hadji & Supervision académique, levée des obstacles \\
\hline
Développeur & Houcine Laghmani & Conception, développement, tests de tous les modules \\
\hline
\end{tabularx}
\end{table}


% ────────────────────────────────────────────────────────────
\section{Identification des besoins}
\label{sec:besoins}
% ────────────────────────────────────────────────────────────

L'analyse du système a permis d'identifier \textbf{quatre acteurs principaux} et \textbf{cinquante besoins fonctionnels} organisés en douze modules, ainsi que \textbf{dix besoins non fonctionnels}.

\subsection{Identification des acteurs}

Le système LLM Dashboard interagit avec quatre catégories d'acteurs, correspondant aux rôles RBAC implémentés dans le module d'authentification~:

\begin{table}[H]
\centering
\caption{Acteurs du système LLM Dashboard}
\label{tab:acteurs}
\begin{tabularx}{\textwidth}{|l|c|X|}
\hline
\textbf{Acteur} & \textbf{Rôle RBAC} & \textbf{Description} \\
\hline
Super Administrateur & \texttt{super\_admin} & Administrateur global de la plateforme. Accès total à tous les tenants, gestion des utilisateurs, configuration système, gestion des contrats et de la facturation. \\
\hline
Administrateur Tenant & \texttt{tenant\_admin} & Responsable d'un tenant spécifique (ex~: QALITAS ou GMAO PRO). Gère les clés API, les règles de gouvernance et consulte les tableaux de bord de son périmètre. \\
\hline
Observateur & \texttt{viewer} & Utilisateur en lecture seule. Consulte les tableaux de bord, les rapports et les métriques de son tenant. \\
\hline
Agent Applicatif & \texttt{api\_key} & Application cliente qui interagit avec la plateforme via des clés API pour router ses appels LLM à travers le Gateway. \\
\hline
\end{tabularx}
\end{table}

\subsection{Besoins fonctionnels}

Les besoins fonctionnels sont organisés par module, reflétant l'architecture modulaire du système. La priorisation suit la méthode MoSCoW~: \textbf{Must} (indispensable), \textbf{Should} (important), \textbf{Could} (souhaitable).

\subsubsection{BF-1~: Authentification et gestion des accès (Module Auth)}

\begin{table}[H]
\centering
\caption{Besoins fonctionnels --- Module Authentification}
\label{tab:bf_auth}
\begin{tabularx}{\textwidth}{|c|X|c|}
\hline
\textbf{ID} & \textbf{Description} & \textbf{Priorité} \\
\hline
BF-1.1 & Le système doit permettre l'inscription et l'authentification des utilisateurs via JWT (access token + refresh token) & Must \\
\hline
BF-1.2 & Le système doit supporter un RBAC à 4 rôles avec isolation stricte par tenant & Must \\
\hline
BF-1.3 & Le système doit permettre la création, la rotation et la révocation de clés API & Must \\
\hline
BF-1.4 & Le système doit supporter les clés virtuelles avec budgets plafonds par clé & Should \\
\hline
BF-1.5 & Le système doit permettre la déconnexion avec blacklisting des tokens JWT dans Redis & Must \\
\hline
\end{tabularx}
\end{table}

\subsubsection{BF-2~: Gateway LLM et routage intelligent (Module Gateway)}

\begin{table}[H]
\centering
\caption{Besoins fonctionnels --- Module Gateway}
\label{tab:bf_gateway}
\begin{tabularx}{\textwidth}{|c|X|c|}
\hline
\textbf{ID} & \textbf{Description} & \textbf{Priorité} \\
\hline
BF-2.1 & Le système doit proxyer les appels LLM vers 5+ fournisseurs (OpenAI, Anthropic, Google, Groq, Cohere) & Must \\
\hline
BF-2.2 & Le système doit journaliser chaque appel (tokens, coût, latence, statut) dans TimescaleDB & Must \\
\hline
BF-2.3 & Le système doit évaluer les règles de gouvernance avant chaque appel & Must \\
\hline
BF-2.4 & Le système doit publier un événement \texttt{llm.call.completed} après chaque appel & Must \\
\hline
\end{tabularx}
\end{table}

\subsubsection{BF-3~: Gouvernance (Module Gateway --- Governance)}

\begin{table}[H]
\centering
\caption{Besoins fonctionnels --- Gouvernance}
\label{tab:bf_gov}
\begin{tabularx}{\textwidth}{|c|X|c|}
\hline
\textbf{ID} & \textbf{Description} & \textbf{Priorité} \\
\hline
BF-3.1 & Supporter 5 types de règles~: \texttt{rate\_limit}, \texttt{budget\_cap}, \texttt{model\_restriction}, \texttt{token\_limit}, \texttt{time\_restriction} & Must \\
\hline
BF-3.2 & Évaluer les règles en cascade et journaliser chaque décision & Must \\
\hline
BF-3.3 & Permettre la gestion CRUD des règles par tenant & Must \\
\hline
BF-3.4 & Utiliser Redis comme compteur temps réel pour les quotas RPM & Should \\
\hline
\end{tabularx}
\end{table}

\subsubsection{BF-4~: Analytics et agrégation (Module Analytics)}

\begin{table}[H]
\centering
\caption{Besoins fonctionnels --- Module Analytics}
\label{tab:bf_analytics}
\begin{tabularx}{\textwidth}{|c|X|c|}
\hline
\textbf{ID} & \textbf{Description} & \textbf{Priorité} \\
\hline
BF-4.1 & Agréger les coûts quotidiennement par tenant, modèle et provider & Must \\
\hline
BF-4.2 & Calculer les rollups mensuels pour la facturation & Must \\
\hline
BF-4.3 & Supporter un pricing versionné avec plages de dates par modèle & Should \\
\hline
BF-4.4 & Exécuter les agrégations automatiquement via Celery Beat & Must \\
\hline
\end{tabularx}
\end{table}

\subsubsection{BF-5~: Détection d'anomalies (Module Detection)}

\begin{table}[H]
\centering
\caption{Besoins fonctionnels --- Module Détection}
\label{tab:bf_detection}
\begin{tabularx}{\textwidth}{|c|X|c|}
\hline
\textbf{ID} & \textbf{Description} & \textbf{Priorité} \\
\hline
BF-5.1 & Détecter les anomalies via 3 algorithmes ML~: STL, Isolation Forest, CUSUM & Must \\
\hline
BF-5.2 & Maintenir des baselines dynamiques par tenant/modèle/agent & Must \\
\hline
BF-5.3 & Classifier les anomalies par sévérité (info, warning, critical) & Should \\
\hline
BF-5.4 & Publier un événement \texttt{anomaly.detected} pour les traitements aval & Must \\
\hline
\end{tabularx}
\end{table}

\subsubsection{BF-6~: Explication IA (Module Explainer)}

\begin{table}[H]
\centering
\caption{Besoins fonctionnels --- Module Explainer}
\label{tab:bf_explainer}
\begin{tabularx}{\textwidth}{|c|X|c|}
\hline
\textbf{ID} & \textbf{Description} & \textbf{Priorité} \\
\hline
BF-6.1 & Générer une explication en langage naturel pour chaque anomalie & Should \\
\hline
BF-6.2 & Supporter plusieurs fournisseurs LLM pour la génération & Should \\
\hline
BF-6.3 & Produire 2--4 recommandations actionnables par anomalie & Should \\
\hline
BF-6.4 & Stocker les embeddings vectoriels pour la recherche sémantique (pgvector) & Could \\
\hline
\end{tabularx}
\end{table}

\subsubsection{BF-7~: Prévision budgétaire (Module Forecasting)}

\begin{table}[H]
\centering
\caption{Besoins fonctionnels --- Module Forecasting}
\label{tab:bf_forecast}
\begin{tabularx}{\textwidth}{|c|X|c|}
\hline
\textbf{ID} & \textbf{Description} & \textbf{Priorité} \\
\hline
BF-7.1 & Prévoir les tokens et coûts futurs selon 3 scénarios (optimiste, moyen, pessimiste) & Should \\
\hline
BF-7.2 & Identifier les risques de dépassement budgétaire et alerter & Should \\
\hline
\end{tabularx}
\end{table}

\subsubsection{BF-8~: Tableau de bord (Module Dashboard)}

\begin{table}[H]
\centering
\caption{Besoins fonctionnels --- Module Dashboard}
\label{tab:bf_dashboard}
\begin{tabularx}{\textwidth}{|c|X|c|}
\hline
\textbf{ID} & \textbf{Description} & \textbf{Priorité} \\
\hline
BF-8.1 & Fournir une vue exécutive avec KPIs (tokens, coût, requêtes, taux d'erreur) & Must \\
\hline
BF-8.2 & Afficher des graphiques de tendance par période (7j, 30j, 90j) & Must \\
\hline
BF-8.3 & Supporter les mises à jour en temps réel via WebSocket & Should \\
\hline
BF-8.4 & Intégrer un assistant IA conversationnel & Could \\
\hline
\end{tabularx}
\end{table}

\subsubsection{BF-9~: Facturation (Module Billing)}

\begin{table}[H]
\centering
\caption{Besoins fonctionnels --- Module Billing}
\label{tab:bf_billing}
\begin{tabularx}{\textwidth}{|c|X|c|}
\hline
\textbf{ID} & \textbf{Description} & \textbf{Priorité} \\
\hline
BF-9.1 & Supporter 3 types de contrats~: pay-as-you-go, forfait, hybride & Must \\
\hline
BF-9.2 & Générer automatiquement des factures mensuelles détaillées par modèle & Must \\
\hline
BF-9.3 & Supporter le cycle de vie des contrats~: brouillon $\rightarrow$ proposé $\rightarrow$ actif & Must \\
\hline
BF-9.4 & Produire des factures au format PDF professionnel & Should \\
\hline
\end{tabularx}
\end{table}

\subsubsection{BF-10~: Traçage (Module Tracing)}

\begin{table}[H]
\centering
\caption{Besoins fonctionnels --- Module Tracing}
\label{tab:bf_tracing}
\begin{tabularx}{\textwidth}{|c|X|c|}
\hline
\textbf{ID} & \textbf{Description} & \textbf{Priorité} \\
\hline
BF-10.1 & Tracer les sessions de conversation LLM avec corrélation parent-enfant & Should \\
\hline
BF-10.2 & Supporter l'A/B testing de prompts & Could \\
\hline
BF-10.3 & Permettre la configuration d'alertes multi-canal & Should \\
\hline
\end{tabularx}
\end{table}

\subsubsection{BF-11~: Gestion des prompts (Module Prompts)}

\begin{table}[H]
\centering
\caption{Besoins fonctionnels --- Module Prompts}
\label{tab:bf_prompts}
\begin{tabularx}{\textwidth}{|c|X|c|}
\hline
\textbf{ID} & \textbf{Description} & \textbf{Priorité} \\
\hline
BF-11.1 & Fournir un CMS de templates de prompts avec versionnement & Should \\
\hline
BF-11.2 & Supporter les variables Jinja2 dans les templates & Should \\
\hline
BF-11.3 & Gérer le cycle de publication~: brouillon $\rightarrow$ publié $\rightarrow$ archivé & Could \\
\hline
\end{tabularx}
\end{table}


\subsection{Besoins non fonctionnels}

\begin{table}[H]
\centering
\caption{Besoins non fonctionnels}
\label{tab:bnf}
\begin{tabularx}{\textwidth}{|c|l|X|}
\hline
\textbf{ID} & \textbf{Catégorie} & \textbf{Description} \\
\hline
BNF-1 & Performance & Latence du proxy Gateway $<$~50\,ms. Dashboard $<$~2 secondes. \\
\hline
BNF-2 & Scalabilité & Ajout de tenants et de modèles LLM sans redéploiement. \\
\hline
BNF-3 & Sécurité & Chiffrement PBKDF2-SHA256, hachage BLAKE2b, TLS, protection XSS. \\
\hline
BNF-4 & Disponibilité & Health checks HTTP, reconnexion automatique, arrêt gracieux. \\
\hline
BNF-5 & Multi-tenant & Isolation stricte des données par \texttt{tenant\_id}. \\
\hline
BNF-6 & Observabilité & Métriques Prometheus, logging JSON structuré. \\
\hline
BNF-7 & Maintenabilité & Clean Architecture, principes SOLID, couverture de tests. \\
\hline
BNF-8 & Portabilité & Conteneurisation Docker complète. \\
\hline
BNF-9 & Auditabilité & Journal d'audit append-only. \\
\hline
BNF-10 & i18n & Interface français / anglais. \\
\hline
\end{tabularx}
\end{table}


\subsection{Diagramme de cas d'utilisation global}

Le diagramme de cas d'utilisation global ci-dessous illustre les interactions entre les quatre acteurs identifiés et les principaux cas d'utilisation du système.

\begin{figure}[H]
    \centering
    \fbox{\parbox{14cm}{\centering\vspace{6cm}\textit{Diagramme de cas d'utilisation global de LLM Dashboard}\\[0.5cm]\textit{(4 acteurs~: Super Admin, Tenant Admin, Viewer, Agent API)}\\[0.5cm]\textit{(Cas d'utilisation~: S'authentifier, Consulter le dashboard, Gérer la gouvernance,}\\
    \textit{Consulter les anomalies, Gérer la facturation, Envoyer un appel LLM, etc.)}\vspace{3cm}}}
    \caption{Diagramme de cas d'utilisation global de LLM Dashboard}
    \label{fig:uc_global}
\end{figure}


% ────────────────────────────────────────────────────────────
\section{Product Backlog}
\label{sec:backlog}
% ────────────────────────────────────────────────────────────

Le Product Backlog ci-dessous recense l'ensemble des User Stories du projet, organisées par Epic et affectées aux sprints correspondants~:

\begin{longtable}{|c|c|p{7.5cm}|c|}
\hline
\textbf{Epic} & \textbf{ID} & \textbf{User Story} & \textbf{Sprint} \\
\hline
\endfirsthead
\hline
\textbf{Epic} & \textbf{ID} & \textbf{User Story} & \textbf{Sprint} \\
\hline
\endhead
\hline
\endfoot

Auth & US-1 & En tant qu'utilisateur, je veux m'authentifier pour accéder à la plateforme & 1 \\
\hline
Auth & US-2 & En tant qu'admin, je veux gérer les clés API pour contrôler l'accès des applications & 1 \\
\hline
Gateway & US-3 & En tant qu'agent, je veux envoyer un appel LLM via le proxy pour bénéficier du monitoring & 1 \\
\hline
Gateway & US-4 & En tant qu'admin, je veux configurer les règles de gouvernance pour contrôler les coûts & 1 \\
\hline
Dashboard & US-5 & En tant que manager, je veux consulter le tableau de bord exécutif pour suivre les KPIs & 1 \\
\hline
Dashboard & US-6 & En tant que viewer, je veux voir les graphiques de tendance pour analyser l'évolution & 1 \\
\hline
Detection & US-7 & En tant qu'admin, je veux être alerté des anomalies de consommation pour réagir rapidement & 2 \\
\hline
Explainer & US-8 & En tant que manager, je veux comprendre les causes d'une anomalie grâce à l'IA & 2 \\
\hline
Forecasting & US-9 & En tant que décideur, je veux prévoir les coûts futurs pour planifier mon budget & 2 \\
\hline
Analytics & US-10 & En tant qu'admin, je veux consulter les coûts agrégés par modèle et par jour & 2 \\
\hline
Optimizer & US-11 & En tant qu'admin, je veux recevoir des recommandations d'optimisation des coûts & 2 \\
\hline
Billing & US-12 & En tant que super admin, je veux gérer les contrats de facturation des tenants & 3 \\
\hline
Billing & US-13 & En tant que super admin, je veux générer automatiquement les factures mensuelles & 3 \\
\hline
Tracing & US-14 & En tant qu'admin, je veux tracer les sessions de conversation LLM & 3 \\
\hline
Prompts & US-15 & En tant qu'admin, je veux gérer un catalogue de templates de prompts & 3 \\
\hline
Tenants & US-16 & En tant que super admin, je veux gérer les tenants de la plateforme & 3 \\
\hline
Deploy & US-17 & En tant que DevOps, je veux déployer la plateforme via Docker Compose & 3 \\
\hline
\end{longtable}


% ────────────────────────────────────────────────────────────
\section{Environnement de travail}
\label{sec:environnement}
% ────────────────────────────────────────────────────────────

\subsection{Architecture logicielle}

La plateforme LLM Dashboard adopte une \textbf{architecture trois-tiers} (\textit{Three-Tier Architecture})~\cite{martin2018}~:

\begin{enumerate}[noitemsep]
    \item \textbf{Couche de présentation}~: application web Next.js~16 (React~19) avec TypeScript~;
    \item \textbf{Couche de logique métier}~: API REST FastAPI (Python~3.13) avec architecture modulaire~;
    \item \textbf{Couche de données}~: PostgreSQL~15 avec TimescaleDB + Redis~7.
\end{enumerate}

\begin{table}[H]
\centering
\caption{Stack technologique complète}
\label{tab:stack}
\begin{tabularx}{\textwidth}{|p{2.8cm}|p{3.5cm}|X|}
\hline
\textbf{Couche} & \textbf{Technologie} & \textbf{Justification} \\
\hline
Frontend & Next.js~16, React~19 & SSR/CSR hybride, routing natif, écosystème React \\
\hline
State Mgmt & React Query & Cache côté client, synchronisation, retry auto \\
\hline
Visualisation & Recharts~3.8 & Graphiques React performants et déclaratifs \\
\hline
Backend API & FastAPI (Python~3.13) & Performance async, validation Pydantic, Swagger auto~\cite{fastapi2024} \\
\hline
ORM & SQLAlchemy~2.x (async) & ORM mature, support async, migrations Alembic \\
\hline
Base de données & PostgreSQL~15 & Fiabilité ACID, types riches (JSONB, UUID)\\
\hline
Séries temporelles & TimescaleDB & Hypertables, compression, agrégations rapides~\cite{timescaledb2024}  \\
\hline
Vecteurs & pgvector & Embeddings pour la recherche sémantique \\
\hline
Cache / Broker & Redis~7 & Cache mémoire, pub/sub, compteurs atomiques \\
\hline
Tâches async & Celery~5.x & Orchestration de tâches avec scheduling~\cite{celery2024} \\
\hline
Reverse Proxy & Nginx & TLS, rate limiting, WebSocket, compression \\
\hline
Conteneurisation & Docker & Isolation, reproductibilité, déploiement \\
\hline
CI/CD & GitHub Actions & Pipeline automatisé, intégration GHCR \\
\hline
\end{tabularx}
\end{table}

\subsection{Architecture matérielle (infrastructure de déploiement)}

L'infrastructure de production est entièrement conteneurisée via Docker Compose et comprend \textbf{six services}~:

\begin{table}[H]
\centering
\caption{Services Docker de la plateforme}
\label{tab:docker_services}
\begin{tabularx}{\textwidth}{|l|l|X|}
\hline
\textbf{Service} & \textbf{Image} & \textbf{Rôle} \\
\hline
\texttt{postgres} & \texttt{timescale/timescaledb:latest-pg15} & Base de données avec extensions TimescaleDB et pgvector \\
\hline
\texttt{redis} & \texttt{redis:7-alpine} & Cache, compteurs de quotas, broker Celery \\
\hline
\texttt{api} & Image custom (Dockerfile) & Serveur API FastAPI (Uvicorn, 4 workers) \\
\hline
\texttt{worker} & Image custom (Dockerfile) & Worker Celery pour les tâches asynchrones \\
\hline
\texttt{scheduler} & Image custom (Dockerfile) & Celery Beat pour les tâches périodiques \\
\hline
\texttt{nginx} & \texttt{nginx:alpine} & Reverse proxy, terminaison TLS, rate limiting \\
\hline
\end{tabularx}
\end{table}

\begin{figure}[H]
    \centering
    \fbox{\parbox{14cm}{\centering\vspace{4cm}\textit{Diagramme de déploiement Docker}\\[0.5cm]\textit{(6 services : PostgreSQL + Redis + API + Worker + Scheduler + Nginx)}\\[0.5cm]\textit{(Réseau interne : llm-net / Volumes persistants : postgres\_data, redis\_data)}\vspace{3cm}}}
    \caption{Architecture de déploiement Docker Compose}
    \label{fig:docker_arch}
\end{figure}

\subsection{Modèle de données physique}

Le schéma de base de données de la plateforme comprend \textbf{19 tables} organisées en domaines fonctionnels. Le diagramme entité-relation ci-dessous présente les entités principales et leurs relations~:

\begin{figure}[H]
    \centering
    \fbox{\parbox{14cm}{\centering\vspace{5cm}\textit{Diagramme Entité-Relation --- Modèle physique}\\[0.5cm]\textit{(19 tables : llm\_tenant, llm\_auth\_user, llm\_api\_key, llm\_token\_log,}\\
    \textit{governance\_rule, llm\_cost\_daily, llm\_cost\_monthly, llm\_model\_pricing,}\\
    \textit{llm\_anomaly, anomaly\_explanation, llm\_token\_baseline, forecast\_result,}\\
    \textit{billing\_contract, invoice, invoice\_line\_item, llm\_session, prompt\_template, ...)}\\[0.5cm]\textit{(Voir fichier PlantUML~: ch3\_er\_database.puml)}\vspace{3cm}}}
    \caption{Modèle de données physique de la plateforme LLM Dashboard}
    \label{fig:er_database}
\end{figure}



\section*{Conclusion}

Ce chapitre a permis de définir le cadre complet du Sprint~0. L'identification des rôles Scrum, la spécification exhaustive de 50 besoins fonctionnels répartis en 11 modules, la définition de 10 besoins non fonctionnels, la modélisation du cas d'utilisation global et l'élaboration du Product Backlog fournissent une base solide pour les sprints de développement. L'environnement de travail, fondé sur une stack technologique moderne et une infrastructure Docker conteneurisée, garantit la reproductibilité et la scalabilité de la solution. Le chapitre suivant détaille le Sprint~1, consacré à la mise en place de l'infrastructure technique et des modules fondamentaux.


% Chapitre 4 — Sprint 1 : Infrastructure et Gateway LLM
% ============================================================
% CHAPITRE 4 — SPRINT 1 : INFRASTRUCTURE ET GATEWAY LLM
% ============================================================

\chapter{Sprint 1 --- Infrastructure et Gateway LLM}

\section*{Plan}

\begin{tabular}{r l r}
\textbf{1} & \textbf{Sprint Planning} & \textbf{\pageref{sec:s1_planning}} \\
\textbf{2} & \textbf{Analyse des User Stories} & \textbf{\pageref{sec:s1_stories}} \\
\textbf{3} & \textbf{Conception} & \textbf{\pageref{sec:s1_conception}} \\
\textbf{4} & \textbf{Réalisation et tests} & \textbf{\pageref{sec:s1_realisation}} \\
\textbf{5} & \textbf{Sprint Rétrospective} & \textbf{\pageref{sec:s1_retro}} \\
\end{tabular}

\section*{Introduction}

Ce chapitre couvre le premier sprint de développement, consacré à la mise en place de l'infrastructure technique fondamentale de la plateforme. Les modules développés durant ce sprint constituent le socle sur lequel reposent tous les traitements ultérieurs~: l'authentification et le contrôle d'accès, le Gateway LLM avec son moteur de gouvernance, le système d'analytics et le tableau de bord exécutif.

% ────────────────────────────────────────────────────────────
\section{Sprint Planning}
\label{sec:s1_planning}
% ────────────────────────────────────────────────────────────

Le Sprint~1 couvre les User Stories suivantes, sélectionnées en fonction de leur criticité pour le fonctionnement de base de la plateforme~:

\begin{table}[H]
\centering
\caption{Sprint Backlog --- Sprint 1}
\label{tab:sprint1_backlog}
\begin{tabularx}{\textwidth}{|c|l|X|c|}
\hline
\textbf{ID} & \textbf{Module} & \textbf{User Story} & \textbf{Priorité} \\
\hline
US-1 & Auth & S'authentifier et gérer les rôles RBAC & Must \\
\hline
US-2 & Auth & Gérer les clés API et les clés virtuelles & Must \\
\hline
US-3 & Gateway & Envoyer un appel LLM via le proxy intelligent & Must \\
\hline
US-4 & Gateway & Configurer et appliquer les règles de gouvernance & Must \\
\hline
US-5 & Dashboard & Consulter le tableau de bord exécutif avec KPIs & Must \\
\hline
US-6 & Dashboard & Visualiser les graphiques de tendance & Must \\
\hline
\end{tabularx}
\end{table}


% ────────────────────────────────────────────────────────────
\section{Analyse des User Stories}
\label{sec:s1_stories}
% ────────────────────────────────────────────────────────────

\subsection{User Story~: S'authentifier (US-1)}

\begin{table}[H]
\centering
\caption{User Story US-1~: Authentification}
\label{tab:us1}
\begin{tabularx}{\textwidth}{|l|X|}
\hline
\textbf{Titre} & S'authentifier et accéder à la plateforme \\
\hline
\textbf{En tant que} & Utilisateur de la plateforme \\
\hline
\textbf{Je veux} & M'inscrire et m'authentifier avec un email et un mot de passe \\
\hline
\textbf{Afin de} & Accéder aux fonctionnalités selon mon rôle RBAC \\
\hline
\textbf{Critères d'acceptation} & \begin{minipage}[t]{\linewidth}\vspace{2pt}
\begin{itemize}[noitemsep,topsep=0pt]
\item L'authentification retourne un access token JWT (30~min) et un refresh token (7~jours)
\item Le mot de passe est haché avec PBKDF2-SHA256 (600\,000 itérations)
\item La déconnexion blackliste le token dans Redis
\item L'accès est restreint selon les 4 rôles RBAC
\end{itemize}\vspace{2pt}
\end{minipage} \\
\hline
\end{tabularx}
\end{table}


\subsection{User Story~: Gateway LLM (US-3)}

\begin{table}[H]
\centering
\caption{User Story US-3~: Gateway LLM}
\label{tab:us3}
\begin{tabularx}{\textwidth}{|l|X|}
\hline
\textbf{Titre} & Envoyer un appel LLM via le proxy intelligent \\
\hline
\textbf{En tant que} & Agent applicatif (QALITAS, GMAO PRO) \\
\hline
\textbf{Je veux} & Envoyer mes requêtes LLM à travers le Gateway \\
\hline
\textbf{Afin de} & Bénéficier du monitoring, de la gouvernance et du cache sémantique \\
\hline
\textbf{Critères d'acceptation} & \begin{minipage}[t]{\linewidth}\vspace{2pt}
\begin{itemize}[noitemsep,topsep=0pt]
\item Le Gateway route vers 5+ fournisseurs (OpenAI, Anthropic, Google, Groq, Cohere)
\item Chaque appel est journalisé dans la table \texttt{llm\_token\_log} (hypertable TimescaleDB)
\item Les règles de gouvernance sont évaluées avant l'exécution
\item Un événement \texttt{llm.call.completed} est publié sur le bus d'événements
\item Le cache sémantique réduit les appels redondants de 30--60\,\%
\end{itemize}\vspace{2pt}
\end{minipage} \\
\hline
\end{tabularx}
\end{table}


\subsection{User Story~: Gouvernance (US-4)}

\begin{table}[H]
\centering
\caption{User Story US-4~: Gouvernance temps réel}
\label{tab:us4}
\begin{tabularx}{\textwidth}{|l|X|}
\hline
\textbf{Titre} & Configurer et appliquer les règles de gouvernance \\
\hline
\textbf{En tant que} & Administrateur tenant \\
\hline
\textbf{Je veux} & Définir des règles de quota, de budget et de restriction de modèles \\
\hline
\textbf{Afin de} & Contrôler les coûts et l'utilisation des LLM de mon organisation \\
\hline
\textbf{Critères d'acceptation} & \begin{minipage}[t]{\linewidth}\vspace{2pt}
\begin{itemize}[noitemsep,topsep=0pt]
\item 5 types de règles supportés~: \texttt{rate\_limit}, \texttt{budget\_cap}, \texttt{model\_restriction}, \texttt{token\_limit}, \texttt{time\_restriction}
\item Évaluation en cascade avec journalisation de chaque décision
\item Compteurs Redis atomiques pour le rate limiting en temps réel
\item Interface CRUD complète pour la gestion des règles
\end{itemize}\vspace{2pt}
\end{minipage} \\
\hline
\end{tabularx}
\end{table}


\subsection{User Story~: Tableau de bord (US-5)}

\begin{table}[H]
\centering
\caption{User Story US-5~: Tableau de bord exécutif}
\label{tab:us5}
\begin{tabularx}{\textwidth}{|l|X|}
\hline
\textbf{Titre} & Consulter le tableau de bord exécutif \\
\hline
\textbf{En tant que} & Manager / Administrateur \\
\hline
\textbf{Je veux} & Visualiser les KPIs de consommation LLM en temps réel \\
\hline
\textbf{Afin de} & Prendre des décisions éclairées sur l'utilisation des LLM \\
\hline
\textbf{Critères d'acceptation} & \begin{minipage}[t]{\linewidth}\vspace{2pt}
\begin{itemize}[noitemsep,topsep=0pt]
\item Affichage des 4 KPIs principaux~: tokens totaux, coût total, nombre de requêtes, taux d'erreur
\item Graphiques de tendance sur 7j, 30j, 90j (Recharts)
\item Mises à jour en temps réel via WebSocket
\item Interface responsive avec animations fluides (Framer Motion)
\end{itemize}\vspace{2pt}
\end{minipage} \\
\hline
\end{tabularx}
\end{table}


\subsection{Spécification des cas d'utilisation --- Sprint 1}

Le diagramme de cas d'utilisation suivant présente les interactions spécifiques planifiées pour le Sprint~1, en mettant en évidence les rôles des acteurs et les opérations qu'ils peuvent effectuer.

\begin{figure}[H]
    \centering
    \fbox{\parbox{14cm}{\centering\vspace{3cm}\textit{Diagramme de cas d'utilisation --- Sprint 1}\\[0.5cm]\textit{(Voir fichier PlantUML~: ch4\_uc\_sprint1.puml)}\vspace{3cm}}}
    \caption{Diagramme de cas d'utilisation du Sprint 1}
    \label{fig:uc_sprint1}
\end{figure}


\subsection{Description textuelle des cas d'utilisation}

Les tableaux suivants fournissent une description détaillée de chaque cas d'utilisation développé dans ce sprint, précisant les acteurs, les préconditions, les scénarios principaux et alternatifs, conformément aux bonnes pratiques de spécification UML.

\begin{table}[H]
\centering
\caption{Description détaillée du cas d'utilisation \textit{S'authentifier}}
\label{tab:uc_desc_auth}
\begin{tabularx}{\textwidth}{|l|X|}
\hline
\textbf{Cas d'utilisation} & S'authentifier \\
\hline
\textbf{Acteur principal} & Utilisateur (Super Admin, Tenant Admin, Observateur) \\
\hline
\textbf{Préconditions} & L'utilisateur possède un compte créé dans le système \\
\hline
\textbf{Postconditions} & L'utilisateur reçoit un access token JWT et un refresh token \\
\hline
\textbf{Scénario principal} & \begin{minipage}[t]{\linewidth}\vspace{2pt}
\begin{enumerate}[noitemsep,topsep=0pt]
\item L'utilisateur accède à la page de connexion
\item L'utilisateur saisit son email et son mot de passe
\item Le système vérifie les identifiants (PBKDF2-SHA256, 600\,000 itérations)
\item Le système génère un access token JWT (validité~: 30~min) et un refresh token (7~jours)
\item Le système redirige l'utilisateur vers le tableau de bord selon son rôle RBAC
\end{enumerate}\vspace{2pt}
\end{minipage} \\
\hline
\textbf{Scénario alternatif} & \begin{minipage}[t]{\linewidth}\vspace{2pt}
\begin{itemize}[noitemsep,topsep=0pt]
\item \textbf{3a.} Identifiants invalides~: le système affiche un message d'erreur 401 et incrémente le compteur d'échecs
\item \textbf{4a.} Token expiré~: le client envoie le refresh token pour obtenir un nouveau access token
\item \textbf{5a.} Déconnexion~: le token est ajouté à la liste noire Redis avec un TTL de 30~minutes
\end{itemize}\vspace{2pt}
\end{minipage} \\
\hline
\end{tabularx}
\end{table}


\begin{table}[H]
\centering
\caption{Description détaillée du cas d'utilisation \textit{Envoyer un appel LLM}}
\label{tab:uc_desc_gateway}
\begin{tabularx}{\textwidth}{|l|X|}
\hline
\textbf{Cas d'utilisation} & Envoyer un appel LLM via le Gateway \\
\hline
\textbf{Acteur principal} & Agent applicatif (QALITAS, GMAO PRO) \\
\hline
\textbf{Préconditions} & L'agent possède une clé API valide et le tenant est actif \\
\hline
\textbf{Postconditions} & La réponse LLM est retournée et l'appel est journalisé dans TimescaleDB \\
\hline
\textbf{Scénario principal} & \begin{minipage}[t]{\linewidth}\vspace{2pt}
\begin{enumerate}[noitemsep,topsep=0pt]
\item L'agent envoie une requête POST vers \texttt{/v1/gateway/chat/completions}
\item Le Gateway évalue les 5 types de règles de gouvernance (rate limit, budget, modèle, tokens, horaire)
\item Le Gateway vérifie le cache sémantique (similarité cosinus > 0.95)
\item Cache miss~: le Gateway route la requête vers le fournisseur LLM approprié
\item Le Gateway calcule le coût (tokens $\times$ pricing par modèle)
\item Le Gateway journalise l'appel dans \texttt{llm\_token\_log} (hypertable TimescaleDB)
\item Le Gateway publie l'événement \texttt{llm.call.completed} sur le bus
\item La réponse LLM est retournée à l'agent avec les métadonnées d'usage
\end{enumerate}\vspace{2pt}
\end{minipage} \\
\hline
\textbf{Scénario alternatif} & \begin{minipage}[t]{\linewidth}\vspace{2pt}
\begin{itemize}[noitemsep,topsep=0pt]
\item \textbf{2a.} Règle bloquante~: le Gateway retourne 429 avec le motif de blocage
\item \textbf{3a.} Cache hit~: la réponse cachée est retournée (coût = 0, latence < 10~ms)
\item \textbf{4a.} Fournisseur indisponible~: le Gateway tente un fallback vers un fournisseur alternatif
\end{itemize}\vspace{2pt}
\end{minipage} \\
\hline
\end{tabularx}
\end{table}


\begin{table}[H]
\centering
\caption{Description détaillée du cas d'utilisation \textit{Configurer la gouvernance}}
\label{tab:uc_desc_governance}
\begin{tabularx}{\textwidth}{|l|X|}
\hline
\textbf{Cas d'utilisation} & Configurer les règles de gouvernance \\
\hline
\textbf{Acteur principal} & Administrateur tenant \\
\hline
\textbf{Préconditions} & L'utilisateur est authentifié avec le rôle \texttt{TENANT\_ADMIN} \\
\hline
\textbf{Postconditions} & La règle est persistée et active pour les prochains appels LLM \\
\hline
\textbf{Scénario principal} & \begin{minipage}[t]{\linewidth}\vspace{2pt}
\begin{enumerate}[noitemsep,topsep=0pt]
\item L'administrateur accède à l'interface de gouvernance
\item L'administrateur sélectionne le type de règle (rate limit, budget cap, model restriction, token limit, time restriction)
\item L'administrateur configure le seuil, la portée (tenant/utilisateur/modèle) et l'action (bloquer/alerter)
\item Le système valide les paramètres et persiste la règle
\item Le système confirme l'activation de la règle
\end{enumerate}\vspace{2pt}
\end{minipage} \\
\hline
\textbf{Scénario alternatif} & \begin{minipage}[t]{\linewidth}\vspace{2pt}
\begin{itemize}[noitemsep,topsep=0pt]
\item \textbf{4a.} Paramètres invalides~: le système affiche les erreurs de validation
\item \textbf{5a.} Règle en conflit~: le système alerte sur les conflits potentiels avec les règles existantes
\end{itemize}\vspace{2pt}
\end{minipage} \\
\hline
\end{tabularx}
\end{table}


\begin{table}[H]
\centering
\caption{Description détaillée du cas d'utilisation \textit{Consulter le tableau de bord}}
\label{tab:uc_desc_dashboard}
\begin{tabularx}{\textwidth}{|l|X|}
\hline
\textbf{Cas d'utilisation} & Consulter le tableau de bord exécutif \\
\hline
\textbf{Acteur principal} & Administrateur tenant / Observateur \\
\hline
\textbf{Préconditions} & L'utilisateur est authentifié et son tenant possède des données d'utilisation \\
\hline
\textbf{Postconditions} & Les KPIs et graphiques sont affichés avec les données les plus récentes \\
\hline
\textbf{Scénario principal} & \begin{minipage}[t]{\linewidth}\vspace{2pt}
\begin{enumerate}[noitemsep,topsep=0pt]
\item L'utilisateur accède au tableau de bord via le menu principal
\item Le système charge les 4 KPIs (tokens, coût, requêtes, taux d'erreur) via l'API REST
\item Le système établit une connexion WebSocket pour les mises à jour temps réel
\item L'utilisateur sélectionne la période d'analyse (7j, 30j, 90j)
\item Le système affiche les graphiques de tendance (Recharts) avec animations (Framer Motion)
\item Les nouvelles données sont poussées automatiquement via WebSocket
\end{enumerate}\vspace{2pt}
\end{minipage} \\
\hline
\textbf{Scénario alternatif} & \begin{minipage}[t]{\linewidth}\vspace{2pt}
\begin{itemize}[noitemsep,topsep=0pt]
\item \textbf{2a.} Aucune donnée~: le système affiche un état vide avec un message d'invitation à effectuer un premier appel LLM
\item \textbf{3a.} WebSocket indisponible~: le système bascule en mode polling (rafraîchissement toutes les 30~secondes)
\end{itemize}\vspace{2pt}
\end{minipage} \\
\hline
\end{tabularx}
\end{table}


% ────────────────────────────────────────────────────────────
\section{Conception}
\label{sec:s1_conception}
% ────────────────────────────────────────────────────────────

\subsection{Diagramme de classes --- Module Auth}

Le diagramme de classes du module Auth modélise les entités nécessaires à la gestion de l'authentification et du contrôle d'accès~:

\begin{figure}[H]
    \centering
    \fbox{\parbox{14cm}{\centering\vspace{4cm}\textit{Diagramme de classes --- Module Auth}\\[0.5cm]\textit{(Classes~: LLMAuthUser, LLMApiKey, LLMVirtualKey)}\\[0.5cm]\textit{(Relations~: User 1..* ApiKey, User 1..* VirtualKey)}\vspace{3cm}}}
    \caption{Diagramme de classes du module Auth}
    \label{fig:class_auth}
\end{figure}

\subsection{Diagramme de classes --- Module Gateway et Gouvernance}

Le diagramme de classes du module Gateway modélise les composants essentiels du proxy LLM~: le moteur de gouvernance avec ses cinq types de règles, le cache sémantique, le registre de fournisseurs et le journal d'appels~:

\begin{figure}[H]
    \centering
    \fbox{\parbox{14cm}{\centering\vspace{4cm}\textit{Diagramme de classes --- Module Gateway}\\[0.5cm]\textit{(Classes~: GatewayService, GovernanceEngine, GovernanceRule, SemanticCacheService,}\\
    \textit{ProviderRegistry, LLMProvider, LLMTokenLog, GovernanceDecision)}\\[0.5cm]\textit{(5 providers~: OpenAI, Anthropic, Google, Groq, Cohere)}\\[0.3cm]\textit{(Voir fichier PlantUML~: ch4\_class\_gateway.puml)}\vspace{3cm}}}
    \caption{Diagramme de classes du module Gateway et Gouvernance}
    \label{fig:class_gateway}
\end{figure}




\subsection{Diagramme de séquence --- Flux d'authentification}

\begin{figure}[H]
    \centering
    \fbox{\parbox{14cm}{\centering\vspace{4cm}\textit{Diagramme de séquence --- Authentification JWT}\\[0.5cm]\textit{(Client $\rightarrow$ API $\rightarrow$ AuthService $\rightarrow$ DB $\rightarrow$ Redis)}\\[0.5cm]\textit{(1. POST /v1/auth/login \quad 2. Vérification PBKDF2 \quad 3. Génération JWT \quad 4. Retour tokens)}\vspace{3cm}}}
    \caption{Diagramme de séquence --- Flux d'authentification}
    \label{fig:seq_auth}
\end{figure}

\subsection{Diagramme de séquence --- Flux Gateway LLM}

\begin{figure}[H]
    \centering
    \fbox{\parbox{14cm}{\centering\vspace{4cm}\textit{Diagramme de séquence --- Appel LLM via le Gateway}\\[0.5cm]\textit{(Agent $\rightarrow$ Gateway $\rightarrow$ GovernanceEngine $\rightarrow$ SemanticCache $\rightarrow$ LLMProvider $\rightarrow$ TokenLog)}\\[0.5cm]\textit{(1. Évaluation des règles \quad 2. Vérification cache sémantique \quad 3. Appel fournisseur \quad 4. Journalisation \quad 5. Publication événement)}\vspace{3cm}}}
    \caption{Diagramme de séquence --- Flux d'un appel LLM à travers le Gateway}
    \label{fig:seq_gateway}
\end{figure}

\subsection{Diagramme de composants --- Architecture modulaire}

\begin{figure}[H]
    \centering
    \fbox{\parbox{14cm}{\centering\vspace{4cm}\textit{Diagramme de composants --- Architecture modulaire du backend}\\[0.5cm]\textit{(12 modules~: Auth, Gateway, Analytics, Detection, Forecasting,}\\
    \textit{Dashboard, Billing, Prompts, Tracing, Observability, Tenants, Catalog)}\\[0.5cm]\textit{(Communication via Event Bus asynchrone et injection de dépendances)}\vspace{3cm}}}
    \caption{Diagramme de composants de la plateforme}
    \label{fig:component_diagram}
\end{figure}


% ────────────────────────────────────────────────────────────
\section{Réalisation et tests}
\label{sec:s1_realisation}
% ────────────────────────────────────────────────────────────

\subsection{Module Auth --- Implémentation}

Le module d'authentification implémente le standard JWT (RFC~7519)~\cite{rfc7519} avec les caractéristiques suivantes~:

\begin{lstlisting}[language=Python3, caption={Extrait du service d'authentification (modules/auth/service.py)}]
class AuthService:
    async def authenticate(self, email: str, password: str) -> TokenPair:
        user = await self.repo.get_by_email(email)
        if not user or not verify_password(password, user.hashed_password):
            raise InvalidCredentialsError()
        
        access_token = create_access_token(
            user_id=str(user.id),
            role=user.role,
            tenant_id=str(user.tenant_id),
            expires_delta=timedelta(minutes=30),
        )
        refresh_token = create_refresh_token(
            user_id=str(user.id),
            expires_delta=timedelta(days=7),
        )
        return TokenPair(access_token=access_token, refresh_token=refresh_token)
\end{lstlisting}

\subsection{Module Gateway --- Implémentation}

Le Gateway LLM constitue le c\oe ur fonctionnel de la plateforme. Il intercepte les requêtes, applique la gouvernance, interroge le cache sémantique, route vers le fournisseur approprié et journalise l'interaction~:

\begin{lstlisting}[language=Python3, caption={Flux simplifié du Gateway (modules/gateway/services/gateway\_service.py)}]
class GatewayService:
    async def process_llm_call(self, request: LLMCallRequest, 
                                tenant_id: str) -> LLMCallResponse:
        # 1. Evaluer les regles de gouvernance
        decision = await self.governance.evaluate(request, tenant_id)
        if decision.action == "block":
            raise GovernanceBlockError(decision.reason)
        
        # 2. Verifier le cache semantique
        cached = await self.semantic_cache.get(request.prompt, tenant_id)
        if cached:
            return cached
        
        # 3. Router vers le fournisseur LLM
        response = await self.provider.call(request)
        
        # 4. Journaliser dans TimescaleDB
        await self.log_repository.create(token_log_entry)
        
        # 5. Publier l'evenement
        publish_background("llm.call.completed", call_data)
        
        return response
\end{lstlisting}

\subsection{Module Dashboard --- Captures d'écran}

\begin{figure}[H]
    \centering
    \fbox{\parbox{14cm}{\centering\vspace{4cm}\textit{Capture d'écran --- Tableau de bord exécutif}\\[0.5cm]\textit{(4 cartes KPI : Tokens totaux, Coût total, Requêtes, Taux d'erreur)}\\[0.5cm]\textit{(Graphique de tendance Recharts avec sélecteur de période)}\vspace{3cm}}}
    \caption{Tableau de bord exécutif --- Vue d'ensemble}
    \label{fig:dashboard_exec}
\end{figure}

\begin{figure}[H]
    \centering
    \fbox{\parbox{14cm}{\centering\vspace{4cm}\textit{Capture d'écran --- Interface de gouvernance}\\[0.5cm]\textit{(Liste des règles actives avec type, seuil et statut)}\\[0.5cm]\textit{(Formulaire de création de règle)}\vspace{3cm}}}
    \caption{Interface de gestion de la gouvernance}
    \label{fig:governance_ui}
\end{figure}

\subsection{Tests du Sprint~1}

Les tests ont été exécutés avec \textbf{pytest} et couvrent les trois modules développés~:

\begin{table}[H]
\centering
\caption{Résultats des tests --- Sprint 1}
\label{tab:tests_s1}
\begin{tabular}{|l|c|c|c|}
\hline
\textbf{Module} & \textbf{Tests unitaires} & \textbf{Tests d'intégration} & \textbf{Résultat} \\
\hline
Auth & 12 & 5 & \textcolor{ForestGreen}{17/17 PASS} \\
\hline
Gateway & 15 & 8 & \textcolor{ForestGreen}{23/23 PASS} \\
\hline
Dashboard & 8 & 3 & \textcolor{ForestGreen}{11/11 PASS} \\
\hline
\textbf{Total Sprint 1} & \textbf{35} & \textbf{16} & \textcolor{ForestGreen}{\textbf{51/51 PASS}} \\
\hline
\end{tabular}
\end{table}


% ────────────────────────────────────────────────────────────
\section{Sprint Rétrospective}
\label{sec:s1_retro}
% ────────────────────────────────────────────────────────────

\begin{table}[H]
\centering
\caption{Rétrospective du Sprint 1}
\label{tab:retro_s1}
\begin{tabularx}{\textwidth}{|l|X|}
\hline
\textbf{Ce qui a bien fonctionné} & \begin{minipage}[t]{\linewidth}\vspace{2pt}
\begin{itemize}[noitemsep,topsep=0pt]
\item L'architecture modulaire a facilité le développement parallèle des modules
\item Le bus d'événements asynchrone a permis un découplage efficace entre les modules
\item La documentation Swagger auto-générée a accéléré le développement frontend
\end{itemize}\vspace{2pt}
\end{minipage} \\
\hline
\textbf{Ce qui peut être amélioré} & \begin{minipage}[t]{\linewidth}\vspace{2pt}
\begin{itemize}[noitemsep,topsep=0pt]
\item La configuration initiale de TimescaleDB a nécessité des ajustements
\item Le cache sémantique nécessite un fine-tuning du seuil de similarité cosinus
\end{itemize}\vspace{2pt}
\end{minipage} \\
\hline
\textbf{Actions pour le prochain sprint} & \begin{minipage}[t]{\linewidth}\vspace{2pt}
\begin{itemize}[noitemsep,topsep=0pt]
\item Implémenter les algorithmes de détection d'anomalies
\item Développer le module de prévision budgétaire
\item Mettre en place le pipeline d'agrégation Celery
\end{itemize}\vspace{2pt}
\end{minipage} \\
\hline
\end{tabularx}
\end{table}

\section*{Conclusion}

Le Sprint~1 a permis de mettre en place les fondations techniques de la plateforme LLM Dashboard. Les modules d'authentification, de Gateway LLM, de gouvernance et de tableau de bord constituent un socle fonctionnel solide. L'architecture modulaire adoptée, combinée au bus d'événements asynchrone et à l'injection de dépendances FastAPI, garantit l'extensibilité du système. Les 51 tests unitaires et d'intégration validés confirment la robustesse des implémentations. Le chapitre suivant détaille le Sprint~2, consacré aux composants d'intelligence artificielle de la plateforme.


% Chapitre 5 — Sprint 2 : Intelligence Artificielle et Analytics
% ============================================================
% CHAPITRE 5 — SPRINT 2 : INTELLIGENCE ARTIFICIELLE ET ANALYTICS
% ============================================================

\chapter{Sprint 2 --- Intelligence Artificielle et Analytics}

\section*{Plan}

\begin{tabular}{r l r}
\textbf{1} & \textbf{Sprint Planning} & \textbf{\pageref{sec:s2_planning}} \\
\textbf{2} & \textbf{Analyse des User Stories} & \textbf{\pageref{sec:s2_stories}} \\
\textbf{3} & \textbf{Conception} & \textbf{\pageref{sec:s2_conception}} \\
\textbf{4} & \textbf{Réalisation et tests} & \textbf{\pageref{sec:s2_realisation}} \\
\textbf{5} & \textbf{Sprint Rétrospective} & \textbf{\pageref{sec:s2_retro}} \\
\end{tabular}

\section*{Introduction}

Le Sprint~2 constitue le c\oe ur intelligent de la plateforme LLM Dashboard. Il est consacré au développement des composants d'apprentissage automatique et d'analyse avancée~: la détection d'anomalies de consommation par trois algorithmes complémentaires, l'explication automatique des anomalies par intelligence artificielle, la prévision budgétaire à trois scénarios, l'optimisation des coûts et le pipeline d'agrégation des données via Celery. Ce sprint transforme la plateforme d'un simple outil de monitoring en un système \textbf{prédictif et prescriptif}.

% ────────────────────────────────────────────────────────────
\section{Sprint Planning}
\label{sec:s2_planning}
% ────────────────────────────────────────────────────────────

\begin{table}[H]
\centering
\caption{Sprint Backlog --- Sprint 2}
\label{tab:sprint2_backlog}
\begin{tabularx}{\textwidth}{|c|l|X|c|}
\hline
\textbf{ID} & \textbf{Module} & \textbf{User Story} & \textbf{Priorité} \\
\hline
US-7 & Detection & Détecter les anomalies de consommation via ML & Must \\
\hline
US-8 & Explainer & Expliquer automatiquement les anomalies par IA & Should \\
\hline
US-9 & Forecasting & Prévoir les coûts futurs selon trois scénarios & Should \\
\hline
US-10 & Analytics & Agréger et analyser les coûts par modèle et par jour & Must \\
\hline
US-11 & Optimizer & Recommander des optimisations de coûts & Should \\
\hline
\end{tabularx}
\end{table}


% ────────────────────────────────────────────────────────────
\section{Analyse des User Stories}
\label{sec:s2_stories}
% ────────────────────────────────────────────────────────────

\subsection{User Story~: Détection d'anomalies (US-7)}

\begin{table}[H]
\centering
\caption{User Story US-7~: Détection d'anomalies}
\label{tab:us7}
\begin{tabularx}{\textwidth}{|l|X|}
\hline
\textbf{Titre} & Détecter automatiquement les anomalies de consommation \\
\hline
\textbf{En tant que} & Administrateur tenant \\
\hline
\textbf{Je veux} & Être alerté lorsqu'une consommation anormale de tokens est détectée \\
\hline
\textbf{Afin de} & Réagir rapidement aux dérives de coûts et aux usages non conformes \\
\hline
\textbf{Critères d'acceptation} & \begin{minipage}[t]{\linewidth}\vspace{2pt}
\begin{itemize}[noitemsep,topsep=0pt]
\item Trois algorithmes ML complémentaires~: STL decomposition~\cite{cleveland1990}, Isolation Forest~\cite{liu2008}, CUSUM~\cite{page1954}
\item Baselines dynamiques par tenant, modèle et agent
\item Classification par sévérité~: info, warning, critical
\item Publication d'un événement \texttt{anomaly.detected} sur le bus d'événements
\item Stockage des anomalies avec métadonnées complètes dans la table \texttt{llm\_anomaly}
\end{itemize}\vspace{2pt}
\end{minipage} \\
\hline
\end{tabularx}
\end{table}

\subsection{User Story~: Explication IA des anomalies (US-8)}

\begin{table}[H]
\centering
\caption{User Story US-8~: Explication IA}
\label{tab:us8}
\begin{tabularx}{\textwidth}{|l|X|}
\hline
\textbf{Titre} & Comprendre les causes d'une anomalie grâce à l'IA \\
\hline
\textbf{En tant que} & Manager / Administrateur \\
\hline
\textbf{Je veux} & Recevoir une explication en langage naturel pour chaque anomalie \\
\hline
\textbf{Afin de} & Comprendre les causes racines et agir en conséquence \\
\hline
\textbf{Critères d'acceptation} & \begin{minipage}[t]{\linewidth}\vspace{2pt}
\begin{itemize}[noitemsep,topsep=0pt]
\item Génération automatique d'une explication structurée (contexte, cause, impact)
\item Support multi-fournisseur LLM (Groq, OpenAI, Anthropic, Gemini, Ollama)
\item Production de 2--4 recommandations actionnables classées par impact
\item Stockage des embeddings vectoriels (384 dimensions) via pgvector
\item Cache des explications dans Redis pour éviter les appels LLM redondants
\end{itemize}\vspace{2pt}
\end{minipage} \\
\hline
\end{tabularx}
\end{table}

\subsection{User Story~: Prévision budgétaire (US-9)}

\begin{table}[H]
\centering
\caption{User Story US-9~: Prévision budgétaire}
\label{tab:us9}
\begin{tabularx}{\textwidth}{|l|X|}
\hline
\textbf{Titre} & Prévoir les coûts futurs pour planifier le budget \\
\hline
\textbf{En tant que} & Décideur / Manager \\
\hline
\textbf{Je veux} & Visualiser les prévisions de coûts à 30, 60 et 90 jours \\
\hline
\textbf{Afin de} & Anticiper les dépenses et éviter les dépassements budgétaires \\
\hline
\textbf{Critères d'acceptation} & \begin{minipage}[t]{\linewidth}\vspace{2pt}
\begin{itemize}[noitemsep,topsep=0pt]
\item Trois scénarios de prévision~: optimiste, moyen, pessimiste~\cite{hyndman2018}
\item Calcul de la tendance quotidienne, du taux de croissance et de la saisonnalité
\item Identification des risques de dépassement avec alertes
\item Visualisation graphique des intervalles de confiance
\end{itemize}\vspace{2pt}
\end{minipage} \\
\hline
\end{tabularx}
\end{table}

\subsection{User Story~: Analytics et agrégation (US-10)}

\begin{table}[H]
\centering
\caption{User Story US-10~: Analytics}
\label{tab:us10}
\begin{tabularx}{\textwidth}{|l|X|}
\hline
\textbf{Titre} & Consulter les coûts agrégés par modèle et par jour \\
\hline
\textbf{En tant que} & Administrateur \\
\hline
\textbf{Je veux} & Visualiser les coûts agrégés quotidiennement et mensuellement \\
\hline
\textbf{Afin de} & Comprendre la répartition des dépenses LLM \\
\hline
\textbf{Critères d'acceptation} & \begin{minipage}[t]{\linewidth}\vspace{2pt}
\begin{itemize}[noitemsep,topsep=0pt]
\item Agrégation nocturne automatique via Celery Beat~\cite{celery2024}
\item Rollup mensuel pour alimenter la facturation
\item Pricing versionné avec plages de dates (table \texttt{llm\_model\_pricing})
\item Endpoint API \texttt{/v1/analytics/costs} avec filtres par période, modèle et provider
\end{itemize}\vspace{2pt}
\end{minipage} \\
\hline
\end{tabularx}
\end{table}

\subsection{User Story~: Optimisation (US-11)}

\begin{table}[H]
\centering
\caption{User Story US-11~: Optimisation des coûts}
\label{tab:us11}
\begin{tabularx}{\textwidth}{|l|X|}
\hline
\textbf{Titre} & Recevoir des recommandations d'optimisation \\
\hline
\textbf{En tant que} & Administrateur tenant \\
\hline
\textbf{Je veux} & Obtenir des recommandations pour réduire mes coûts LLM \\
\hline
\textbf{Afin de} & Optimiser mon budget sans compromettre la qualité de service \\
\hline
\textbf{Critères d'acceptation} & \begin{minipage}[t]{\linewidth}\vspace{2pt}
\begin{itemize}[noitemsep,topsep=0pt]
\item Recommandations de substitution de modèles avec simulation ROI
\item Recommandations d'optimisation structurelle de prompts
\item Cycle de vie~: proposée $\rightarrow$ validée $\rightarrow$ déployée / rejetée
\item Scan hebdomadaire automatique via Celery Beat
\end{itemize}\vspace{2pt}
\end{minipage} \\
\hline
\end{tabularx}
\end{table}


% ────────────────────────────────────────────────────────────
\section{Conception}
\label{sec:s2_conception}
% ────────────────────────────────────────────────────────────

\subsection{Diagramme de classes --- Module Detection}

\begin{figure}[H]
    \centering
    \fbox{\parbox{14cm}{\centering\vspace{4cm}\textit{Diagramme de classes --- Module Detection}\\[0.5cm]\textit{(Classes~: LLMAnomaly, LLMTokenBaseline, AnomalyDetectionService)}\\[0.5cm]\textit{(Attributs~: anomaly\_type, severity, z\_score, algorithm\_used, details JSONB)}\vspace{3cm}}}
    \caption{Diagramme de classes du module Detection}
    \label{fig:class_detection}
\end{figure}

\subsection{Diagramme de classes --- Modules Analytics, Forecasting et Optimizer}

Le diagramme de classes suivant modélise les entités centrales des modules d'agrégation de coûts, de prévision budgétaire et d'optimisation~:

\begin{figure}[H]
    \centering
    \fbox{\parbox{14cm}{\centering\vspace{4cm}\textit{Diagramme de classes --- Modules Analytics, Forecasting et Optimizer}\\[0.5cm]\textit{(Classes~: LLMCostDaily, LLMCostMonthly, LLMModelPricing, AnalyticsService,}\\
    \textit{ForecastResult, ForecastService, OptimizationRecommendation, OptimizerService)}\\[0.5cm]\textit{(Voir fichier PlantUML~: ch5\_class\_analytics.puml)}\vspace{3cm}}}
    \caption{Diagramme de classes des modules Analytics, Forecasting et Optimizer}
    \label{fig:class_analytics}
\end{figure}

\subsection{Diagramme de séquence --- Pipeline de détection d'anomalies}

\begin{figure}[H]
    \centering
    \fbox{\parbox{14cm}{\centering\vspace{4cm}\textit{Diagramme de séquence --- Pipeline de détection d'anomalies}\\[0.5cm]\textit{(EventBus $\rightarrow$ DetectionService $\rightarrow$ STL $\rightarrow$ IsolationForest $\rightarrow$ CUSUM)}\\[0.5cm]\textit{(1. Événement reçu \quad 2. Requête baseline \quad 3. Exécution des 3 algorithmes)}\\
    \textit{(4. Classification de sévérité \quad 5. Persistance \quad 6. Publication événement)}\vspace{3cm}}}
    \caption{Diagramme de séquence du pipeline de détection d'anomalies}
    \label{fig:seq_detection}
\end{figure}

\subsection{Diagramme de séquence --- Explication IA}

\begin{figure}[H]
    \centering
    \fbox{\parbox{14cm}{\centering\vspace{4cm}\textit{Diagramme de séquence --- Explication IA d'une anomalie}\\[0.5cm]\textit{(1. Événement anomaly.detected reçu \quad 2. Construction du prompt contextuel)}\\[0.5cm]\textit{(3. Appel LLM (Groq/OpenAI) \quad 4. Parsing de la réponse structurée)}\\
    \textit{(5. Génération des embeddings \quad 6. Stockage dans pgvector)}\vspace{3cm}}}
    \caption{Diagramme de séquence de l'explication IA}
    \label{fig:seq_explainer}
\end{figure}

\subsection{Diagramme de séquence --- Prévision budgétaire}

Le diagramme suivant détaille le flux de prévision budgétaire à trois scénarios (optimiste, moyen, pessimiste), incluant l'évaluation des risques de dépassement~:

\begin{figure}[H]
    \centering
    \fbox{\parbox{14cm}{\centering\vspace{4cm}\textit{Diagramme de séquence --- Prévision budgétaire à 3 scénarios}\\[0.5cm]\textit{(Admin $\rightarrow$ Frontend $\rightarrow$ API $\rightarrow$ ForecastService $\rightarrow$ DB $\rightarrow$ RiskAssessor)}\\[0.5cm]\textit{(1. Récupération des coûts sur 90j \quad 2. Calcul de tendance et croissance)}\\
    \textit{(3. Projection 3 scénarios \quad 4. Évaluation des risques \quad 5. Persistance)}\\[0.3cm]\textit{(Voir fichier PlantUML~: ch5\_seq\_forecast.puml)}\vspace{3cm}}}
    \caption{Diagramme de séquence de la prévision budgétaire}
    \label{fig:seq_forecast}
\end{figure}

\subsection{Diagramme de séquence --- Agrégation nocturne des coûts}

Le pipeline d'agrégation, orchestré par Celery Beat, s'exécute automatiquement chaque nuit et chaque premier du mois pour produire les rollups nécessaires à la facturation~:

\begin{figure}[H]
    \centering
    \fbox{\parbox{14cm}{\centering\vspace{4cm}\textit{Diagramme de séquence --- Pipeline d'agrégation Celery}\\[0.5cm]\textit{(Celery Beat $\rightarrow$ Worker $\rightarrow$ AnalyticsService $\rightarrow$ llm\_token\_log $\rightarrow$ llm\_cost\_daily $\rightarrow$ llm\_cost\_monthly)}\\[0.5cm]\textit{(1. Agrégation quotidienne \quad 2. Déclenchement détection anomalies)}\\
    \textit{(3. Rollup mensuel \quad 4. Publication événement pour facturation)}\\[0.3cm]\textit{(Voir fichier PlantUML~: ch5\_seq\_aggregation.puml)}\vspace{3cm}}}
    \caption{Diagramme de séquence du pipeline d'agrégation nocturne}
    \label{fig:seq_aggregation}
\end{figure}

\subsection{Diagramme d'activité --- Pipeline Celery}

\begin{figure}[H]
    \centering
    \fbox{\parbox{14cm}{\centering\vspace{4cm}\textit{Diagramme d'activité --- Pipeline d'agrégation Celery}\\[0.5cm]\textit{(Tâches périodiques orchestrées par Celery Beat~:)}\\[0.3cm]\textit{- Chaque nuit~: aggregate\_daily\_costs}\\
    \textit{- Chaque 1\textsuperscript{er} du mois~: rollup\_monthly\_costs}\\
    \textit{- Chaque semaine~: run\_optimization\_scan}\vspace{3cm}}}
    \caption{Diagramme d'activité du pipeline Celery}
    \label{fig:activity_celery}
\end{figure}


% ────────────────────────────────────────────────────────────
\section{Réalisation et tests}
\label{sec:s2_realisation}
% ────────────────────────────────────────────────────────────

\subsection{Module Detection --- Implémentation}

Le module de détection d'anomalies implémente les trois algorithmes présentés dans l'état de l'art (chapitre~2). Le service principal orchestre l'exécution séquentielle des détecteurs~:

\begin{lstlisting}[language=Python3, caption={Service de détection d'anomalies (modules/detection/services/detection\_service.py)}]
class AnomalyDetectionService:
    async def detect(self, tenant_id: str, 
                      lookback_days: int = 14) -> list[LLMAnomaly]:
        # 1. Recuperer les donnees historiques
        logs = await self.repo.get_token_logs(tenant_id, lookback_days)
        baseline = await self.repo.get_baseline(tenant_id)
        
        anomalies = []
        
        # 2. STL decomposition — detecte les deviations saisonnieres
        stl_results = self._run_stl(logs, baseline)
        anomalies.extend(stl_results)
        
        # 3. Isolation Forest — detecte les points aberrants
        if_results = self._run_isolation_forest(logs)
        anomalies.extend(if_results)
        
        # 4. CUSUM — detecte les derives progressives
        cusum_results = self._run_cusum(logs, baseline)
        anomalies.extend(cusum_results)
        
        # 5. Classifier par severite et dedupliquer
        classified = self._classify_severity(anomalies)
        
        # 6. Publier les evenements
        for anomaly in classified:
            await publish("anomaly.detected", anomaly.to_dict())
        
        return classified
\end{lstlisting}


\subsection{Module Explainer --- Implémentation}

Le module Explainer utilise un LLM pour générer des explications structurées des anomalies détectées. Un prompt contextuel est construit dynamiquement avec les données de l'anomalie~:

\begin{lstlisting}[language=Python3, caption={Service d'explication IA (modules/detection/services/explainer\_service.py)}]
class ExplainerService:
    async def explain(self, anomaly: LLMAnomaly) -> Explanation:
        # Construire le prompt contextuel
        prompt = self._build_prompt(anomaly)
        
        # Appeler le LLM (fallback automatique entre providers)
        for provider in [GroqProvider, OpenAIProvider, AnthropicProvider]:
            try:
                response = await provider.generate(prompt)
                break
            except ProviderError:
                continue
        
        # Parser la reponse structuree
        explanation = self._parse_explanation(response)
        
        # Generer l'embedding pour la recherche semantique
        embedding = self.embedder.encode(explanation.summary)
        
        # Stocker dans pgvector
        await self.repo.save(explanation, embedding)
        
        return explanation
\end{lstlisting}


\subsection{Module Forecasting --- Implémentation}

Le module de prévision budgétaire calcule trois scénarios (optimiste, moyen, pessimiste) fondés sur l'analyse de la tendance et de la saisonnalité des données historiques~\cite{hyndman2018}~:

\begin{lstlisting}[language=Python3, caption={Service de prévision budgétaire (modules/forecasting/services/forecast\_service.py)}]
class ForecastService:
    async def forecast(self, tenant_id: str, 
                        horizon_days: int = 30) -> ForecastResult:
        # Recuperer les couts quotidiens historiques
        daily_costs = await self.repo.get_daily_costs(tenant_id, days=90)
        
        # Calculer la tendance et le taux de croissance
        trend = self._compute_trend(daily_costs)
        growth_rate = self._compute_growth_rate(daily_costs)
        
        # Generer les 3 scenarios
        optimistic = self._project(trend, growth_rate * 0.5, horizon_days)
        baseline   = self._project(trend, growth_rate,       horizon_days)
        pessimistic = self._project(trend, growth_rate * 1.5, horizon_days)
        
        # Evaluer les risques de depassement
        risks = self._assess_budget_risks(pessimistic, tenant_id)
        
        return ForecastResult(
            optimistic=optimistic,
            baseline=baseline,
            pessimistic=pessimistic,
            risks=risks,
        )
\end{lstlisting}


\subsection{Pipeline Celery --- Tâches périodiques}

Les tâches d'agrégation et d'analyse sont orchestrées par \textbf{Celery Beat}~\cite{celery2024}, qui déclenche automatiquement les traitements selon un calendrier configuré~:

\begin{lstlisting}[language=Python3, caption={Configuration de Celery Beat (celery\_app.py)}]
app.conf.beat_schedule = {
    "aggregate-daily-costs": {
        "task": "modules.analytics.tasks.aggregate_daily_costs",
        "schedule": crontab(hour=2, minute=0),   # Chaque nuit a 2h
    },
    "rollup-monthly-costs": {
        "task": "modules.analytics.tasks.rollup_monthly_costs",
        "schedule": crontab(day_of_month=1, hour=3),  # 1er du mois
    },
    "run-optimization-scan": {
        "task": "modules.optimizer.tasks.run_weekly_scan",
        "schedule": crontab(day_of_week=1, hour=4),  # Chaque lundi
    },
}
\end{lstlisting}


\subsection{Captures d'écran}

\begin{figure}[H]
    \centering
    \fbox{\parbox{14cm}{\centering\vspace{4cm}\textit{Capture d'écran --- Page des anomalies détectées}\\[0.5cm]\textit{(Liste des anomalies avec sévérité, algorithme, date et tenant)}\\[0.5cm]\textit{(Détail d'une anomalie avec explication IA et recommandations)}\vspace{3cm}}}
    \caption{Interface de visualisation des anomalies détectées}
    \label{fig:anomalies_ui}
\end{figure}

\begin{figure}[H]
    \centering
    \fbox{\parbox{14cm}{\centering\vspace{4cm}\textit{Capture d'écran --- Page de prévision budgétaire}\\[0.5cm]\textit{(Graphique avec 3 scénarios~: optimiste (vert), moyen (bleu), pessimiste (rouge))}\\[0.5cm]\textit{(Tableau des risques de dépassement par tenant)}\vspace{3cm}}}
    \caption{Interface de prévision budgétaire à trois scénarios}
    \label{fig:forecasting_ui}
\end{figure}

\begin{figure}[H]
    \centering
    \fbox{\parbox{14cm}{\centering\vspace{4cm}\textit{Capture d'écran --- Page des coûts et analytics}\\[0.5cm]\textit{(Répartition des coûts par modèle, par provider, par jour)}\\[0.5cm]\textit{(Graphiques en barres et en aires empilées)}\vspace{3cm}}}
    \caption{Interface d'analyse des coûts}
    \label{fig:costs_ui}
\end{figure}


\subsection{Tests du Sprint~2}

\begin{table}[H]
\centering
\caption{Résultats des tests --- Sprint 2}
\label{tab:tests_s2}
\begin{tabular}{|l|c|c|c|}
\hline
\textbf{Module} & \textbf{Tests unitaires} & \textbf{Tests d'intégration} & \textbf{Résultat} \\
\hline
Detection & 18 & 6 & \textcolor{ForestGreen}{24/24 PASS} \\
\hline
Explainer & 10 & 4 & \textcolor{ForestGreen}{14/14 PASS} \\
\hline
Forecasting & 8 & 3 & \textcolor{ForestGreen}{11/11 PASS} \\
\hline
Analytics & 12 & 5 & \textcolor{ForestGreen}{17/17 PASS} \\
\hline
Optimizer & 7 & 3 & \textcolor{ForestGreen}{10/10 PASS} \\
\hline
\textbf{Total Sprint 2} & \textbf{55} & \textbf{21} & \textcolor{ForestGreen}{\textbf{76/76 PASS}} \\
\hline
\end{tabular}
\end{table}


% ────────────────────────────────────────────────────────────
\section{Sprint Rétrospective}
\label{sec:s2_retro}
% ────────────────────────────────────────────────────────────

\begin{table}[H]
\centering
\caption{Rétrospective du Sprint 2}
\label{tab:retro_s2}
\begin{tabularx}{\textwidth}{|l|X|}
\hline
\textbf{Ce qui a bien fonctionné} & \begin{minipage}[t]{\linewidth}\vspace{2pt}
\begin{itemize}[noitemsep,topsep=0pt]
\item La complémentarité des trois algorithmes de détection assure une couverture complète
\item Le bus d'événements permet un chaînage automatique~: détection $\rightarrow$ explication $\rightarrow$ notification
\item Le pipeline Celery gère efficacement les tâches d'agrégation nocturnes
\item L'intégration de pgvector permet la recherche sémantique d'anomalies similaires
\end{itemize}\vspace{2pt}
\end{minipage} \\
\hline
\textbf{Ce qui peut être amélioré} & \begin{minipage}[t]{\linewidth}\vspace{2pt}
\begin{itemize}[noitemsep,topsep=0pt]
\item Le seuil de détection de l'Isolation Forest nécessite un calibrage par tenant
\item Les appels LLM pour l'explication peuvent être coûteux~; le cache sémantique atténue ce problème
\end{itemize}\vspace{2pt}
\end{minipage} \\
\hline
\textbf{Actions pour le prochain sprint} & \begin{minipage}[t]{\linewidth}\vspace{2pt}
\begin{itemize}[noitemsep,topsep=0pt]
\item Développer le module de facturation contractuelle
\item Implémenter le traçage des sessions et la gestion des prompts
\item Mettre en place le déploiement Docker et le pipeline CI/CD
\end{itemize}\vspace{2pt}
\end{minipage} \\
\hline
\end{tabularx}
\end{table}


\section*{Conclusion}

Le Sprint~2 a doté la plateforme LLM Dashboard de capacités d'intelligence artificielle avancées. Le pipeline de détection d'anomalies, combinant STL, Isolation Forest et CUSUM, identifie de manière proactive les comportements de consommation anormaux. Le module Explainer produit des explications intelligibles et des recommandations actionnables grâce aux LLM. Le moteur de prévision budgétaire offre une visibilité prospective à trois scénarios, tandis que le pipeline Celery automatise l'ensemble des traitements d'agrégation. Les 76 tests validés confirment la fiabilité des implémentations. Le chapitre suivant détaille le Sprint~3, consacré aux fonctionnalités entreprise et au déploiement de la plateforme.


% Chapitre 6 — Sprint 3 : Fonctionnalités Entreprise et Déploiement
% ============================================================
% CHAPITRE 6 — SPRINT 3 : FONCTIONNALITÉS ENTREPRISE ET DÉPLOIEMENT
% ============================================================

\chapter{Sprint 3 --- Fonctionnalités Entreprise et Déploiement}

\section*{Plan}

\begin{tabular}{r l r}
\textbf{1} & \textbf{Sprint Planning} & \textbf{\pageref{sec:s3_planning}} \\
\textbf{2} & \textbf{Analyse des User Stories} & \textbf{\pageref{sec:s3_stories}} \\
\textbf{3} & \textbf{Conception} & \textbf{\pageref{sec:s3_conception}} \\
\textbf{4} & \textbf{Fonctionnalités du système} & \textbf{\pageref{sec:s3_features}} \\
\textbf{5} & \textbf{Implémentation et tests} & \textbf{\pageref{sec:s3_realisation}} \\
\textbf{6} & \textbf{Sprint Rétrospective} & \textbf{\pageref{sec:s3_retro}} \\
\end{tabular}

\section*{Introduction}

Le Sprint~3 parachève le développement de la plateforme LLM Dashboard en ajoutant les fonctionnalités entreprise indispensables à son exploitation commerciale~: la facturation contractuelle avec génération automatique de factures, le traçage des sessions conversationnelles, la gestion des templates de prompts, l'administration des tenants et le déploiement conteneurisé avec intégration continue. Ce sprint transforme le prototype technique en un \textbf{produit SaaS opérationnel} prêt pour la mise en production.

% ────────────────────────────────────────────────────────────
\section{Sprint Planning}
\label{sec:s3_planning}
% ────────────────────────────────────────────────────────────

\begin{table}[H]
\centering
\caption{Sprint Backlog --- Sprint 3}
\label{tab:sprint3_backlog}
\begin{tabularx}{\textwidth}{|c|l|X|c|}
\hline
\textbf{ID} & \textbf{Module} & \textbf{User Story} & \textbf{Priorité} \\
\hline
US-12 & Billing & Gérer les contrats de facturation des tenants & Must \\
\hline
US-13 & Billing & Générer automatiquement les factures mensuelles & Must \\
\hline
US-14 & Tracing & Tracer les sessions de conversation LLM & Should \\
\hline
US-15 & Prompts & Gérer un catalogue de templates de prompts & Should \\
\hline
US-16 & Tenants & Gérer les tenants de la plateforme & Must \\
\hline
US-17 & Deploy & Déployer la plateforme via Docker et CI/CD & Must \\
\hline
\end{tabularx}
\end{table}


% ────────────────────────────────────────────────────────────
\section{Analyse des User Stories}
\label{sec:s3_stories}
% ────────────────────────────────────────────────────────────

\subsection{User Story~: Facturation contractuelle (US-12)}

\begin{table}[H]
\centering
\caption{User Story US-12~: Gestion des contrats}
\label{tab:us12}
\begin{tabularx}{\textwidth}{|l|X|}
\hline
\textbf{Titre} & Gérer les contrats de facturation des tenants \\
\hline
\textbf{En tant que} & Super Administrateur \\
\hline
\textbf{Je veux} & Créer et gérer les contrats de facturation pour chaque tenant \\
\hline
\textbf{Afin de} & Définir les conditions tarifaires adaptées à chaque client \\
\hline
\textbf{Critères d'acceptation} & \begin{minipage}[t]{\linewidth}\vspace{2pt}
\begin{itemize}[noitemsep,topsep=0pt]
\item Trois types de contrats~: \textbf{pay-as-you-go} (facturation à l'usage), \textbf{forfait} (enveloppe mensuelle fixe), \textbf{hybride} (forfait + dépassement)
\item Cycle de vie des contrats~: brouillon $\rightarrow$ proposé $\rightarrow$ accepté $\rightarrow$ actif
\item Paramètres configurables~: frais de base, seuil de tokens inclus, taux de dépassement, remise
\item Interface CRUD complète avec recherche et filtrage
\end{itemize}\vspace{2pt}
\end{minipage} \\
\hline
\end{tabularx}
\end{table}

\subsection{User Story~: Génération de factures (US-13)}

\begin{table}[H]
\centering
\caption{User Story US-13~: Génération de factures}
\label{tab:us13}
\begin{tabularx}{\textwidth}{|l|X|}
\hline
\textbf{Titre} & Générer automatiquement les factures mensuelles \\
\hline
\textbf{En tant que} & Super Administrateur \\
\hline
\textbf{Je veux} & Que le système génère automatiquement les factures en fin de mois \\
\hline
\textbf{Afin de} & Automatiser le processus de facturation et réduire les erreurs \\
\hline
\textbf{Critères d'acceptation} & \begin{minipage}[t]{\linewidth}\vspace{2pt}
\begin{itemize}[noitemsep,topsep=0pt]
\item Calcul automatique basé sur les agrégations mensuelles (table \texttt{llm\_cost\_monthly})
\item Lignes de facture détaillées par modèle LLM avec tokens et coûts unitaires
\item Support du cycle de vie~: brouillon $\rightarrow$ finalisée $\rightarrow$ envoyée $\rightarrow$ payée
\item Application automatique des conditions contractuelles (forfait, dépassement, remise)
\end{itemize}\vspace{2pt}
\end{minipage} \\
\hline
\end{tabularx}
\end{table}

\subsection{User Story~: Traçage des sessions (US-14)}

\begin{table}[H]
\centering
\caption{User Story US-14~: Traçage des sessions}
\label{tab:us14}
\begin{tabularx}{\textwidth}{|l|X|}
\hline
\textbf{Titre} & Tracer les sessions de conversation LLM \\
\hline
\textbf{En tant que} & Administrateur \\
\hline
\textbf{Je veux} & Suivre l'historique des sessions de conversation avec les LLM \\
\hline
\textbf{Afin de} & Analyser les patterns d'usage et déboguer les interactions \\
\hline
\textbf{Critères d'acceptation} & \begin{minipage}[t]{\linewidth}\vspace{2pt}
\begin{itemize}[noitemsep,topsep=0pt]
\item Corrélation parent-enfant entre les appels d'une même conversation
\item Visualisation en cascade (\textit{waterfall}) des latences
\item Suivi du nombre de tours (\textit{turns}), tokens totaux et coût par session
\item Support de l'A/B testing de prompts avec allocation de trafic
\end{itemize}\vspace{2pt}
\end{minipage} \\
\hline
\end{tabularx}
\end{table}

\subsection{User Story~: Gestion des prompts (US-15)}

\begin{table}[H]
\centering
\caption{User Story US-15~: Gestion des prompts}
\label{tab:us15}
\begin{tabularx}{\textwidth}{|l|X|}
\hline
\textbf{Titre} & Gérer un catalogue de templates de prompts \\
\hline
\textbf{En tant que} & Administrateur \\
\hline
\textbf{Je veux} & Créer, versionner et publier des templates de prompts réutilisables \\
\hline
\textbf{Afin de} & Standardiser et optimiser les prompts utilisés par l'organisation \\
\hline
\textbf{Critères d'acceptation} & \begin{minipage}[t]{\linewidth}\vspace{2pt}
\begin{itemize}[noitemsep,topsep=0pt]
\item CMS de templates avec versionnement (v1, v2, v3\ldots)
\item Support des variables Jinja2 avec rendu dynamique
\item Cycle de publication~: brouillon $\rightarrow$ publié $\rightarrow$ archivé
\item Catégorisation par tags et recherche textuelle
\end{itemize}\vspace{2pt}
\end{minipage} \\
\hline
\end{tabularx}
\end{table}


% ────────────────────────────────────────────────────────────
\section{Conception}
\label{sec:s3_conception}
% ────────────────────────────────────────────────────────────

\subsection{Diagramme de classes --- Module Billing}

Le diagramme de classes du module Billing modélise le système de facturation contractuelle avec ses trois types de contrats et le cycle de vie des factures~:

\begin{figure}[H]
    \centering
    \fbox{\parbox{14cm}{\centering\vspace{4cm}\textit{Diagramme de classes --- Module Billing}\\[0.5cm]\textit{(Classes~: BillingContract, Invoice, InvoiceLineItem, BillingService)}\\[0.5cm]\textit{(Enums~: ContractType, ContractStatus, InvoiceStatus)}\\[0.3cm]\textit{(Voir fichier PlantUML~: ch6\_class\_billing.puml)}\vspace{3cm}}}
    \caption{Diagramme de classes du module Billing}
    \label{fig:class_billing}
\end{figure}

\subsection{Diagramme de classes --- Modules Tracing et Prompts}

Le diagramme suivant modélise les entités nécessaires au traçage des sessions conversationnelles, à la gestion des templates de prompts et aux alertes multi-canal~:

\begin{figure}[H]
    \centering
    \fbox{\parbox{14cm}{\centering\vspace{4cm}\textit{Diagramme de classes --- Modules Tracing et Prompts}\\[0.5cm]\textit{(Classes~: LLMSession, LLMTokenLog, TracingService, ABTestConfig,}\\
    \textit{PromptTemplate, PromptService, AlertRule)}\\[0.5cm]\textit{(Enums~: SessionStatus, PromptStatus, AlertChannel)}\\[0.3cm]\textit{(Voir fichier PlantUML~: ch6\_class\_tracing.puml)}\vspace{3cm}}}
    \caption{Diagramme de classes des modules Tracing et Prompts}
    \label{fig:class_tracing}
\end{figure}

\subsection{Diagramme de séquence --- Traçage d'une session LLM}

Le diagramme suivant illustre le cycle de vie complet d'une session conversationnelle~: création automatique, ajout de tours successifs, fermeture et consultation de la vue en cascade~:

\begin{figure}[H]
    \centering
    \fbox{\parbox{14cm}{\centering\vspace{4cm}\textit{Diagramme de séquence --- Traçage d'une session}\\[0.5cm]\textit{(Agent $\rightarrow$ Gateway $\rightarrow$ TracingService $\rightarrow$ DB $\rightarrow$ EventBus $\rightarrow$ Dashboard)}\\[0.5cm]\textit{(1. Création session \quad 2. Ajout tours successifs \quad 3. Fermeture)}\\
    \textit{(4. Consultation waterfall avec latences et tokens par tour)}\\[0.3cm]\textit{(Voir fichier PlantUML~: ch6\_seq\_tracing.puml)}\vspace{3cm}}}
    \caption{Diagramme de séquence du traçage d'une session LLM}
    \label{fig:seq_tracing}
\end{figure}

\subsection{Diagramme de séquence --- Gestion des prompts}

Le flux de gestion des templates de prompts comprend la création, la publication, le versionnement et le rendu dynamique avec des variables Jinja2~:

\begin{figure}[H]
    \centering
    \fbox{\parbox{14cm}{\centering\vspace{4cm}\textit{Diagramme de séquence --- Gestion des templates de prompts}\\[0.5cm]\textit{(Admin $\rightarrow$ Frontend $\rightarrow$ API $\rightarrow$ PromptService $\rightarrow$ DB $\rightarrow$ Jinja2Engine)}\\[0.5cm]\textit{(1. Création DRAFT \quad 2. Publication \quad 3. Nouvelle version)}\\
    \textit{(4. Rendu dynamique avec variables)}\\[0.3cm]\textit{(Voir fichier PlantUML~: ch6\_seq\_prompt\_mgmt.puml)}\vspace{3cm}}}
    \caption{Diagramme de séquence de la gestion des prompts}
    \label{fig:seq_prompts}
\end{figure}

\subsection{Diagramme de séquence --- Facturation}

\begin{figure}[H]
    \centering
    \fbox{\parbox{14cm}{\centering\vspace{4cm}\textit{Diagramme de séquence --- Génération automatique de factures}\\[0.5cm]\textit{(Celery Beat $\rightarrow$ BillingService $\rightarrow$ ContractRepo $\rightarrow$ CostMonthlyRepo)}\\[0.5cm]\textit{(1. Récupération contrat actif \quad 2. Calcul selon type)}\\
    \textit{(3. Application remise \quad 4. Création facture avec lignes détaillées)}\\[0.3cm]\textit{(Voir fichier PlantUML~: ch6\_seq\_billing.puml)}\vspace{3cm}}}
    \caption{Diagramme de séquence de la génération de factures}
    \label{fig:seq_billing}
\end{figure}

\subsection{Diagrammes d'états --- Cycles de vie}

Les diagrammes suivants modélisent les machines à états des trois entités dont le cycle de vie est régi par des transitions formalisées~:

\begin{figure}[H]
    \centering
    \fbox{\parbox{14cm}{\centering\vspace{3cm}\textit{Diagramme d'états --- Cycle de vie d'un contrat}\\[0.5cm]\textit{(DRAFT $\rightarrow$ PROPOSED $\rightarrow$ ACCEPTED $\rightarrow$ ACTIVE $\rightarrow$ TERMINATED)}\\[0.3cm]\textit{(Voir fichier PlantUML~: ch6\_state\_contract.puml)}\vspace{2cm}}}
    \caption{Machine à états du cycle de vie d'un contrat}
    \label{fig:state_contract}
\end{figure}

\begin{figure}[H]
    \centering
    \fbox{\parbox{14cm}{\centering\vspace{3cm}\textit{Diagramme d'états --- Cycle de vie d'une facture}\\[0.5cm]\textit{(DRAFT $\rightarrow$ FINALIZED $\rightarrow$ SENT $\rightarrow$ PAID / OVERDUE)}\\[0.3cm]\textit{(Voir fichier PlantUML~: ch6\_state\_invoice.puml)}\vspace{2cm}}}
    \caption{Machine à états du cycle de vie d'une facture}
    \label{fig:state_invoice}
\end{figure}

\begin{figure}[H]
    \centering
    \fbox{\parbox{14cm}{\centering\vspace{3cm}\textit{Diagramme d'états --- Cycle de vie d'un template de prompt}\\[0.5cm]\textit{(DRAFT $\rightarrow$ PUBLISHED $\rightarrow$ ARCHIVED / NEW\_VERSION)}\\[0.3cm]\textit{(Voir fichier PlantUML~: ch6\_state\_prompt.puml)}\vspace{2cm}}}
    \caption{Machine à états du cycle de vie d'un template de prompt}
    \label{fig:state_prompt}
\end{figure}


% ────────────────────────────────────────────────────────────
\section{Fonctionnalités du système}
\label{sec:s3_features}
% ────────────────────────────────────────────────────────────

Cette section présente les principales interfaces de la plateforme LLM Dashboard, offrant une vue d'ensemble des fonctionnalités accessibles aux différents utilisateurs.

\subsection{Page d'accueil --- Tableau de bord exécutif}

\begin{figure}[H]
    \centering
    \fbox{\parbox{14cm}{\centering\vspace{4cm}\textit{Capture d'écran --- Tableau de bord exécutif}\\[0.5cm]\textit{(KPIs : Tokens totaux, Coût total, Requêtes, Taux d'erreur)}\\[0.5cm]\textit{(Graphiques de tendance, répartition par modèle, top tenants)}\vspace{3cm}}}
    \caption{Page d'accueil --- Tableau de bord exécutif}
    \label{fig:feat_dashboard}
\end{figure}

\subsection{Page Anomalies --- Détection et explication IA}

\begin{figure}[H]
    \centering
    \fbox{\parbox{14cm}{\centering\vspace{4cm}\textit{Capture d'écran --- Page des anomalies}\\[0.5cm]\textit{(Liste filtrable par sévérité, algorithme, date)}\\[0.5cm]\textit{(Détail avec explication IA, recommandations et anomalies similaires)}\vspace{3cm}}}
    \caption{Page Anomalies --- Détection et explication IA}
    \label{fig:feat_anomalies}
\end{figure}

\subsection{Page Gouvernance --- Configuration des règles}

\begin{figure}[H]
    \centering
    \fbox{\parbox{14cm}{\centering\vspace{4cm}\textit{Capture d'écran --- Page de gouvernance}\\[0.5cm]\textit{(Tableau des règles actives avec type, seuil et scope)}\\[0.5cm]\textit{(Formulaire de création : rate\_limit, budget\_cap, model\_restriction, etc.)}\vspace{3cm}}}
    \caption{Page Gouvernance --- Configuration des règles}
    \label{fig:feat_governance}
\end{figure}

\subsection{Page Facturation --- Contrats et factures}

\begin{figure}[H]
    \centering
    \fbox{\parbox{14cm}{\centering\vspace{4cm}\textit{Capture d'écran --- Page de facturation}\\[0.5cm]\textit{(Onglets : Factures / Paramètres contrat / Tarification)}\\[0.5cm]\textit{(Détail d'une facture avec lignes par modèle et montants)}\vspace{3cm}}}
    \caption{Page Facturation --- Contrats et factures}
    \label{fig:feat_billing}
\end{figure}

\subsection{Page Playground --- Test d'appels LLM}

\begin{figure}[H]
    \centering
    \fbox{\parbox{14cm}{\centering\vspace{4cm}\textit{Capture d'écran --- Playground LLM}\\[0.5cm]\textit{(Éditeur Monaco pour le prompt, sélecteur de modèle/provider)}\\[0.5cm]\textit{(Résultat avec tokens consommés, coût et latence)}\vspace{3cm}}}
    \caption{Page Playground --- Interface de test d'appels LLM}
    \label{fig:feat_playground}
\end{figure}

\subsection{Page Prévisions --- Forecasting budgétaire}

\begin{figure}[H]
    \centering
    \fbox{\parbox{14cm}{\centering\vspace{4cm}\textit{Capture d'écran --- Page de prévisions budgétaires}\\[0.5cm]\textit{(3 courbes~: optimiste, moyen, pessimiste)}\\[0.5cm]\textit{(Alertes de risque et horizon configurable)}\vspace{3cm}}}
    \caption{Page Prévisions --- Forecasting budgétaire}
    \label{fig:feat_forecasting}
\end{figure}

\subsection{Page Sessions --- Traçage des conversations}

\begin{figure}[H]
    \centering
    \fbox{\parbox{14cm}{\centering\vspace{4cm}\textit{Capture d'écran --- Page de traçage des sessions}\\[0.5cm]\textit{(Liste des sessions avec durée, tours, tokens totaux)}\\[0.5cm]\textit{(Vue waterfall des appels individuels)}\vspace{3cm}}}
    \caption{Page Sessions --- Traçage des conversations LLM}
    \label{fig:feat_sessions}
\end{figure}

\subsection{Page Tenants --- Administration multi-tenant}

\begin{figure}[H]
    \centering
    \fbox{\parbox{14cm}{\centering\vspace{4cm}\textit{Capture d'écran --- Page de gestion des tenants}\\[0.5cm]\textit{(Tableau des tenants avec statut, nombre d'utilisateurs, consommation)}\\[0.5cm]\textit{(Actions : créer, activer, désactiver, configurer le contrat)}\vspace{3cm}}}
    \caption{Page Tenants --- Administration multi-tenant}
    \label{fig:feat_tenants}
\end{figure}


% ────────────────────────────────────────────────────────────
\section{Implémentation et tests}
\label{sec:s3_realisation}
% ────────────────────────────────────────────────────────────

\subsection{Module Billing --- Implémentation}

Le module de facturation implémente un système complet de gestion des contrats et de génération de factures~:

\begin{lstlisting}[language=Python3, caption={Service de génération de factures (modules/billing/services/billing\_service.py)}]
class BillingService:
    async def generate_monthly_invoice(self, tenant_id: str,
                                        month: int, year: int) -> Invoice:
        # 1. Recuperer le contrat actif du tenant
        contract = await self.contract_repo.get_active(tenant_id)
        
        # 2. Recuperer les couts mensuels agreges
        monthly_costs = await self.cost_repo.get_monthly(
            tenant_id, month, year
        )
        
        # 3. Calculer selon le type de contrat
        if contract.contract_type == "pay_as_you_go":
            total = sum(c.total_cost for c in monthly_costs)
        elif contract.contract_type == "flat_rate":
            total = contract.base_fee
        else:  # hybrid
            usage = sum(c.total_cost for c in monthly_costs)
            if usage > contract.included_amount:
                overage = (usage - contract.included_amount) 
                          * contract.overage_rate
                total = contract.base_fee + overage
            else:
                total = contract.base_fee
        
        # 4. Appliquer la remise
        total *= (1 - contract.discount_pct / 100)
        
        # 5. Creer la facture avec les lignes detaillees
        invoice = await self.invoice_repo.create(
            tenant_id=tenant_id, total=total,
            line_items=self._build_line_items(monthly_costs),
        )
        return invoice
\end{lstlisting}


\subsection{Déploiement Docker et CI/CD}

Le déploiement de la plateforme repose sur \textbf{Docker Compose} pour l'orchestration des six services conteneurisés et sur \textbf{GitHub Actions} pour l'intégration continue~:

\begin{lstlisting}[language=bash, caption={Déploiement Docker Compose (extrait de docker-compose.yml)}]
services:
  postgres:
    image: timescale/timescaledb:latest-pg15
    healthcheck:
      test: ["CMD-SHELL", "pg_isready -U llmuser"]
  
  redis:
    image: redis:7-alpine
    command: redis-server --maxmemory 256mb
  
  api:
    build: .
    depends_on:
      postgres: { condition: service_healthy }
      redis:    { condition: service_healthy }
  
  worker:
    build: .
    command: celery -A celery_app worker --concurrency=2
  
  scheduler:
    build: .
    command: celery -A celery_app beat
  
  nginx:
    image: nginx:alpine
    ports: ["80:80", "443:443"]
\end{lstlisting}

Le pipeline CI/CD comprend quatre étapes exécutées automatiquement à chaque push~:

\begin{enumerate}[noitemsep]
    \item \textbf{Lint}~: vérification syntaxique de tous les fichiers Python~;
    \item \textbf{Test}~: exécution de pytest avec services TimescaleDB et Redis éphémères~;
    \item \textbf{Build}~: construction de l'image Docker multi-stage et push vers GHCR~;
    \item \textbf{Deploy}~: notification et déploiement sur le serveur de production.
\end{enumerate}


\subsection{Tests du Sprint~3}

\begin{table}[H]
\centering
\caption{Résultats des tests --- Sprint 3}
\label{tab:tests_s3}
\begin{tabular}{|l|c|c|c|}
\hline
\textbf{Module} & \textbf{Tests unitaires} & \textbf{Tests d'intégration} & \textbf{Résultat} \\
\hline
Billing & 14 & 6 & \textcolor{ForestGreen}{20/20 PASS} \\
\hline
Tracing & 8 & 4 & \textcolor{ForestGreen}{12/12 PASS} \\
\hline
Prompts & 7 & 3 & \textcolor{ForestGreen}{10/10 PASS} \\
\hline
Tenants & 5 & 2 & \textcolor{ForestGreen}{7/7 PASS} \\
\hline
Déploiement & -- & 4 & \textcolor{ForestGreen}{4/4 PASS} \\
\hline
\textbf{Total Sprint 3} & \textbf{34} & \textbf{19} & \textcolor{ForestGreen}{\textbf{53/53 PASS}} \\
\hline
\end{tabular}
\end{table}

\subsection{Synthèse globale des tests}

\begin{table}[H]
\centering
\caption{Synthèse globale des tests de la plateforme}
\label{tab:tests_global}
\begin{tabular}{|l|c|c|c|}
\hline
\textbf{Sprint} & \textbf{Tests unitaires} & \textbf{Tests d'intégration} & \textbf{Total} \\
\hline
Sprint 1 (Infrastructure) & 35 & 16 & 51 \\
\hline
Sprint 2 (Intelligence Artificielle) & 55 & 21 & 76 \\
\hline
Sprint 3 (Entreprise) & 34 & 19 & 53 \\
\hline
\textbf{Total plateforme} & \textbf{124} & \textbf{56} & \textcolor{ForestGreen}{\textbf{180/180 PASS}} \\
\hline
\end{tabular}
\end{table}


% ────────────────────────────────────────────────────────────
\section{Sprint Rétrospective}
\label{sec:s3_retro}
% ────────────────────────────────────────────────────────────

\begin{table}[H]
\centering
\caption{Rétrospective du Sprint 3}
\label{tab:retro_s3}
\begin{tabularx}{\textwidth}{|l|X|}
\hline
\textbf{Ce qui a bien fonctionné} & \begin{minipage}[t]{\linewidth}\vspace{2pt}
\begin{itemize}[noitemsep,topsep=0pt]
\item Le module de facturation contractuelle couvre les trois types de contrats industriels
\item Le déploiement Docker Compose permet un démarrage complet en une seule commande
\item Le pipeline CI/CD automatise l'ensemble du processus de livraison
\item L'interface de gestion des tenants offre une vue centralisée aux super administrateurs
\end{itemize}\vspace{2pt}
\end{minipage} \\
\hline
\textbf{Ce qui peut être amélioré} & \begin{minipage}[t]{\linewidth}\vspace{2pt}
\begin{itemize}[noitemsep,topsep=0pt]
\item La génération PDF des factures peut être enrichie avec un template plus élaboré
\item Le système d'alertes multi-canal (email, Slack, webhook) nécessite une configuration plus poussée
\item Les tests de charge en conditions réelles doivent être réalisés avant le déploiement à grande échelle
\end{itemize}\vspace{2pt}
\end{minipage} \\
\hline
\end{tabularx}
\end{table}


\section*{Conclusion}

Le Sprint~3 a achevé le développement de la plateforme LLM Dashboard en ajoutant les fonctionnalités entreprise essentielles à son exploitation commerciale. Le module de facturation contractuelle, avec ses trois types de contrats et sa génération automatique de factures, répond aux besoins spécifiques de TIM Tunisie pour la facturation de ses clients QALITAS et GMAO PRO. Le traçage des sessions, la gestion des prompts et l'administration des tenants complètent l'offre fonctionnelle. Le déploiement conteneurisé via Docker Compose et le pipeline CI/CD GitHub Actions garantissent une mise en production reproductible et automatisée. Au total, la plateforme est validée par \textbf{180 tests} couvrant l'ensemble des 11 modules fonctionnels.


% Conclusion générale et perspectives
% ============================================================
% CONCLUSION GÉNÉRALE ET PERSPECTIVES
% ============================================================

\chapter*{Conclusion Générale et Perspectives}
\addcontentsline{toc}{chapter}{Conclusion Générale et Perspectives}
\markboth{Conclusion Générale et Perspectives}{Conclusion Générale et Perspectives}

\vspace{0.5cm}

\section*{Bilan du projet}

Le présent travail de fin d'études, réalisé au sein de \textbf{TIM Tunisie}, avait pour objectif de concevoir et de développer une plateforme SaaS de gouvernance et de monitoring des modèles de langage de grande taille (LLM). Au terme de trois sprints de développement conduits selon la méthodologie Scrum, nous avons livré \textbf{LLM Dashboard}, une solution intégrée et opérationnelle répondant de manière exhaustive à la problématique identifiée.

La plateforme développée se distingue par la richesse et la complémentarité de ses \textbf{onze modules fonctionnels}~:

\begin{itemize}[noitemsep]
    \item Le \textbf{Gateway LLM} assure le proxying transparent des appels vers cinq fournisseurs (OpenAI, Anthropic, Google, Groq, Cohere), avec journalisation complète dans TimescaleDB et évaluation des règles de gouvernance en temps réel~;
    \item Le \textbf{moteur de gouvernance} applique cinq types de règles configurables (quota, budget, restriction de modèle, limite de tokens, restriction temporelle) avec des compteurs atomiques Redis~;
    \item Le \textbf{pipeline de détection d'anomalies} combine trois algorithmes de machine learning complémentaires --- STL decomposition~\cite{cleveland1990}, Isolation Forest~\cite{liu2008} et CUSUM~\cite{page1954} --- pour une couverture maximale des types d'anomalies~;
    \item Le \textbf{module Explainer} génère automatiquement des explications intelligibles et des recommandations actionnables à l'aide de LLM, avec stockage des embeddings vectoriels pour la recherche sémantique~;
    \item Le \textbf{moteur de prévision budgétaire} projette les coûts futurs selon trois scénarios (optimiste, moyen, pessimiste) avec identification des risques de dépassement~;
    \item Le \textbf{système de facturation contractuelle} supporte trois types de contrats (pay-as-you-go, forfait, hybride) avec génération automatique de factures détaillées~;
    \item Le \textbf{tableau de bord interactif}, développé avec Next.js~16 et React~19, offre une visualisation en temps réel via WebSocket de toutes les métriques de la plateforme.
\end{itemize}

\vspace{0.3cm}

Sur le plan technique, la plateforme repose sur une \textbf{architecture trois-tiers} robuste~: un backend FastAPI asynchrone, une base de données PostgreSQL enrichie des extensions TimescaleDB et pgvector, un cache multi-niveaux Redis, un pipeline de tâches asynchrones Celery, et un déploiement conteneurisé via Docker Compose avec six services orchestrés. La qualité du code est validée par \textbf{180 tests} (124 unitaires et 56 d'intégration) et un pipeline CI/CD automatisé via GitHub Actions.

\vspace{0.3cm}

Du point de vue académique, ce projet a permis de mobiliser et d'approfondir des connaissances dans des domaines variés~: l'architecture logicielle (Clean Architecture, principes SOLID)~\cite{martin2018}, l'apprentissage automatique non supervisé pour la détection d'anomalies~\cite{chandola2009}, les architectures multi-tenant SaaS~\cite{bezemer2010}, les API REST et les protocoles de sécurité (JWT, PBKDF2, TLS)~\cite{rfc7519,owasp2021}, et les systèmes de traitement de séries temporelles~\cite{hyndman2018}.


\section*{Limites identifiées}

Malgré les résultats obtenus, plusieurs limites ont été identifiées~:

\begin{enumerate}[noitemsep]
    \item \textbf{Tests de charge}~: les tests de performance sous charge élevée (milliers de requêtes simultanées) n'ont pas été réalisés en conditions réelles de production~;
    \item \textbf{Seuils de détection}~: les seuils des algorithmes de détection d'anomalies sont actuellement statiques et nécessiteraient un calibrage adaptatif par tenant~;
    \item \textbf{Génération PDF}~: le format des factures PDF pourrait être enrichi avec un template graphiquement plus élaboré~;
    \item \textbf{Alertes multi-canal}~: le système de notification est fonctionnel mais limité au webhook~; l'intégration avec les canaux email et Slack reste à compléter~;
    \item \textbf{Modèles de prévision}~: les algorithmes de prévision budgétaire utilisent des méthodes linéaires et pourraient bénéficier de modèles plus sophistiqués (ARIMA, Prophet).
\end{enumerate}


\section*{Perspectives d'évolution}

Les perspectives d'évolution de la plateforme LLM Dashboard sont multiples et s'inscrivent dans une logique d'amélioration continue~:

\begin{enumerate}[noitemsep]
    \item \textbf{Auto-calibrage des seuils ML}~: implémenter un mécanisme d'apprentissage en ligne (\textit{online learning}) pour ajuster automatiquement les seuils de détection d'anomalies en fonction du comportement de chaque tenant~;
    \item \textbf{Prévision avancée}~: intégrer des modèles de prévision de séries temporelles plus sophistiqués tels que ARIMA~\cite{box2015}, Prophet ou des réseaux de neurones récurrents (LSTM)~;
    \item \textbf{Analyse de sentiment des prompts}~: exploiter les modèles NLP pour analyser le contenu des prompts et détecter les usages potentiellement problématiques~;
    \item \textbf{Marketplace de modèles}~: développer une place de marché interne permettant aux tenants de découvrir et de comparer les modèles LLM disponibles~;
    \item \textbf{Déploiement Kubernetes}~: migrer de Docker Compose vers Kubernetes pour une scalabilité horizontale automatique et une haute disponibilité~;
    \item \textbf{Conformité réglementaire}~: intégrer des contrôles de conformité RGPD et AI Act pour garantir une utilisation responsable des LLM~;
    \item \textbf{Extension multi-cloud}~: supporter le déploiement sur les trois principaux clouds (AWS, GCP, Azure) avec abstraction de l'infrastructure.
\end{enumerate}

\vspace{0.5cm}

En définitive, ce projet de fin d'études nous a permis de concevoir et de livrer une plateforme complète et opérationnelle qui répond aux besoins réels de TIM Tunisie en matière de gouvernance des LLM. L'expérience acquise, tant sur le plan technique que méthodologique, constitue un socle solide pour les évolutions futures de la solution et pour notre parcours professionnel dans le domaine de l'ingénierie logicielle et de l'intelligence artificielle.


% ── Bibliographie ──────────────────────────────────────────
% ============================================================
% BIBLIOGRAPHIE
% ============================================================

\begin{thebibliography}{99}

% ── Modèles de langage et Transformers ─────────────────────

\bibitem{vaswani2017}
A.~Vaswani, N.~Shazeer, N.~Parmar, J.~Uszkoreit, L.~Jones, A.~N.~Gomez, Ł.~Kaiser et I.~Polosukhin,
\textit{Attention Is All You Need},
Advances in Neural Information Processing Systems (NeurIPS), vol.~30, pp.~5998--6008, 2017.

\bibitem{brown2020}
T.~Brown, B.~Mann, N.~Ryder \textit{et al.},
\textit{Language Models are Few-Shot Learners},
Advances in Neural Information Processing Systems (NeurIPS), vol.~33, pp.~1877--1901, 2020.

\bibitem{openai2023}
OpenAI,
\textit{GPT-4 Technical Report},
arXiv preprint arXiv:2303.08774, 2023.

\bibitem{touvron2023}
H.~Touvron, T.~Lavril, G.~Izacard \textit{et al.},
\textit{LLaMA: Open and Efficient Foundation Language Models},
arXiv preprint arXiv:2302.13971, Meta AI Research, 2023.

\bibitem{anthropic2024}
Anthropic,
\textit{The Claude Model Family: Capabilities and Safety},
Rapport technique, Anthropic, 2024.

% ── Gouvernance et LLMOps ─────────────────────────────────

\bibitem{mckinsey2024}
McKinsey \& Company,
\textit{The State of AI in Early 2024: Gen AI Adoption Spikes and Starts to Generate Value},
McKinsey Global Survey, mars 2024.

\bibitem{bommasani2022}
R.~Bommasani, D.~A.~Hudson, E.~Adeli \textit{et al.},
\textit{On the Opportunities and Risks of Foundation Models},
arXiv preprint arXiv:2108.07258, Stanford University, 2022.

\bibitem{chen2024llmops}
M.~Chen, J.~Tworek, H.~Jun \textit{et al.},
\textit{Practices and Challenges of LLM Operations in Production},
Proceedings of the 2024 IEEE International Conference on Software Engineering (ICSE), pp.~112--125, 2024.

% ── Détection d'anomalies ─────────────────────────────────

\bibitem{liu2008}
F.~T.~Liu, K.~M.~Ting et Z.-H.~Zhou,
\textit{Isolation Forest},
Proceedings of the 8th IEEE International Conference on Data Mining (ICDM), pp.~413--422, 2008.

\bibitem{cleveland1990}
R.~B.~Cleveland, W.~S.~Cleveland, J.~E.~McRae et I.~Terpenning,
\textit{STL: A Seasonal-Trend Decomposition Procedure Based on Loess},
Journal of Official Statistics, vol.~6, no.~1, pp.~3--33, 1990.

\bibitem{page1954}
E.~S.~Page,
\textit{Continuous Inspection Schemes},
Biometrika, vol.~41, no.~1--2, pp.~100--115, 1954.

\bibitem{chandola2009}
V.~Chandola, A.~Banerjee et V.~Kumar,
\textit{Anomaly Detection: A Survey},
ACM Computing Surveys, vol.~41, no.~3, article~15, pp.~1--58, 2009.

% ── Architecture logicielle ───────────────────────────────

\bibitem{fielding2000}
R.~T.~Fielding,
\textit{Architectural Styles and the Design of Network-Based Software Architectures},
Thèse de doctorat, University of California, Irvine, 2000.

\bibitem{martin2018}
R.~C.~Martin,
\textit{Clean Architecture: A Craftsman's Guide to Software Structure and Design},
Prentice Hall, 2018.

\bibitem{bezemer2010}
C.-P.~Bezemer et A.~Zaidman,
\textit{Multi-Tenant SaaS Applications: Maintenance Dream or Nightmare?},
Proceedings of the Joint ERCIM Workshop on Software Evolution and International Workshop on Principles of Software Evolution (IWPSE-EVOL), pp.~88--92, ACM, 2010.

% ── Méthodologie Scrum ─────────────────────────────────────

\bibitem{schwaber2020}
K.~Schwaber et J.~Sutherland,
\textit{The Scrum Guide: The Definitive Guide to Scrum -- The Rules of the Game},
Scrum.org, novembre 2020.

\bibitem{rubin2012}
K.~S.~Rubin,
\textit{Essential Scrum: A Practical Guide to the Most Popular Agile Process},
Addison-Wesley Professional, 2012.

% ── Frameworks et technologies ─────────────────────────────

\bibitem{fastapi2024}
S.~Ramírez,
\textit{FastAPI: Modern, Fast (High-Performance), Web Framework for Building APIs with Python},
Documentation officielle, \url{https://fastapi.tiangolo.com/}, consulté en avril 2025.

\bibitem{nextjs2024}
Vercel,
\textit{Next.js: The React Framework for the Web},
Documentation officielle, \url{https://nextjs.org/docs}, consulté en avril 2025.

\bibitem{timescaledb2024}
Timescale Inc.,
\textit{TimescaleDB Documentation: An Open-Source Time-Series SQL Database},
Documentation officielle, \url{https://docs.timescale.com/}, consulté en mars 2025.

\bibitem{celery2024}
Ask Solem \textit{et al.},
\textit{Celery: Distributed Task Queue},
Documentation officielle, \url{https://docs.celeryq.dev/}, consulté en avril 2025.

% ── Sécurité ──────────────────────────────────────────────

\bibitem{rfc7519}
M.~Jones, J.~Bradley et N.~Sakimura,
\textit{JSON Web Token (JWT)},
RFC 7519, Internet Engineering Task Force (IETF), mai 2015.

\bibitem{owasp2021}
OWASP Foundation,
\textit{OWASP Top Ten -- 2021: The Ten Most Critical Web Application Security Risks},
OWASP, 2021.

% ── Séries temporelles et prévision ───────────────────────

\bibitem{hyndman2018}
R.~J.~Hyndman et G.~Athanasopoulos,
\textit{Forecasting: Principles and Practice},
3\textsuperscript{e} édition, OTexts, Melbourne, Australie, 2021.

\bibitem{box2015}
G.~E.~P.~Box, G.~M.~Jenkins, G.~C.~Reinsel et G.~M.~Ljung,
\textit{Time Series Analysis: Forecasting and Control},
5\textsuperscript{e} édition, John Wiley \& Sons, 2015.
% ── Méthodologies d'ingénierie ────────────────────────────

\bibitem{evans2003}
E.~Evans,
\textit{Domain-Driven Design: Tackling Complexity in the Heart of Software},
Addison-Wesley Professional, 2003.

\bibitem{beyer2016}
B.~Beyer, C.~Jones, J.~Petoff et N.~R.~Murphy,
\textit{Site Reliability Engineering: How Google Runs Production Systems},
O'Reilly Media, 2016.

\bibitem{sculley2015}
D.~Sculley, G.~Holt, D.~Golovin \textit{et al.},
\textit{Hidden Technical Debt in Machine Learning Systems},
Advances in Neural Information Processing Systems (NeurIPS), vol.~28, pp.~2503--2511, 2015.

\bibitem{ford2017}
N.~Ford, R.~Parsons et P.~Kua,
\textit{Building Evolutionary Architectures: Support Constant Change},
O'Reilly Media, 2017.

\bibitem{limoncelli2014}
T.~A.~Limoncelli,
\textit{The Practice of Cloud System Administration: DevOps and SRE Practices for Web Services},
Addison-Wesley Professional, 2014.

\bibitem{gamma1994}
E.~Gamma, R.~Helm, R.~Johnson et J.~Vlissides,
\textit{Design Patterns: Elements of Reusable Object-Oriented Software},
Addison-Wesley Professional, 1994.

% ── API Gateway et architectures distribuées ──────────────

\bibitem{richardson2018}
C.~Richardson,
\textit{Microservices Patterns: With Examples in Java},
Manning Publications, 2018.

\bibitem{nginx2024}
F5 Networks,
\textit{NGINX Documentation: Reverse Proxy, Load Balancing, and API Gateway},
Documentation officielle, \url{https://docs.nginx.com/}, consulté en avril 2025.

% ── Séries temporelles ────────────────────────────────────

\bibitem{jensen2017}
S.~K.~Jensen, T.~B.~Pedersen et C.~Thomsen,
\textit{Time Series Management Systems: A Survey},
IEEE Transactions on Knowledge and Data Engineering, vol.~29, no.~11, pp.~2581--2600, 2017.

% ── Embeddings et recherche vectorielle ───────────────────

\bibitem{mikolov2013}
T.~Mikolov, K.~Sutskever, K.~Chen, G.~Corrado et J.~Dean,
\textit{Distributed Representations of Words and Phrases and Their Compositionality},
Advances in Neural Information Processing Systems (NeurIPS), vol.~26, pp.~3111--3119, 2013.

\bibitem{johnson2019}
J.~Johnson, M.~Douze et H.~Jégou,
\textit{Billion-Scale Similarity Search with GPUs},
IEEE Transactions on Big Data, vol.~7, no.~3, pp.~535--547, 2021.

% ── Architecture événementielle ──────────────────────────

\bibitem{hohpe2003}
G.~Hohpe et B.~Woolf,
\textit{Enterprise Integration Patterns: Designing, Building, and Deploying Messaging Solutions},
Addison-Wesley Professional, 2003.

\bibitem{kleppmann2017}
M.~Kleppmann,
\textit{Designing Data-Intensive Applications: The Big Ideas Behind Reliable, Scalable, and Maintainable Systems},
O'Reilly Media, 2017.

% ── Sécurité ──────────────────────────────────────────────

\bibitem{rfc2898}
B.~Kaliski,
\textit{PKCS \#5: Password-Based Cryptography Specification Version 2.0},
RFC 2898, Internet Engineering Task Force (IETF), septembre 2000.

\bibitem{oauth2012}
D.~Hardt,
\textit{The OAuth 2.0 Authorization Framework},
RFC 6749, Internet Engineering Task Force (IETF), octobre 2012.

% ── Conteneurisation ─────────────────────────────────────

\bibitem{merkel2014}
D.~Merkel,
\textit{Docker: Lightweight Linux Containers for Consistent Development and Deployment},
Linux Journal, vol.~2014, no.~239, article~2, 2014.

\bibitem{burns2018}
B.~Burns, J.~Beda et K.~Hightower,
\textit{Kubernetes: Up and Running --- Dive into the Future of Infrastructure},
2\textsuperscript{e} édition, O'Reilly Media, 2019.

% ── Prévision ─────────────────────────────────────────────

\bibitem{makridakis2018}
S.~Makridakis, E.~Spiliotis et V.~Assimakopoulos,
\textit{Statistical and Machine Learning Forecasting Methods: Concerns and Ways Forward},
PLoS ONE, vol.~13, no.~3, e0194889, 2018.

\bibitem{taylor2018}
S.~J.~Taylor et B.~Letham,
\textit{Forecasting at Scale},
The American Statistician, vol.~72, no.~1, pp.~37--45, 2018.

\end{thebibliography}


% ── Annexes ────────────────────────────────────────────────
% ============================================================
% ANNEXES
% ============================================================

\appendix

\chapter{Annexes}

\section{Structure complète du backend}

L'arborescence suivante illustre l'organisation modulaire du backend FastAPI~:

\begin{lstlisting}[language=bash, caption={Arborescence du projet backend}]
llm-dashboard/
|-- main.py                     # Point d'entree FastAPI
|-- celery_app.py               # Configuration Celery
|-- core/
|   |-- database.py             # Pool de connexions PostgreSQL
|   |-- redis.py                # 4 pools Redis dedies
|   |-- settings.py             # Configuration Pydantic
|   |-- event_bus.py            # Bus d'evenements asynchrone
|   |-- observability.py        # Metriques et monitoring
|   |-- rate_limit.py           # Middleware de rate limiting
|   |-- security_headers.py     # En-tetes de securite HTTP
|-- modules/
|   |-- auth/                   # M7 : Authentification et RBAC
|   |   |-- models.py           #   Entites SQLAlchemy
|   |   |-- schemas.py          #   DTO Pydantic
|   |   |-- router.py           #   Endpoints FastAPI
|   |   |-- service.py          #   Logique metier
|   |   |-- dependencies.py     #   Injection de dependances
|   |-- gateway/                # M1 : Proxy LLM + Gouvernance
|   |   |-- models.py
|   |   |-- router.py
|   |   |-- services/
|   |   |   |-- gateway_service.py
|   |   |   |-- governance_engine.py
|   |   |   |-- llm_provider.py
|   |   |   |-- semantic_cache.py
|   |-- analytics/              # M2 : Agregation des couts
|   |-- detection/              # M3 : Detection d'anomalies ML
|   |-- forecasting/            # M6 : Prevision budgetaire
|   |-- billing/                # M10 : Facturation contractuelle
|   |-- dashboard/              # M8 : Tableau de bord + WebSocket
|   |-- tracing/                # M9 : Tracage des sessions
|   |-- prompts/                # M11 : CMS de templates
|   |-- tenants/                # Gestion multi-tenant
|   |-- observability/          # Metriques Prometheus
|-- migrations/                 # Migrations Alembic
|-- tests/                      # 21 fichiers de tests pytest
|-- docker-compose.yml          # Orchestration des 6 services
|-- Dockerfile                  # Image multi-stage
|-- .github/workflows/ci.yml   # Pipeline CI/CD
\end{lstlisting}


\section{Configuration de sécurité Nginx}

\begin{lstlisting}[language=bash, caption={Extrait de la configuration Nginx (nginx/nginx.conf)}]
# Rate limiting multi-niveaux
limit_req_zone $binary_remote_addr zone=api:10m rate=30r/s;
limit_req_zone $binary_remote_addr zone=login:10m rate=1r/s;

server {
    listen 443 ssl http2;
    server_name llm-dashboard.tim-tunisie.com;
    
    # TLS 1.2/1.3 avec ciphers Mozilla Intermediate
    ssl_protocols TLSv1.2 TLSv1.3;
    ssl_prefer_server_ciphers on;
    
    # HSTS 1 an
    add_header Strict-Transport-Security 
               "max-age=31536000; includeSubDomains" always;
    
    # Proxy vers FastAPI
    location /v1/ {
        limit_req zone=api burst=20 nodelay;
        proxy_pass http://api:8000;
    }
    
    # WebSocket pour le dashboard temps reel
    location /v1/ws/ {
        proxy_pass http://api:8000;
        proxy_http_version 1.1;
        proxy_set_header Upgrade $http_upgrade;
        proxy_set_header Connection "upgrade";
    }
}
\end{lstlisting}


\section{Modèle de données --- Extrait SQLAlchemy}

\begin{lstlisting}[language=Python3, caption={Modèle LLMTokenLog (modules/gateway/models.py)}]
class LLMTokenLog(Base):
    """Journal des appels LLM — hypertable TimescaleDB."""
    __tablename__ = "llm_token_log"
    
    id: Mapped[uuid.UUID] = mapped_column(primary_key=True)
    tenant_id: Mapped[uuid.UUID] = mapped_column(index=True)
    timestamp: Mapped[datetime] = mapped_column(index=True)
    
    # Identite de l'appel
    provider: Mapped[str]         # openai, anthropic, google...
    model: Mapped[str]            # gpt-4o, claude-3.5-sonnet...
    agent_id: Mapped[str]         # Identifiant de l'application
    
    # Metriques
    input_tokens: Mapped[int]
    output_tokens: Mapped[int]
    total_tokens: Mapped[int]
    cost_usd: Mapped[Decimal]     # Cout calcule en USD
    latency_ms: Mapped[int]       # Latence en millisecondes
    
    # Statut
    status: Mapped[str]           # success, error, cached
    error_message: Mapped[Optional[str]]
    
    # Metadata
    session_id: Mapped[Optional[uuid.UUID]]  # Correlation session
    metadata_json: Mapped[Optional[dict]]    # JSONB flexible
\end{lstlisting}

\end{document}


\end{document}
